%  LaTeX support: latex@mdpi.com
%=================================================================
\documentclass[entropy,article,submit,pdftex,oneauthor]{Definitions/mdpi}

% Axiom environment
\newtheorem{Axiom}{Axiom}

% QED symbol
\renewcommand{\qedsymbol}{\textnormal{Q.E.D.}}

% Preprint mode: suppress MDPI auto-generated text
\makeatletter
\AtBeginDocument{%
  \renewcommand{\cright}{%
    \fontdimen2\font=1.2pt
    \raggedright\textbf{Copyright:} \copyright\ \the\year\ by the author.
    Preprint released under the Creative Commons Attribution
    (\href{https://creativecommons.org/licenses/by/4.0/}{CC~BY}) license.%
  }%
  \fancypagestyle{plain}{%
    \renewcommand{\headrulewidth}{0pt}%
    \renewcommand{\footrulewidth}{0pt}%
    \fancyhead{}%
    \fancyfoot[L]{\fontsize{8}{8}\selectfont Preprint --- \today}%
    \fancyfoot[C]{\thepage}%
    \fancyfoot[R]{}%
  }%
  \fancypagestyle{fancy}{%
    \renewcommand{\headrulewidth}{0pt}%
    \renewcommand{\footrulewidth}{0pt}%
    \fancyhead{}%
    \fancyfoot[L]{\fontsize{8}{8}\selectfont Preprint --- \today}%
    \fancyfoot[C]{\thepage}%
    \fancyfoot[R]{}%
  }%
  \pagestyle{plain}%
}
\makeatother

%=================================================================
% MDPI internal commands - do not modify
\firstpage{1}
\makeatletter
\setcounter{page}{\@firstpage}
\makeatother
\pubvolume{1}
\issuenum{1}
\articlenumber{}
\pubyear{2026}
\copyrightyear{2026}
\datereceived{}
\daterevised{}
\dateaccepted{}
\datepublished{}

%=================================================================
\Title{An Axiomatic Model: Axiomatic Derivation and Interpretation of the $(5,2)$ Signature in Wyler's Fine-Structure Constant Formula}

\newcommand{\orcidauthorA}{0009-0007-4132-1707}
\Author{Hyukjin Han \orcidA{}}

\AuthorNames{Hyukjin Han}

\address{%
Independent Researcher, Hwaseong, Gyeonggi-do, Republic of Korea; bokkamsun@gmail.com}

\corres{Correspondence: bokkamsun@gmail.com}

%=================================================================
\abstract{Four axioms (orthogonality, atomicity, minimal cost, causality preservation) and one proposition define a discrete computational system. Because no free parameters enter the deduction, fitting is excluded; because the axioms close the search space, numerological tuning is blocked. The complete axiomatic system comprises 15 axioms~\cite{han2026}; the present paper uses only 4 axioms and 1 proposition. The irreversible/reversible classification of cost corresponds to the signature of a quadratic form (Theorem~\ref{thm:cost-signature}), forcing the signature partition $(5,2)$ and excluding $\mathrm{SO}(7)$ and $\mathrm{SO}(4,3)$. All four factors of the Wyler formula receive an account within the axiomatic structure, though their epistemic status differs (I/II/III). Results: $1/\alpha = 137.036\,082$ ($6 \times 10^{-7}$), $\sin^2\theta_W = 0.23122$ ($4 \times 10^{-6}$); the baryon-to-photon ratio $\eta_B = 6.14 \times 10^{-10}$ ($0.5\sigma$) is an independent forward-testable production at axiom-described level (classification~III). The internal degrees of freedom $7 = 4+3$ coincide with the Hamming $[7,4,3]$ code, and by the \emph{Cl(0,7) Clifford form proposition} of the complete axiomatic system v1.6~\cite{han2026}, this coincidence is elevated to an algebraic isomorphism (data type $128 = 2^7 = \dim \mathrm{Cl}(0,7)$ multivector space). Table~\ref{tab:axiom-results} classifies 38 results (I:~13, II:~9, III:~15, IV:~1). Existing axiomatic reconstruction programs reconstruct the structure of quantum mechanics but do not produce physical constants. The present axiomatic system is a model designed to produce physical constants from axioms. The derivation process yields 5 unique predictions (absence of $Z'$, CAS pattern in higher-loop $\beta$-coefficients, nearest integer to $1/\alpha(M_Z) = 2^7$, etc.), each falsifiable by specific experiments (Table~\ref{tab:predictions}).}

\keyword{quantum information; Compare-And-Swap; fine-structure constant; Wyler formula; metric signature SO(5,2); Weinberg angle; quantum measurement; Landauer principle; Hamming error correcting code; computational primitives in physics}

%%%%%%%%%%%%%%%%%%%%%%%%%%%%%%%%%%%%%%%%%%
\begin{document}

%%%%%%%%%%%%%%%%%%%%%%%%%%%%%%%%%%%%%%%%%%
\section{Introduction}
\label{sec:intro}

The fine-structure constant is defined as
\begin{linenomath}
\begin{equation}
\alpha = \frac{e^2}{4\pi\varepsilon_0 \hbar c} \approx \frac{1}{137.036}
\label{eq:alpha-def}
\end{equation}
\end{linenomath}
and determines the strength of the electromagnetic interaction (Eq.~\eqref{eq:alpha-def}). This value corresponds to the low-energy (Thomson) limit at momentum transfer $q^2 \to 0$. Under QED renormalization, $\alpha$ varies with energy scale ($\alpha(M_Z) \approx 1/128$); the present paper derives this low-energy limiting value.

The origin of this particular value remains one of the deepest open problems in physics. Feynman called it ``one of the greatest damn mysteries of physics: a magic number that comes to us with no understanding by man''~\cite{feynman1985}.
In 1969, Wyler proposed a geometric formula that reproduces $\alpha$ as a volume ratio of the bounded symmetric domain $D_5 = \mathrm{SO}_0(5,2)/[\mathrm{SO}(5)\times\mathrm{SO}(2)]$ ($\mathrm{SO}_0$ denotes the identity component)~\cite{wyler1969,wyler1971}. It yields $1/\alpha = 137.036\,082$, matching five digits beyond the integer part. However, the formula was met with three criticisms~\cite{robertson1971,gilmore1972}:
\begin{itemize}
\item \textbf{(R1) Lack of physical motivation}: Why should $D_5$ be related to electromagnetism?
\item \textbf{(R2) Arbitrariness of group selection}: Why $\mathrm{SO}(5,2)$ and not some other group?
\item \textbf{(R3) Non-uniqueness of measure}: Choosing a different symmetric space can yield a different value (Gilmore).
\end{itemize}

These criticisms have remained unresolved for over half a century.


\textbf{Position of the present paper}. The complete axiomatic system comprises 15 axioms~\cite{han2026}. The present paper uses \emph{only 4 axioms and 1 proposition} to produce the following results. Starting from an axiomatic computational system whose primitive object is the Compare-And-Swap (CAS) atomic operation, we show that the $(5,2)$ metric signature is forced, describe all four factors of the Wyler formula within the axiomatic structure, and obtain $\alpha$, $\sin^2\theta_W$, and $\eta_B$ to precisions consistent with experimental values. Table~\ref{tab:axiom-results} catalogues a total of 38 results produced by the present axiomatic model.

The status of the four factors of the Wyler formula within the axiomatic system: $9$ is derived from the axioms (classification~I); $\pi^5/(2^4 \cdot 5!)$ is axiom-determined as a mathematical identity from the signature $(5,2)$ forced by the axioms (classification~II); $(\cdot)^{1/4}$ is axiom-determined as the geometric mean of 4 orthogonal independent axes (classification~II); and $8\pi^4$ is axiom-determined as the group volume determined by $(5,2)$ (classification~II). All four factors are accounted for by the axiomatic structure (see Table~\ref{tab:factors}).

All three Robertson--Gilmore criticisms (R1: lack of physical motivation, R2: arbitrariness of group choice, R3: non-uniqueness of measure) are answered. R1 and R2 are answered directly from the axioms; R3 is answered by Theorem~\ref{thm:cost-signature}, which eliminates the freedom in choosing symmetric spaces. Existing axiomatic reconstruction programs reconstruct the structure of quantum mechanics but do not produce physical constants. The present axiomatic system is a model designed to produce physical constants from axioms (see Table~\ref{tab:comparison}). As the classification in Table~\ref{tab:axiom-results} shows, the status of formal derivation~(I) and axiom-described results~(III) differs, and the present paper makes this distinction explicit.

Organization of the present paper: Section~\ref{sec:axioms} presents the design principles, methodological transparency, and the 4 axioms + 1 proposition. Section~\ref{sec:theorem} defines the cost--signature correspondence and proves the uniqueness of the cost classification. Section~\ref{sec:derivation} derives $\alpha$. Section~\ref{sec:prediction} presents the structural consistency check via $\sin^2\theta_W$. Section~\ref{sec:robertson} answers the Robertson/Gilmore criticisms. Section~\ref{sec:auxiliary} presents auxiliary results (Hamming code and $\eta_B$). Section~\ref{sec:discussion} discusses implications. Section~\ref{sec:conclusion} concludes.

\textbf{The constraints imposed on the axioms are powerful.} The four axioms of the present axiomatic system declare, respectively, minimal cost, atomicity, orthogonality, and causality preservation. These constraints force a logical minimal convergence: the structure satisfying all constraints converges to a unique point, and that point coincides with physical constants. Because the deduction starts from constraints, there are no free parameters and fitting is impossible; because the search space is closed by the axioms, numerology is blocked. The constraints, not the target, determine the output.

\textbf{We show that axioms describe nature.} The Banya equation was designed to be compatible with existing physics (Section~\ref{sec:transparency}, point~6); that compatibility is a design criterion, not a discovery. The non-trivial claim is that the cost structure of this equation produces specific numerical values matching experiment \emph{without fitting}. The present axiomatic system answers precisely this question: if the principle of least action is faithfully translated into a discrete cost language, do physical constants emerge from axioms? The answer is yes. Table~\ref{tab:axiom-results} summarizes all results produced by the 4 axioms and 1 proposition. The status of each result is classified into four levels: \textbf{I.~Axiom-derived} (formal proof or forward chain completed), \textbf{II.~Axiom-determined} (the axioms force the structure, which is then mathematically determined), \textbf{III.~Axiom-described} (derivation path and interpretation are given in the text, formal systematization reserved for dedicated papers), \textbf{IV.~Out of scope} (not included in the scope of the present paper).

%%%%%%%%%%%%%%%%%%%%%%%%%%%%%%%%%%%%%%%%%%
\section{Axiomatic System}
\label{sec:axioms}

\subsection{Design Principles}
\label{sec:framework}

This section explains what the axiomatic system models and why it takes this form. The explanation is intended to help the reader grasp the overall structure before encountering the axioms.

\subsection{Least Action and Least Cost}

The principle of least action in physics declares that ``nature selects the path that extremizes the action.'' The present axiomatic system reformulates this principle in a discrete cost language: ``every change follows the operational path that minimizes cost.'' Here, cost is the minimal unit ($+1$) paid when crossing an orthogonal boundary in the prescribed sequence.

The most important consequence of this reformulation is that \emph{contention creates order, and order creates discreteness}. When multiple local states coexist, contention arises; causality preservation forces an ordering; and that ordering breaks the process into discrete steps. This is the core principle that pervades the entire axiomatic system.

The direct benefit of this principle is that the presence or absence of cost ($0$ or $>0$) automatically classifies the 7 degrees of freedom into two categories---reversible and irreversible. This classification determines the metric signature, and this is the central result of the present paper.

Table~\ref{tab:least-cost-action} summarizes this correspondence. The key is the last row: logical minimal convergence---the convergence to a single structure when the minimal conditions of the axiomatic system are collected---follows the same logic as nature's convergence to a specific physical structure when following least action. The convergence point of the axioms is the convergence point of nature.

\begin{table}[H]
\begin{adjustwidth}{-\extralength}{0cm}
\caption{Correspondence between least cost (axiomatic system) and least action (nature).\label{tab:least-cost-action}}
\small
\begin{tabularx}{\fulllength}{p{4.8cm}p{4.8cm}X}
\toprule
\textbf{Axiomatic system (least cost)} & \textbf{Nature (least action)} & \textbf{Remarks} \\
\midrule
Follows the cost-minimizing path (Axiom~\ref{ax:cost}) & Selects the path that extremizes the action & Discrete/continuous expression of the same principle \\
Contention $\to$ forced order $\to$ discrete & Coexisting states $\to$ measurement $\to$ quantum discreteness & Causality preservation creates discreteness \\
CAS = unique least-cost operation (Axiom 2) & State changes in nature are atomic & Fewer than 3 steps violates causality; more than 3 violates least cost \\
Cost $\{0,+1\}$ binary classification & Irreversible/reversible classification & Presence or absence of order is the sole criterion \\
7 d.o.f.\ $\to$ $(5,2)$ unique convergence & Metric signature $(5,2)$ & Collecting minimal conditions yields convergence to one structure \\
\emph{Logical minimal convergence $\to$ $\alpha$, $\sin^2\theta_W$} & \emph{Least action in nature $\to$ physical constants} & \emph{Convergence point of axioms = convergence point of nature} \\
\bottomrule
\end{tabularx}
\end{adjustwidth}
\end{table}

\subsection{Terminological Note: The Meaning of ``Irreversible''}

In the present paper, ``irreversible'' means \emph{an ordered transition} in which the ordering is forced. This is related to, but not identical with, thermodynamic irreversibility (a process accompanied by entropy production). Landauer~\cite{landauer1961} showed that logical irreversibility (information erasure) entails a thermodynamic cost ($kT\ln 2$), and Bennett~\cite{bennett1973} showed that logically reversible computation can in principle be performed without thermodynamic cost. Whenever ``irreversible'' appears in the present paper, it should be read as ``a transition that cannot be undone because the ordering is forced.''

\textit{Structural correspondence between cost $+1$ and information erasure.} In the complete axiomatic system~\cite{han2026}, the Swap step of CAS (when Compare is true) overwrites the previous state of DATA with a new state. The previous state is irrecoverable---this is precisely the structure of information erasure as defined by Landauer. In the binary encoding of CAS (001$\to$011$\to$111), each transition is a 1-bit flip, and this flip irreversibly alters the bit state of the preceding step. Hence cost $+1$ corresponds to 1 bit of irreversible information erasure. Conversely, when Compare is false, Swap does not execute, no information erasure occurs, the cost is $0$, and the process is reversible. This correspondence is organized into three tiers.

\begin{center}
\small
\begin{tabular}{clp{7.2cm}l}
\toprule
\textbf{Tier} & \textbf{Category} & \textbf{Meaning of cost $+1$} & \textbf{Status} \\
\midrule
1 & Within axioms & Dimensionless counting measure tallying ordered-boundary crossings. Not in energy units & Established \\
2 & Structural correspondence & CAS Swap state overwrite $=$ Landauer 1-bit information erasure. 5 irreversible axes: erasure occurs; 2 reversible axes: no erasure ($c=0$) & Established \\
3 & Thermodynamic consequence & If Tier 2 holds: $kT\ln 2$/bit $\to$ CAS cycle heat lower bound $5 \times kT\ln 2$ & Conditional (III) \\
\bottomrule
\end{tabular}
\end{center}

What the present paper establishes is Tiers 1 and 2; Tier 3 is conditional on the Tier 2 structural correspondence (classification~III, Table~\ref{tab:axiom-results}).

\textit{Information-theoretic perspective (Hartley information).} In one CAS cycle, each of the 7 axes occupies an independent binary degree of freedom (0 or 1) (Axiom~\ref{ax:cas}: 3 mutually orthogonal steps; Axiom~\ref{ax:banya}: 4 mutually orthogonal axes). Orthogonality = independence, so under the maximum-entropy assumption (uniform prior), the pre-cycle state space is $2^7 = 128$.

When Compare is true and the 5 irreversible axes are determined and erased, the remaining uncertainty resides in the 2 reversible axes alone. In terms of Hartley information~\cite{hartley1928} (maximum entropy under a uniform prior; a special case of Shannon entropy), $\Delta H = \log_2(128) - \log_2(4) = 7 - 2 = 5$ bits. This erasure of 5 bits exactly matches the 5 irreversible axes of Theorem~\ref{thm:cost-signature}---the number of erased axes equals the number of erased bits. The 2 reversible axes (observer, superposition) undergo no information erasure and therefore do not contribute to $\Delta H$, corresponding to cost $0$.

\emph{Compatibility with Cl(0,7) algebra --- complete axiomatic system v1.6}. The algebraic identity of $2^7 = 128$ is identified, by the \emph{7-axis orthogonal Clifford form proposition} of the complete axiomatic system v1.6~\cite{han2026}, with the \emph{multivector space dimension of the Cl(0,7) Clifford algebra} (see Eq.~\eqref{eq:cl7-dimension}; Section~\ref{sec:predictions}). The present 5 bits of Hartley information (the erased irreversible 5 axes) are placed on the same 7 axes as either the 5-vector subspace of Cl(0,7) (dimension $\binom{7}{5} = 21$) or the irreversible-5 part of the (5,2) signature partition of Theorem~\ref{thm:cost-signature}, but at separate layers. The algebraic essence of data type 128 being Cl(0,7) is a different layer of classification from the cost-to-signature mapping of Theorem~\ref{thm:cost-signature}; both layers act on the same 7 generators.

The terminology follows Hartley's original formulation ($\log_2 N$ under a uniform prior) and corresponds to the uniform-distribution special case of Shannon's general entropy formula $-\sum p_i \log_2 p_i$.

\subsection{Glossary}

The axiomatic system of the present paper uses a discrete operational structure called ``conditional state transition.'' This structure is formally identical to the Compare-And-Swap (CAS) known in concurrent computing, but in the present paper it is translated into a physical context as summarized in Table~\ref{tab:glossary}.

\begin{table}[H]
\begin{adjustwidth}{-\extralength}{0cm}
\caption{Physical translation of operational terminology.\label{tab:glossary}}
\small
\begin{tabularx}{\fulllength}{p{3.8cm}p{4cm}X}
\toprule
\textbf{Physics term (this paper)} & \textbf{CS counterpart} & \textbf{Remarks} \\
\midrule
Conditional state transition & Compare-And-Swap (CAS) & Indivisible operation: observe$\to$compare$\to$record \\
Observe (Read) & Read & Fetch current state \\
Compare & Compare & Match against expected value \\
Record (Swap) & Swap & Alter state if condition met \\
Local state & Entity & Physical state at a spacetime point \\
Order-enforcing device & Lock & Constraint guaranteeing transition order \\
Causality violation & Race condition & Non-determinism from unsynchronized concurrent writes \\
Indivisible operation & Atomic operation & A single transition that cannot be interrupted \\
Periodic state check & Polling & State inspection at every tick \\
Cyclic register & Ring buffer & Finite-dimensional cyclic structure \\
\bottomrule
\end{tabularx}
\end{adjustwidth}
\end{table}

In what follows, ``CAS'' is used as shorthand for ``conditional state transition,'' and ``entity'' for ``local state.''

\subsection{Three Components}

The axiomatic system consists of three components:
\begin{enumerate}
\item \emph{Structure}: Four mutually orthogonal axes (time, space, observer, superposition) are divided into two brackets---the classical bracket (DATA) and the quantum bracket (OPERATOR).
\item \emph{Operator}: A single indivisible operation CAS (conditional state transition) proceeds in 3 steps---observe (Read), compare (Compare), record (Swap). CAS operates on the OPERATOR bracket side and records the result in the DATA bracket.
\item \emph{Cost}: Crossing an orthogonal boundary $+$ in order incurs cost $+1$; absence of order incurs cost $0$.
\end{enumerate}

Combining these three components yields 7 internal degrees of freedom (4 domain axes $+$ 3 CAS steps), and the cost rule classifies each axis as irreversible (cost $> 0$) or reversible (cost $= 0$). This classification determines the metric signature.

\subsection{Why This Operation}

Why must the primitive operation be a 3-step indivisible operation of
``observe $\to$ compare $\to$ record''?
We now argue that this choice is a physical necessity.

Multiple local states coexist in spacetime. A single global change ($\delta$) is projected through multiple observers into multiple local states. To record changes in these multiple local states \emph{while preserving causality}, one must first check (compare) the current state and then conditionally record. Recording without checking leads to non-determinism from concurrent writes---a causality violation.

CAS is the \emph{least-cost indivisible operation} satisfying this requirement:
\begin{itemize}
\item Observe (Read): Fetch the current state (cost $+1$, unavoidable).
\item Compare: Match against the expected value (cost $+1$, essential for causality preservation).
\item Record (Swap): Alter the state if matched (cost $+1$, the only path for state change).
\end{itemize}
Fewer than 3 steps cannot preserve causality (recording without checking); more than 3 steps violates least cost. ``Observe$\to$compare$\to$record'' is the causality-preserving least-cost indivisible operation: the lower bound (2 steps or fewer violates causality) and upper bound (4 steps or more violates least cost) force exactly 3 steps; LL/SC is isomorphic to this same 3-step structure; TAS fails due to the absence of Compare. Multi-word operations (DCAS, k-CAS, etc.)\ are parallel executions of single CAS and violate least cost. Within the axiomatic structure the 3-step form is unique. Complete classification of all conceivable external primitives is a meta-question analogous to the rule selection problem (Section~\ref{sec:limitations}).

Herlihy~\cite{herlihy1991} proved that in the shared-memory model, CAS is a universal primitive capable of achieving wait-free synchronization. This universality concerns synchronization capability within a register-based computational model and does not claim universality for physical interactions in general. The basis for selecting CAS as the sole operator in the present paper is not Herlihy's CS theorem but the least-cost requirement for causality preservation (the 3-step argument above).

Bennett~\cite{bennett1973} showed that logically reversible computation can in principle be performed without thermodynamic cost, and Landauer~\cite{landauer1961} showed that logical irreversibility (information erasure) entails a thermodynamic cost of $kT\ln 2$. The cost structure of the present axiomatic system (Axiom~\ref{ax:cost}) is the structural counterpart of the Landauer principle. By the structural correspondence of Section~\ref{sec:framework} (Tier 2), each $+$ crossing corresponds to 1 bit of irreversible information erasure, and the correspondence with Landauer's~\cite{landauer1961} $kT\ln 2$ follows (Tier 3, conditional; see Section~\ref{sec:framework}).

\textit{Note on alternative primitives.} In concurrent computing, LL/SC (Load-Linked/Store-Conditional) is another concurrent primitive with the same consensus power (consensus number $\infty$) as CAS. The logical decomposition of LL/SC is Read$\to$(implicit Compare)$\to$Conditional Write, the same 3-step structure as CAS, and applying it to the present axiomatic system yields the same cost structure (5+2). Hence CAS and LL/SC are logically isomorphic in the present axiomatic system. By contrast, Test-And-Set (TAS) has only 2 steps (Read$\to$Set) with no Compare step, so it overwrites the state without matching against the expected value---this leads to non-determinism from unsynchronized writes (causality violation) and fails to satisfy the causality-preservation condition. ``Unique operator'' means not that the specific name CAS is unique, but that the 3-step logical structure ``observe$\to$compare$\to$record'' is unique.

\textit{Multiple behaviors unified by the least-cost principle.} The 3-step uniqueness above concerns the \emph{structure} of CAS. However, the least-cost character of CAS extends beyond structure---CAS is an operational system that selects the least-cost path depending on the access target, and these selections determine the numerical results throughout the paper.

\begin{center}
\small
\begin{tabular}{p{2.8cm}p{4.5cm}l}
\toprule
\textbf{Behavior} & \textbf{What least cost forces} & \textbf{Reference} \\
\midrule
Object registration (ball creation) & Minimum orthogonal boundary crossings $= 13$. Fewer leaves ball incomplete, more violates least cost & Def.~\ref{lem:cost13} \\
Object identification (indexing) & RLU coupled access ($1{+}1/\pi \approx 1.3$) $<$ raw access ($4$). ${\sim}70\%$ reduction $\to$ coupled forced & §\ref{sec:baryogenesis} \\
Internal structural ratios & Cost reading + norm reading are two representations of the same $+$ crossing. Their combination determines $\sin^2\theta_W$ & §\ref{sec:prediction} \\
\bottomrule
\end{tabular}
\end{center}

All three behaviors are different applications of the same least-cost principle. CAS is the unique operator not only because of its 3-step structural uniqueness, but because these multiple behaviors are all unified by a single least-cost principle.

\textit{Ontological status of CAS --- distinction of two tiers.} The status of CAS separates into two tiers.

\emph{Within the axioms}: the 3-step ``observe$\to$compare$\to$record'' structure is unique within the axiomatic system (lower bound: causality violation; upper bound: least-cost violation; alternatives excluded: TAS fails, LL/SC isomorphic, DCAS violates least cost). This uniqueness is a consequence of the axiomatic structure, not an external choice.

\emph{Outside the axioms}: ``why these axioms?'' is a meta-question shared by every axiomatic system---Euclid's five postulates, the Zermelo--Fraenkel axioms, the von Neumann axioms of quantum mechanics. Positing CAS as the fundamental operation is a choice at the same tier as positing ``least action'' or ``minimum entropy production'' as physical principles, and its justification is verified a posteriori by results (agreement with physical constants).

The two tiers do not conflict: that the 3-step CAS is unique \emph{within} the axioms, and that the choice of axioms itself is a meta-question, are statements at different tiers.

\subsection{Scope of the Present Paper}

The complete system comprises 15 axioms~\cite{han2026} and addresses, beyond the fine-structure constant, quark masses, cosmological energy distribution, and many other physical constants. The present paper uses only \emph{4 axioms (Axioms~\ref{ax:banya}, \ref{ax:cas}, \ref{ax:cost}, \ref{ax:delta}) and 1 proposition (Proposition~\ref{ax:dof})}, with $\sin^2\theta_W$ as an additional structural consistency check. This demonstrates that 4 axioms suffice for a self-contained derivation, maximizing verifiability.

\textit{Note on the discrete/continuous distinction.} Axiom 3 of the complete axiomatic system~\cite{han2026} explicitly declares the discrete/continuous distinction of DATA/OPERATOR: DATA (classical bracket) is discrete and OPERATOR (quantum bracket) is continuous (classification~I, Table~\ref{tab:axiom-results}). The point in the present paper where this distinction operates is the derivation of $\sin^2\theta_W$ in Section~\ref{sec:prediction}: the mutual orthogonality of the 3 CAS axes (within the OPERATOR bracket) defines a continuous sphere, and the geometric constant $\pi$ of this sphere enters the formula. If OPERATOR were not continuous, the sphere would not be defined and $\pi$ could not arise.

\subsection{Reader's Guide: Four Common Misconceptions}

We highlight four points where the notation and terminology of the present axiomatic system differ from standard physics:
\begin{enumerate}
\item \emph{The 4 axes are not 4D spacetime}: The 4 axes of the present axiomatic system (time, space, observer, superposition) differ from Minkowski 4D spacetime. Spacetime consists of the \emph{2 axes} time + space (DATA bracket), while observer + superposition are the \emph{2 axes} of the quantum domain (OPERATOR bracket). Not all 4 axes constitute spacetime.
\item \emph{$\delta$ is ``change,'' not energy}: The $\delta$ in Eq.~\eqref{eq:banya} is a unit-free quantity of change; energy, distance, probability, etc.\ are all ways of measuring $\delta$. Units are assigned to $\delta$ itself only \emph{after} substituting constants into the norm.
\item \emph{$+$ denotes orthogonal composition + two readings}: The $+$ in Eq.~\eqref{eq:banya} is a structural notation meaning ``two mutually orthogonal axes belong to the same bracket.'' In the present axiomatic system, the $+$ symbol additionally carries two \emph{reading} modes: (i) \emph{Cost reading}: reading a $+$ crossing as unit cost per Axiom~\ref{ax:cost} $\to$ \emph{integer} (e.g., 9 crossings of $+$ = 9). (ii) \emph{Norm reading}: reading the same $+$ crossing as a dimension of the orthogonal norm $\to$ \emph{geometric measure} $\pi$ (e.g., 9 irreversible $+$ crossings = $9\pi$; see the ``geometric measure'' paragraph following Definition~\ref{def:residual}).
The two readings are different expressions of \emph{the same entity}. However, the norm reading is valid only for \emph{irreversible $+$ crossings} (reversible crossings have cost = 0, so no arc representation---reversible quantities admit only cost-reading integers). Throughout the present paper, the use of $+$ carries both meanings and both readings \emph{depending on context} (see Section~\ref{sec:prediction}, 5-step derivation, step (4)).
\item \emph{Status of the observer axis}: Placing ``observer'' as a fundamental axis of the axiomatic system is a strategic choice that promotes the century-old unsolved measurement problem of quantum mechanics from ``something to be explained'' to ``a structural axis.'' This is analogous to the strategy in general relativity of defining ``gravity'' as spacetime curvature.
\end{enumerate}

%%%%%%%%%%%%%%%%%%%%%%%%%%%%%%%%%%%%%%%%%%


\subsection{Methodological Transparency}
\label{sec:transparency}


\begin{enumerate}
\item \emph{12 items vs.\ ``4 axioms + 1 proposition.''} The label ``4 axioms + 1 proposition'' refers to the logical skeleton. The actual derivation employs, in addition to the 4 axioms and 1 proposition, 2 operational definitions, 5 axiomatic-structure consequences (A1a--A1d, A2a), for a total of $4{+}1{+}2{+}5 = 12$ items (Table~\ref{tab:assumptions} enumerates all of them). These are not hidden inputs but \emph{explicitly listed} inputs. This is structurally analogous to Euclid's use of definitions and common notions alongside his five postulates: the postulates are the logical core, but the derivations require the full inventory. Crucially, none of the 12 items is a free parameter: each is either a logical consequence of the axioms, a standard mathematical operation applied to axiom-defined quantities, or an observed fact explicitly declared. The search space is fully constrained by the axioms, and parameter fitting remains impossible.

\item \emph{3-dimensional spatiality is an imported observed fact.} Three-dimensional space is not derived from the axioms. The classical bracket of Axiom~\ref{ax:banya} declares spacetime by importing classical physics, and three-dimensionality enters as an observed fact (Definition~\ref{lem:cost13}: $\mathrm{space} = x+y+z$). This is a design choice: the axiomatic system explicitly demarcates what it explains from what it presupposes.

\item \emph{All items are axiom-grounded.} Every item in Table~\ref{tab:assumptions} is either an explicit axiom, a proposition, an operational definition, or an axiomatic-structure consequence. No independent items beyond the axioms are used.

\item \emph{The four Wyler-formula factors are not of equal status.} $9$ is axiom-derived (I); $8\pi^4$ and $\pi^5/(2^4 \cdot 5!)$ are axiom-determined (II); $(\cdot)^{1/4}$ is axiom-determined (II) as the geometric mean of 4 independent axes. The statement ``all four factors are described within the axiomatic structure'' means that all four receive an account within the system---not that all four are formally derived. This distinction is the reason Table~\ref{tab:axiom-results} exists: the four-level classification (I--IV) makes the epistemic status of every result explicit.

\item \emph{$\eta_B$ is an auxiliary result at axiom-described level (classification~III).} The correction form $(4+1/\pi)$ of the baryon-to-photon ratio depends on the specific coupling mechanism of RLU coupled access. RLU itself is axiomatically required for CAS operation (defined after Axiom~\ref{ax:cas}), the $+$ in $(4+1/\pi)$ is uniquely determined by Axiom~\ref{ax:banya}'s orthogonal composition (Section~\ref{sec:prediction} step~(4)); the formal systematization of the full RLU mechanism is classification~III. It is included because a forward-testable prediction from the same axiomatic system has value for falsifiability, even at a lower epistemic tier than the core results ($\alpha$, $\sin^2\theta_W$).

\item \emph{What is designed and what is not.} The Banya equation (Axiom~\ref{ax:banya}) was designed so that existing physics equations can be recovered by substituting known physical quantities into the orthogonal norm---the classical bracket imports classical spacetime, and the quantum bracket imports observer and superposition. That the axioms are \emph{compatible} with known physics is therefore a design criterion, not a discovery. What is \emph{not} designed is that the cost structure of this substitution-compatible equation produces specific numerical values ($\alpha = 1/137.036\ldots$, $\sin^2\theta_W = 0.23122$) matching experiment. The design input is ``translate least action faithfully into discrete cost''; the numerical output is an emergent consequence of that translation, not a fitting target.
\end{enumerate}


\subsection{List of Assumptions}
\label{sec:assumption-list}

The present axiomatic system is labeled ``4 axioms + 1 proposition,'' but listing \emph{all} items actually used in deriving $\alpha$, $\sin^2\theta_W$, and $\eta_B$ on an equal footing yields the following. Table~\ref{tab:assumptions} is the \emph{complete item table} required to derive the results of the present paper, and it makes explicit that the ``4 axioms'' label is \emph{not} a complete list of items.

\begin{table}[H]
\begin{adjustwidth}{-\extralength}{0cm}
\caption{List of items used in deriving the results of the present paper.\label{tab:assumptions}}
\small
\begin{tabularx}{\fulllength}{p{1.0cm}p{3.8cm}X}
\toprule
\textbf{ID} & \textbf{Name} & \textbf{Remarks} \\
\midrule
A1 & Axiom 1 (Banya equation) & Explicit axiom. $\delta^2$ equation (Axiom~\ref{ax:banya}) \\
A2 & Axiom 2 (CAS 3-step) & Explicit axiom. R--C--S (Axiom~\ref{ax:cas}) \\
A3 & Axiom 3 (cost $+1$) & Explicit axiom (Axiom~\ref{ax:cost}) \\
A4 & Axiom 4 ($\delta$ flag) & Explicit axiom (Axiom~\ref{ax:delta}) \\
P1 & Proposition 1 (descriptive d.o.f.) & Consequence of A1+A2 (Lemma). ``Lemma demotion'' stated in text (Proposition~\ref{ax:dof}) \\
D1 & Definition 1 ($13 = 8 + 5$) & Operational definition. Cost enumeration (Definition~\ref{lem:cost13}) \\
D2 & Definition 2 (ball value $4 = 1 + 3$) & Operational definition. 3-dimensional decomposition of the space axis of Axiom~\ref{ax:banya} (Definition~\ref{def:residual}) \\
A1c & Prefactor 2 ($\eta_B$) & Bracket count $= 2$ of Axiom~\ref{ax:banya}. Axiomatic-structure consequence (§\ref{sec:baryogenesis}) \\
A1d & Two-reading principle & Axiom~\ref{ax:banya} (norm reading of $+$) + Axiom~\ref{ax:cost} (cost reading of $+$). Axiomatic-structure consequence (§\ref{sec:prediction}) \\
A2a & RLU $1/\pi$ combination & Axioms~\ref{ax:cas} (RLU required) + \ref{ax:banya} (4 axes) + \ref{ax:cost} (irreversible hemisphere). $(4+1/\pi)$ is a pointer-dereference structure. Axiomatic-structure consequence (§\ref{sec:baryogenesis}) \\
A1a & 3-dimensional spatiality & Observational decomposition of the space axis declared by Axiom~\ref{ax:banya}. Made explicit as $x+y+z$ in Definition~\ref{lem:cost13} \\
A1b & Wyler-Robertson form & $8\pi^4$ is axiom-determined (II) from $(5,2)$ group volume. $(\cdot)^{1/4}$ is axiom-determined (II) as the geometric mean of 4 orthogonal independent axes (Table~\ref{tab:factors}) \\
\bottomrule
\end{tabularx}
\end{adjustwidth}
\end{table}

That is, behind the ``4 axioms + 1 proposition'' label, there are in fact \emph{4 explicit axioms + 1 proposition + 2 operational definitions + 5 axiomatic-structure consequences (A1a--A1d, A2a) = 12 items}. A1a (3-dimensional space) is the decomposition of the space axis of Axiom~\ref{ax:banya}, and A1b (Wyler-Robertson form) arises from the axiomatic structure through the $(5,2)$ group volume (II) and the geometric mean of 4 axes (II). This table collects all 12 items in one place.

\textit{Note on axiomatic-structure consequences.} A1a--A1d and A2a are consequences of the explicit axioms: A1a (3D spatiality) from Axiom~\ref{ax:banya}; A1b (Wyler-Robertson form) from $(5,2)$; A1c (prefactor 2) from the bracket count; A1d (two readings) from Axiom~\ref{ax:banya} (norm) + Axiom~\ref{ax:cost} (cost) applied to the same $+$; A2a (RLU combination) from Axioms~\ref{ax:cas},~\ref{ax:banya},~\ref{ax:cost}. No independent items beyond the axioms are used.

\subsection{4 Explicit Axioms + 1 Proposition}

The following 4 explicit axioms + 1 proposition form the core skeleton of the derivation of $\alpha$, used together with the items in Table~\ref{tab:assumptions} (D1, D2, A1a--A1d, A2a). The numbering correspondence with the complete 15-axiom system~\cite{han2026} is:

\begin{center}
\small
\begin{tabular}{ccl}
\toprule
\textbf{This paper} & \textbf{Complete system} & \textbf{Content} \\
\midrule
Axiom 1 & Axiom 1 & Banya equation (4-axis orthogonal norm) \\
Axiom 2 & Axiom 2 & CAS as unique operator \\
Axiom 3 & Axiom 4 & Cost \\
Axiom 4 & Axiom 15 & $\delta$ global flag \\
\bottomrule
\end{tabular}
\end{center}

\noindent The ``Axiom 3 of the complete system''~\cite{han2026} repeatedly referenced in the present paper is the axiom that declares \emph{DATA (classical bracket) is discrete and OPERATOR (quantum bracket) is continuous}. This is a different axiom from Axiom 3 (cost) of the present paper.

\begin{Axiom}[Banya Equation --- 4-Axis Orthogonal Norm]\label{ax:banya}
Every change $\delta$ in the universe is the norm of 4 mutually orthogonal axes (this equation is called the \emph{Banya equation}):
\begin{linenomath}
\begin{equation}
\delta^2 = (\mathrm{time} + \mathrm{space})^2 + (\mathrm{observer} + \mathrm{superposition})^2.
\label{eq:banya}
\end{equation}
\end{linenomath}
The first bracket $(\mathrm{time} + \mathrm{space})$ is called the \emph{classical bracket} (DATA), and the second $(\mathrm{observer} + \mathrm{superposition})$ the \emph{quantum bracket} (OPERATOR). The two brackets are orthogonal. The $+$ denotes not arithmetic addition but orthogonal composition: within the same bracket, the two axes are orthogonal and together span a single subspace, and the $+$ between brackets denotes the orthogonal sum of distinct subspaces.

\textit{Necessity of the 4 axes.} Observer and superposition are not metaphors but structurally essential components. Without the observer axis, there is no destination for the branching result of CAS Compare (Axiom~\ref{ax:cas}), making CAS execution impossible. Without superposition, CAS has no means of referencing DATA, making all operations impossible. Hence the quantum bracket requires at least 2 axes (observer, superposition), and the classical bracket also requires at least 2 axes (time, space)---4 axes is the minimal configuration.
\end{Axiom}

\textit{Construction principle of the Banya equation.} The two brackets of Eq.~\eqref{eq:banya} each import existing physics. The classical bracket $(\mathrm{time} + \mathrm{space})$ imports the spacetime of classical physics directly---time and 3-dimensional space are observed facts that the axiom declares rather than proves. The quantum bracket $(\mathrm{observer} + \mathrm{superposition})$ imports phenomena discovered by quantum physics---the observer effect and superposition. Both are real physical phenomena, and Axiom~\ref{ax:banya} orthogonalizes them. The $\delta$ on the left-hand side is a new physical quantity absent from existing physics: a \emph{quantity of change} obtained by orthogonal composition of the classical and quantum brackets, which is neither energy, nor distance, nor probability, but unit-free change itself. Energy, mass, distance, etc.\ are ways of measuring $\delta$ after assigning units.

\textit{Separation of system time and domain time --- immediate consequence of Axiom~\ref{ax:banya}.}\label{rem:two-times} Equation~\eqref{eq:banya} structurally separates \emph{two kinds of time}. The left-hand side $\delta$ is the \emph{global quantity of change} spanning the entire equation (system time: 1 cycle $=$ 1 tick), while the time axis on the right-hand side is a single \emph{local domain axis} inside the classical bracket. Because $\delta$ is the norm of the entire right-hand side, $\delta$ and time occupy different tiers---$\delta$ \emph{contains} time but is \emph{not equal to} time. This separation is forced the moment Axiom~\ref{ax:banya} is read; no additional axiom is required. Axiom~3 of the complete system~\cite{han2026} (DATA discrete / OPERATOR continuous) further specifies this separation, but the separation itself is an immediate consequence of Axiom~\ref{ax:banya} (classification~I). All subsequent references to ``system time'' and ``domain time'' in this paper are anchored to this paragraph.

\begin{Axiom}[Conditional State Transition (CAS) Is the Unique Operator]\label{ax:cas}
Every change in the universe is a repetition of a single conditional state transition (CAS: Compare-And-Swap) operation. CAS proceeds in 3 steps:
\begin{linenomath}
\begin{equation}
\underset{\text{observe}}{\mathrm{Read}}\ (001) \;\longrightarrow\; \underset{\text{compare}}{\mathrm{Compare}}\ (011) \;\longrightarrow\; \underset{\text{record}}{\mathrm{Swap}}\ (111).
\label{eq:cas-stages}
\end{equation}
\end{linenomath}
Each step occupies an independent binary degree of freedom (the independent bits of each step are 001, 010, 100; the notation in Eq.~\eqref{eq:cas-stages} represents cumulative states). The 3 steps are mutually orthogonal: $\mathrm{R} \perp \mathrm{C} \perp \mathrm{S}$. CAS operates in the OPERATOR bracket and records the result in the DATA bracket at the record (Swap) step. The transition $\mathrm{R} \to \mathrm{C} \to \mathrm{S}$ is order-forced by logical dependence: Compare cannot execute without the result of Read, and Swap cannot execute without the result of Compare.
\end{Axiom}

\textit{CAS 3-axis orthogonality, sphere, and RLU.} The 3 steps of CAS (R, C, S) are mutually orthogonal. In the norm space spanned by 3 independent orthogonal axes, the constraint surface of constant norm is a sphere. When CAS executes Swap on the space axis of DATA, the result is recorded isotropically (spherically) by virtue of the 3-axis orthogonality.

\emph{RLU} is an \emph{addressless virtual indexing unit} operating on this sphere. The name is adapted from LRU (Least Recently Used), the standard cache-management technique in computer science, with the first letter changed to R to emphasize the addressless property. RLU is located within the OPERATOR bracket and can access all 4 domain coordinates of objects registered in DATA slots---2 on the DATA side and 2 on the OPERATOR side. CAS operates by referencing RLU, and in the present axiomatic system all causality is managed by RLU.

\textit{Why addressless.} On a sphere, logical addresses (memory addresses, slot numbers) do not exist---objects are identified only by angular coordinates ($\theta$, $\varphi$). Unlike address-based indices in standard computer science, RLU does not use logical addresses; it is an \emph{addressless angular index} that distinguishes objects by the angular coordinates of the OPERATOR bracket~\cite{han2026}.

\textit{Why $1/\pi$.} Irreversible $+$ crossings allow only the forward direction (Axiom~\ref{ax:cas}), so only half of the sphere---an arc length of $\pi$---is accessible. RLU identifies one object on this hemisphere per $1/\pi$ of arc. This is an axiomatic-system-internal identification rule, not the arithmetic reciprocal of $\pi$: it belongs to a separate category that is neither cost reading nor norm reading. The $1/\pi$ appearing in the $\eta_B$ correction $(4 + 1/\pi)$ originates here.

\begin{Axiom}[Cost]\label{ax:cost}
Crossing an orthogonal boundary $+$ in a prescribed order incurs cost $+1$. In the absence of order, the cost is $0$. Cost is a dimensionless counting measure that tallies the number of order-forced crossings; it is not in energy or temperature units.
\begin{itemize}
\item The CAS transitions $\mathrm{R} \to \mathrm{C} \to \mathrm{S}$ are order-forced (Axiom~\ref{ax:cas}), so each transition incurs cost $+1$.
\item When CAS records in the DATA bracket, it crosses the OPERATOR$\to$DATA bracket boundary, incurring cost $+1$.
\item When CAS accesses axes within its own OPERATOR bracket (observer, superposition), no bracket boundary is crossed, so cost $= 0$.
\end{itemize}
Operations with cost $> 0$ are ordered and hence irreversible. Access with cost $= 0$ is unordered and hence reversible. This branching determines the partition into 5 irreversible and 2 reversible axes, $(5,2)$. When CAS Compare is true, Swap executes, paying cost $+1$, and a definite state is recorded in DATA (collapse). When Compare is false, Swap does not execute, cost is $0$, and superposition is maintained. Quantum is the default; classical results from paying cost.

Cost is the only physical quantity in the present axiomatic system. Energy, mass, force, and entropy are all different names for cost~\cite{han2026}. No physical quantity other than cost exists within the axiomatic system.
\end{Axiom}

\begin{Proposition}[Descriptive Degrees of Freedom]\label{ax:dof}
The total number of internal degrees of freedom divides into \emph{structure} (DATA) and \emph{cost} (OPERATOR), and is a direct consequence of Axiom~\ref{ax:banya} (4 axes) and Axiom~\ref{ax:cas} (3 steps):
\begin{linenomath}
\begin{equation}
\underbrace{4}_{\text{domain axes}} \;+\; \underbrace{3}_{\text{CAS steps}} = 7.
\label{eq:dof7}
\end{equation}
\end{linenomath}
The $+$ has the same meaning as the orthogonal composition in Axiom~\ref{ax:banya}. Since the two sets belong to different brackets, they are combined by orthogonal sum ($+$).
\end{Proposition}

\begin{proof}
Axiom~\ref{ax:banya} declares 4 mutually orthogonal axes (time, space, observer, superposition). Axiom~\ref{ax:cas} declares that CAS consists of 3 mutually orthogonal steps (Read, Compare, Swap). The two sets belong to different brackets (domain axes belong to the space declared by Axiom~\ref{ax:banya}; CAS steps are the operator domain), so they combine by orthogonal sum ($+$). The total internal degrees of freedom are $4 + 3 = 7$.
\end{proof}

\emph{Note}: This proposition is a direct logical consequence of Axioms~\ref{ax:banya} and~\ref{ax:cas} and is therefore not an independent axiom but a lemma. The axiom count is 4: Axiom~\ref{ax:banya} (Banya equation), Axiom~\ref{ax:cas} (CAS as unique operator), Axiom~\ref{ax:cost} (cost), and Axiom~\ref{ax:delta} ($\delta$ global flag).

Table~\ref{tab:dof} classifies the descriptive degrees of freedom derived from the 7 internal degrees of freedom into structure (DATA) and cost (OPERATOR) categories. Every number appearing in the derivations of the present paper originates from this table.

\begin{table}[H]
\begin{adjustwidth}{-\extralength}{0cm}
\caption{Descriptive degrees of freedom --- structure (DATA) and cost (OPERATOR).\label{tab:dof}}
\small
\begin{tabularx}{\fulllength}{p{1.2cm}p{3.5cm}X}
\toprule
\multicolumn{3}{l}{\emph{Structural descriptive d.o.f.\ (DATA category)}} \\
\midrule
\textbf{Value} & \textbf{Origin} & \textbf{Remarks} \\
\midrule
$1$ & Minimum unit & Bit basis \\
$2$ & 2 brackets & DATA, OPERATOR (Axiom~\ref{ax:banya}) \\
$3$ & CAS 3 steps & R, C, S (Axiom~\ref{ax:cas}) \\
$4$ & 4 domains & time, space, observer, superposition (Axiom~\ref{ax:banya}) \\
$7$ & $4+3$ & CAS internal d.o.f.\ (Proposition~\ref{ax:dof}). Numerator of $\sin^2\theta_W$ \\
\midrule
\multicolumn{3}{l}{\emph{Cost descriptive d.o.f.\ (OPERATOR category)}} \\
\midrule
$1$ & Minimum cost per $+$ crossing & Cost basis (Axiom~\ref{ax:cost}) \\
$2$ & 2 reversible axes & observer, superposition. Cost $0$ (Theorem~\ref{thm:cost-signature}) \\
$4$ & Ball value $1+3$ & time write $1$ + space write $3$ (Definition~\ref{def:residual}) \\
$5$ & 5 irreversible axes & $(5,2)$ partition (Theorem~\ref{thm:cost-signature}). Exponent of $\pi^5$ \\
$9$ & Residual cost $13-4$ & $9\pi$ in the denominator of $\sin^2\theta_W$ (Definition~\ref{def:residual}) \\
$13$ & Total cost $8+5$ & 8 reads $+$ 5 writes (Definition~\ref{lem:cost13}) \\
\bottomrule
\end{tabularx}
\parbox{\fulllength}{Structure and cost are orthogonal brackets (Axiom~\ref{ax:banya}). Within the same bracket, only shifts ($2^N$) are possible; combination by $+$ occurs only when crossing between brackets. Every number appearing in the derivations of $\alpha$, $\sin^2\theta_W$, and $\eta_B$ in the present paper originates from this table.}
\end{adjustwidth}
\end{table}

\begin{Definition}[Cost Enumeration of One CAS Cycle --- Working Definition]\label{lem:cost13}
The space axis has 3-dimensional orthogonal sub-axes $(x, y, z)$---the `space' of Axiom~\ref{ax:banya} is a macro notation for these 3 dimensions, and when CAS accesses space it sequentially accesses the 3 sub-axes (an observed fact of our universe and a working assumption of the present paper).

In the present axiomatic system, a \emph{ball} is a single unit of matter registered in spacetime (a particle at the \emph{Planck-mass} scale in the mass dimension---the \emph{minimum mass unit} of the axiomatic system), and a \emph{grasp} (juim) is the act of CAS registering one ball in a DATA slot---i.e., the result of one CAS Swap. The cost required to create one ball is the $13$ $+$ crossings of the present definition; note that the ``minimum cost unit'' ($+1$, one $+$ crossing) and the ``minimum mass unit'' (ball, the result of $13$ $+$ crossings) are different quantities. The cost calculation in the present paper merely counts the dimensionless number of $+$ crossings required to ``register one ball''; mass scale is the dimensional context of the axiomatic system and lies outside the scope of the present paper.

On this microscopic decomposition, we \emph{define} the total number of $+$-boundaries crossed by one CAS cycle (one grasp creation: idle $\to$ R $\to$ C $\to$ S $\to$ DATA commit) to be exactly $13$:
\begin{linenomath}
\begin{equation}
13 = \underbrace{8}_{\text{reads (moving)}} \;+\; \underbrace{5}_{\text{writes (grasping)}}.
\label{eq:cost13}
\end{equation}
\end{linenomath}
Each $+$ crossing incurs cost $+1$ by Axiom~\ref{ax:cost}. The following table enumerates the complete path of one CAS cycle. Each row is a step forced by the axioms; removing any step leaves the ball incomplete, adding any step violates least cost. The enumeration is therefore unique.

\begin{center}
\small
\begin{tabular}{lcc}
\toprule
\textbf{Path} & \textbf{Read (access)} & \textbf{Write (grasp)} \\
\midrule
CAS R entry & $+1$ & --- \\
CAS R $\to$ C transition & $+1$ & --- \\
CAS C $\to$ S transition & $+1$ & --- \\
OPERATOR $\to$ DATA bracket boundary & $+1$ & --- \\
time $\to$ space & $+1$ & --- \\
$x$ access & $+1$ & --- \\
$x \to y$ & $+1$ & --- \\
$y \to z$ & $+1$ & --- \\
time timestamp write & --- & $+1$ \\
$x$ write & --- & $+1$ \\
$y$ write & --- & $+1$ \\
$z$ write & --- & $+1$ \\
Swap $\to$ DATA commit & --- & $+1$ \\
\midrule
\textbf{Total} & \textbf{8} & \textbf{5} \\
\bottomrule
\end{tabular}
\end{center}
\textit{Status note}: This definition is a \emph{working definition} that applies the cost accounting of the complete Banya axiomatic system~\cite{han2026} to the 7-axis (4 domain + 3 CAS) system of the present paper and the 3-dimensional decomposition of the space axis of Axiom~\ref{ax:banya}. The formal distinction of the 13 boundaries has been described (classification~I, Table~\ref{tab:axiom-results}).
\end{Definition}

\textit{Granularity note}: The micro-level decomposition (space $= x + y + z$) in this definition serves cost-accounting purposes only and does not affect the signature $(5,2)$ classification of Theorem~\ref{thm:cost-signature}. Theorem~\ref{thm:cost-signature} determines the signature by applying the cost function to the 4 macro axes of Axiom~\ref{ax:banya} and the 3 CAS stages of Axiom~\ref{ax:cas}, totalling 7 macro DOF; it operates independently of the micro decomposition of space. The two levels (macro signature, micro enumeration) are descriptions at different granularities of the same axiomatic system.

\emph{Three-dimensional spatiality}. (1)~The classical bracket $(\mathrm{time} + \mathrm{space})$ in Axiom~\ref{ax:banya} inherits classical-physics spacetime as given. (2)~Because space is three-dimensional in classical physics, declaring space automatically imports a three-dimensional structure. (3)~When the CAS creates a ball, it writes into a 4-dimensional spacetime: 1 time axis $+$ 3 space axes. Consequently, the ``ball value $4 = 1 + 3$'' and ``residual cost $9 = 13 - 4$'' of Definition~\ref{def:residual} depend on this 4-dimensional spacetime structure.

\begin{Definition}[Ball value and residual cost]\label{def:residual}
The 13 $+$ crossings of Definition~\ref{lem:cost13} are partitioned as follows:
\begin{linenomath}
\begin{align}
\text{ball value} &= 4 = \underbrace{1}_{\text{time write}} + \underbrace{3}_{x,y,z\text{ write}}, \\
\text{residual cost} &= 9 = 13 - 4.
\end{align}
\end{linenomath}
The \emph{ball value} $4$ is precisely the \emph{cost consumed in creating one ball} --- time $1$ + space $3$ (the registration cost of a 4-axis coordinate). When the ball value $4$ is recovered, the ball is released from DATA.

The residual cost $9$ is recorded in the RLU index (defined immediately after Axiom~\ref{ax:cas}) and recovered incrementally; it consists of $8$ read crossings and $1$ commit write.
\end{Definition}

\emph{Geometric measure of the residual cost}\label{rem:geometric-measure}. Each $+$ crossing is irreversible by the forward directionality $\mathrm{R} \to \mathrm{C} \to \mathrm{S}$ of Axiom~\ref{ax:cas} (reverse-order crossings are undefined in the axiom). The geometric measure of a single irreversible $+$ crossing is the \emph{hemispherical arc length} $\pi$ on the unit sphere---a reversible (bidirectional) crossing would traverse the full sphere $2\pi$, but only the forward direction is permitted, yielding half. The total geometric arc length of the residual cost $9$ in Definition~\ref{def:residual} is therefore $9\pi$, the additive measure of $9$ independent $+$ crossings (a sum of $1$-dimensional arc lengths).

This measure differs in kind from the $5$-dimensional complex-ball volume $\pi^5/5!$ of $D_5$ (a multiplicative measure)---arc length versus volume are measures of different dimension in differential geometry, and addition versus multiplication are their respective natural combination rules.

Recovery of the residual cost within the RLU proceeds as \emph{geometric decay} along this arc. \emph{Why the geometric form}: objects inside the RLU cannot reference anything outside the RLU; they can reference only \emph{themselves} (self-reference). Self-reference implies that each step depends solely on its own predecessor ($a_{n+1} = r \cdot a_n$), which naturally generates a geometric series. Other decay forms (linear, exponential, etc.)\ would require external references and are therefore incompatible with the self-reference constraint of the RLU. This mechanism is analogous to the LRU (Least Recently Used) cache eviction rule and is employed in greater detail in the RLU coupled-access analysis of Section~\ref{sec:baryogenesis}.

\emph{Relation of this mechanism to the results of the paper}: The geometric-decay mechanism described above serves as an \emph{operational explanation} of the RLU and is not directly used in deriving the quantitative results ($\alpha$, $\sin^2\theta_W$, $\eta_B$) of this paper---it serves only as motivation for the existence of RLU coupled access. This paragraph provides the operational foundation for understanding the RLU mechanism and is independent of the main derivations. The geometric-decay mechanism is not used in the quantitative results of this paper and is presented as an axiom-described result (classification~III).

\begin{Axiom}[$\delta$ is a global flag]\label{ax:delta}
The left-hand side $\delta$ of Eq.~\eqref{eq:banya} lies outside the 7-bit internal machine defined by the right-hand side. When $\delta = 1$ the entire right-hand side is valid; when $\delta = 0$ it is void. The number of internal degrees of freedom is therefore exactly $7$, not $8$.
\end{Axiom}

\textit{Use-case classification --- Raw vs Coupled}\label{rem:stage-classification}. When the two $+$-reading principles of this axiomatic system (see common misconception \#3 in the Reader's note of Section~\ref{sec:framework}) are used to produce physical quantities, two use cases arise:
\begin{itemize}
\item \emph{Raw use case}: $+$ crossings are expressed directly through cost reading and norm reading alone. The resulting form is $n\pi$ (norm reading) or $n + m$ (cost-reading sum) --- e.g., the denominator $9\pi$ of $\sin^2\theta_W$ in Section~\ref{sec:prediction} (norm reading of the residual cost $9$).
\item \emph{Coupled use case}: the RLU coupled-access mechanism, forced by the minimum-cost principle (Section~\ref{sec:framework}), is used in addition to the two readings. The resulting form is $(m + 1/\pi)$ or similar --- e.g., the $\eta_B$ correction $(4 + 1/\pi)$ of Section~\ref{sec:baryogenesis}.
\end{itemize}
The raw use case is satisfied by the two reading principles alone and sits naturally within the forward chain; the coupled use case invokes the additional RLU mechanism and therefore belongs to the axiom-described tier (classification~III; RLU itself is axiomatically required for CAS; the $+$ in $(4+1/\pi)$ is uniquely determined by Axiom~\ref{ax:banya}'s orthogonal composition --- see Section~\ref{sec:prediction} step~(4) and Section~\ref{sec:baryogenesis}; the formal systematization of the full RLU mechanism is classification~III). In this paper, $\sin^2\theta_W$ (raw) and $\eta_B$ (coupled) accordingly carry different epistemic status under this classification.

%%%%%%%%%%%%%%%%%%%%%%%%%%%%%%%%%%%%%%%%%%
\section{Cost--Signature Correspondence Theorem}
\label{sec:theorem}

In this section we state and prove the central result of the paper.

\begin{Definition}[Cost function]
We denote the cost function defined by Axiom~\ref{ax:cost} as $c: \{x_1, \ldots, x_7\} \to \{0, +1\}$. For each axis $x_i$, $c(x_i) = +1$ if the CAS must cross the parenthesis boundary $+$ in sequence when accessing $x_i$, and $c(x_i) = 0$ otherwise.
\end{Definition}

\begin{Definition}[Signature function]
In the quadratic form $Q = \sum_{i=1}^{7} \sigma_i x_i^2$ on a 7-dimensional real vector space, we denote the sign of each axis $\sigma_i \in \{+1, -1\}$ as the signature function $\sigma: \{x_1, \ldots, x_7\} \to \{+1, -1\}$. Let $p$ be the number of positive ($+1$) signs and $q$ the number of negative ($-1$) signs; then $(p,q)$ is the signature of the quadratic form.
\end{Definition}

\begin{Definition}[Cost--signature map]\label{def:cost-sig}
In the 7-dimensional internal space satisfying Axioms~\ref{ax:banya}--\ref{ax:delta}, the signature function $\sigma$ is defined in terms of the cost function $c$ as follows:
\begin{linenomath}
\begin{equation}
\sigma(x_i) = \begin{cases} +1 & \text{if } c(x_i) > 0 \quad (\text{irreversible}) \\ -1 & \text{if } c(x_i) = 0 \quad (\text{reversible}) \end{cases}
\label{eq:cost-sig}
\end{equation}
\end{linenomath}
Physical motivation: irreversible axes have $c > 0$ and are the \emph{data axes} that the axiomatic system \emph{records} --- they are marked with positive ($+$) sign. Reversible axes have $c = 0$ and serve as \emph{filter axes} of the CAS, which the axiomatic system does not record --- they are marked with negative ($-$) sign. The \emph{filter} interpretation of the negative sign is elaborated in part~(B) of the proof below and applied in the derivation of $\sin^2\theta_W$ (Section~\ref{sec:prediction}).
\end{Definition}

\emph{Why a quadratic form}: the standard mathematical device that connects a binary classification of axes (cost present/absent) to a continuous geometric object (symmetric space) is the signature of a quadratic form --- and it is the only one. Given a binary classification $\{c > 0, c = 0\}$, the unique standard path to a continuous group is $Q = \sum \sigma_i x_i^2$ with signature $(p,q)$ $\to$ $\mathrm{SO}(p,q)$. No alternative path exists; Definition~\ref{def:cost-sig} is therefore not a choice but a necessity.

\begin{Theorem}[Uniqueness of the cost partition]\label{thm:cost-signature}
Under Axioms~\ref{ax:banya}--\ref{ax:delta}, the cost classification $c(x_i)$ of the seven axes is determined by the axioms, and the partition into five irreversible axes and two reversible axes, $(5,2)$, is unique. Because rotations between axes of differing cost properties violate the cost-property boundary, the unique maximal compact subgroup preserving this partition is $\mathrm{SO}(5) \times \mathrm{SO}(2)$. The bounded symmetric domain $D_5 = \mathrm{SO}_0(5,2)/[\mathrm{SO}(5) \times \mathrm{SO}(2)]$ (where $\mathrm{SO}_0(5,2)$ is the identity component) is thereby determined.
\end{Theorem}

\emph{Note on the status of the definition and theorem}: Definition~\ref{def:cost-sig} (cost--signature map) assigns positive ($+$) sign to irreversible axes and negative ($-$) sign to reversible axes. Since $\mathrm{SO}(5,2) \cong \mathrm{SO}(2,5)$, reversing the sign convention yields the same $D_5$. Theorem~\ref{thm:cost-signature} establishes the uniqueness of the \emph{partition} $(5,2)$; this partition alone determines $D_5$ and $\alpha$.

\emph{Formal status of the theorem}: This theorem applies the constraints of Axioms~\ref{ax:banya}--\ref{ax:delta} to the seven axes and shows that the partition $(5,2)$ is forced. Because the axiomatic constraints are strong, only one of the eight possible partitions survives --- this is the essence of axiom-derived classification (Class~I).

\begin{proof}
The proof is divided into three parts: (A) uniqueness of the cost classification, (B) correspondence between cost and metric sign, and (C) formal exclusion of alternative partitions. (The well-definedness of the map $\sigma$ follows directly from (A) and (C) and is addressed in a remark after the proof.)

\emph{(A) Uniqueness of the cost classification.} We determine $c(x_i)$ for each of the seven axes:
\begin{itemize}
\item \textbf{time, space}: By Axiom~\ref{ax:banya}, these belong to the DATA parenthesis. Because the CAS operates in the OPERATOR parenthesis (Axiom~\ref{ax:cas}), accessing DATA requires crossing the parenthesis boundary $+$. By Axiom~\ref{ax:cost}, $c = +1$. Irreversible.
\item \textbf{R, C, S} (the three CAS stages, each an independent binary degree of freedom): By Axiom~\ref{ax:cas} the ordering $\mathrm{R} \to \mathrm{C} \to \mathrm{S}$ is forced by logical dependence. Each stage occupies an independent binary degree of freedom (001, 011, 111), and each order-forced transition has $c = +1$ by Axiom~\ref{ax:cost}. Irreversible.
\item \textbf{observer, superposition}: By Axiom~\ref{ax:banya}, these belong to the OPERATOR parenthesis. Because the CAS also operates in OPERATOR (Axiom~\ref{ax:cas}), access within the same parenthesis does not cross $+$. By Axiom~\ref{ax:cost}, $c = 0$. Reversible.
\end{itemize}
The classification of each axis is completely determined by the parenthesis to which the axis belongs (Axiom~\ref{ax:banya}) and the location of the CAS (Axiom~\ref{ax:cas}). Because no axis can be moved to a different parenthesis, the classification is unique. Result: five axes with $c > 0$ and two axes with $c = 0$.

\emph{(B) Correspondence between cost and metric sign (data/filter category markers).} In the quadratic form $Q = \sum \sigma_i x_i^2$, the physical meaning of the sign $\sigma_i$ is that it designates whether the axis is a \emph{data dimension} or a \emph{filter dimension}.

In the present axiomatic system the cost function $c$ directly determines this distinction:
\begin{itemize}
\item \textbf{Irreversible axes} ($c > 0$): data dimensions that the axiomatic system \emph{records}. The Swap result of the CAS is permanently registered in the DATA slot. The time and space axes are the targets of timestamp and coordinate recording, respectively, and the R/C/S CAS stages are the irreversible steps that \emph{carry out} that recording. Hence $\sigma_i = +1$ --- the category marker for data dimensions.
\item \emph{Reversible axes} ($c = 0$): dimensions that the axiomatic system \emph{does not record} and that serve as \emph{filter/selector} roles of the CAS. The observer domain acts as the CAS \emph{branch filter}, and the superposition domain acts as the CAS \emph{RLU index filter} (see the RLU definition immediately following Axiom~\ref{ax:cas}) --- both merely select \emph{which data to fetch or which branch to take} without themselves being recorded as data. Hence $\sigma_i = -1$ --- the category marker for filter dimensions.
\end{itemize}

We emphasize the key point: \emph{the negative ($-$) sign is not ``subtraction'' or ``monotonic decrease'' but the category marker for ``filter dimension''}. $(-3)^2 = (+3)^2 = 9$ so that the absolute value of $\sigma_i x_i^2$ is the same regardless of sign, and the sign is merely a \emph{type marker} distinguishing data dimensions ($p$) from filter dimensions ($q$) in the SO($p,q$) group structure, carrying no arithmetic meaning. In the derivation of $\sin^2\theta_W$ (Section~\ref{sec:prediction}, ``set complement $=$ filter'' paragraph) the negative-sign interpretation is likewise used as a \emph{set-filter} operation, consistent with the $(5,2)$ sign assignment of the present theorem --- in both cases the negative sign means ``filter/selection,'' not ``subtraction.''

Comparison with Minkowski spacetime. In the Minkowski metric, two different categories (time vs.\ space) also receive different signs, but the $(5, 2)$ sign assignment of the present axiomatic system arises from the \emph{cost category} (data vs.\ filter), not from a kinematic causal structure. The negative sign in the present axiomatic system marks ``filter dimensions'' (reversible, $c = 0$) and differs in meaning from the spatial negative sign in Minkowski space --- the similarity is analogical, not identical.

\emph{(C) Formal exclusion of alternative partitions.} We show that the partition $(5,2)$ obtained in (A) is \emph{unique} by explicit comparison with every other possible partition. For seven axes the partition is $(p, q)$ with $p + q = 7$, giving eight possibilities: $(7,0), (6,1), (5,2), (4,3), (3,4), (2,5), (1,6), (0,7)$. Apart from $(5,2)$ and its isomorphic counterpart $(2,5)$, all remaining six are excluded by the axioms of the axiomatic system:
\begin{center}
\small
\begin{tabular}{ccp{8.5cm}}
\toprule
\textbf{Partition} & \textbf{Verdict} & \textbf{Exclusion rationale} \\
\midrule
$(7,0)$ & Excluded & observer, superposition $\in$ OPERATOR $\Rightarrow$ no boundary crossing $\Rightarrow$ $c=0$ forced. Violates Axioms~\ref{ax:banya} and \ref{ax:cost} \\
$(0,7)$ & Excluded & time, space $\in$ DATA, CAS $\in$ OPERATOR $\Rightarrow$ boundary must be crossed $\Rightarrow$ $c=+1$ forced. Violates Axioms~\ref{ax:banya} and \ref{ax:cost} \\
$(6,1)$ & Excluded & Requires an OPERATOR axis to have $c>0$, but (A) fixes $c=0$. Axiom violation \\
$(4,3)$ & Excluded & Requires a DATA axis or CAS stage to have $c=0$, but (A) fixes $c=+1$. Axiom violation \\
$(3,4)$ & Excluded & Same logic as $(4,3)$ \\
$(1,6)$ & Excluded & Same logic as $(6,1)$ \\
$(5,2)$ & \textbf{Adopted} & Uniquely compatible with the cost classification of (A). Includes isomorphic $(2,5)$ \\
\bottomrule
\end{tabular}
\end{center}
Result: of the eight possible partitions, the only one compatible with Axioms~\ref{ax:banya}, \ref{ax:cas}, and \ref{ax:cost} of the axiomatic system is $(5,2)$ (including its isomorphic counterpart $(2,5)$). Since $\mathrm{SO}(5,2) \cong \mathrm{SO}(2,5)$, the choice of sign direction does not affect the result, and what this theorem establishes is the \emph{uniqueness of the partition itself}.

From (A), (B), and (C) the partition is determined to be five irreversible and two reversible axes, giving $(p, q) = (5, 2)$.
\end{proof}

\emph{Remark (well-definedness of the map)}. The map $\sigma: c \to \{+1, -1\}$ of Definition~\ref{def:cost-sig} possesses the following three formal properties as a direct consequence of (A) and (C): (i) \emph{Total}: the map is defined on all seven axes; (ii) \emph{Deterministic}: given the $c$ value of each axis, the $\sigma$ value is determined (no external input); (iii) \emph{Count-bijective}: the resulting partition $(p,q)$ is determined by the number of $c=+1$ axes and $c=0$ axes. ``Structural equivalence between the cost classification and the quadratic-form signature'' refers to these three properties.

\emph{Exclusion of alternative 21-dimensional groups: a categorical distinction}. SO(5,2) is a noncompact simple Lie group of dimension 21, but other simple Lie groups of the same dimension 21 exist --- SO(7) (the compact form, signature $(7,0)$) and SO(4,3) (signature $(4,3)$ or $(3,4)$). Within the present axiomatic system the distinction among these three reduces to the question of which \emph{category of degrees of freedom} (Proposition~\ref{ax:dof}: structural vs.\ cost) the metric signature is derived from.

\begin{itemize}
\item \emph{SO(5,2) is derived from \emph{cost degrees of freedom}}: Definition~\ref{def:cost-sig} specifies that the cost function produces two distinct categories ($c > 0$ irreversible and $c = 0$ reversible) and assigns different signs to them. Because Theorem~\ref{thm:cost-signature} derives the sizes of the two categories as 5 and 2, the signature $(5,2)$ is forced.

\item \emph{SO(7) is the compact form of \emph{structural degrees of freedom}}: SO(7) is the compact group in which all seven axes carry the same sign, treating the seven axes as \emph{structurally equivalent} objects. However, the present axiomatic system derives the metric signature from the \emph{cost category}, not from the structural category, and the cost category forces two different signs (see item above). SO(7) is therefore incompatible with the metric-signature derivation principle of the axiomatic system.

\item \emph{SO(4,3) is the split form of \emph{structural degrees of freedom}}: The partition $(4,3)$ required for SO(4,3) coincides numerically with the structural decomposition of Proposition~\ref{ax:dof} (4 domain axes $+$ 3 CAS stages $= 7$), but this is a partition of the \emph{structural} category, not the \emph{cost} category. Because the axiomatic system derives the metric signature from the cost category (Definition~\ref{def:cost-sig} and Theorem~\ref{thm:cost-signature}), the structural partition $(4,3)$ is not an input to the signature derivation. Since Theorem~\ref{thm:cost-signature} forces the \emph{cost} partition $(5,2)$, SO(4,3) is incompatible with the axiomatic system.
\end{itemize}

Among 21-dimensional simple Lie groups, the only one compatible with the \emph{signature-from-cost} principle of the axiomatic system is therefore SO(5,2) (including its isomorphic counterpart SO(2,5)). SO(7) and SO(4,3) are objects belonging to a different category (structural) within the axiomatic system and are not used in the metric-signature derivation itself.

%%%%%%%%%%%%%%%%%%%%%%%%%%%%%%%%%%%%%%%%%%
\section{Derivation of $\alpha$}
\label{sec:derivation}

Once the signature $(5,2)$ is fixed by Theorem~\ref{thm:cost-signature}, the chain $\mathrm{SO}_0(5,2) \to D_5 \to$ the factors of the Wyler formula \eqref{eq:wyler-formula} follows. Table~\ref{tab:factors} classifies the axiomatic status of each factor: $9$ is axiom-derived (I), $\pi^5/(2^4 \cdot 5!)$ is axiom-determined (II), $(\cdot)^{1/4}$ is axiom-determined (II) as the geometric mean of 4 independent axes, and $8\pi^4$ is the group volume determined by $(5,2)$ and hence mathematically axiom-determined (II). Their epistemic status differs (see Table~\ref{tab:factors}).

\subsection{SO(5,2) and the Maximal Compact Subgroup}

The group of linear transformations preserving the quadratic form of signature $(5,2)$ is $\mathrm{SO}(5,2)$, and its identity component $\mathrm{SO}_0(5,2)$ is the natural transformation group of the bounded symmetric domain $D_5$~\cite{helgason1978}. The maximal compact subgroup is $\mathrm{SO}(5) \times \mathrm{SO}(2)$.

Why this subgroup is uniquely selected: the cost classification partitions the 7 axes into $\{$5 irreversible$\} \cup \{$2 reversible$\}$. The 5 irreversible axes share the same property (cost $> 0$) and are therefore mutually rotatable---this is $\mathrm{SO}(5)$. The 2 reversible axes likewise share the same property (cost $= 0$) and are mutually rotatable---this is $\mathrm{SO}(2)$. However, rotations between irreversible and reversible axes are forbidden because their cost properties differ (cost $> 0$ vs.\ cost $= 0$). Consequently, the maximal compact subgroup is $\mathrm{SO}(5) \times \mathrm{SO}(2)$, and alternative subgroups such as $\mathrm{U}(5)$ or $\mathrm{Sp}(4) \times \mathrm{U}(1)$ are excluded because they violate the cost-property boundary.

\subsection{Quotient Space $D_5$}

The quotient space
\begin{linenomath}
\begin{equation}
D_5 = \frac{\mathrm{SO}_0(5,2)}{\mathrm{SO}(5) \times \mathrm{SO}(2)}
\label{eq:d5}
\end{equation}
\end{linenomath}
is a Cartan type IV bounded symmetric domain of complex dimension 5~\cite{helgason1978,hua1963}. Forming the quotient by the internal rotations within each sector leaves only the transformations that \emph{connect} the irreversible and reversible sectors---that is, transformations that cross the bracket boundary $+$. $D_5$ is therefore ``the space of all possible configurations that incur cost.''

\subsection{Wyler Volume Ratio}

The ratio of the Shilov boundary volume to the $D_5$ volume determines $\alpha$~\cite{wyler1969,wyler1971,hua1963}:
\begin{linenomath}
\begin{equation}
\alpha = \frac{9}{8\pi^4} \left( \frac{\pi^5}{2^4 \cdot 5!} \right)^{1/4}.
\label{eq:wyler-formula}
\end{equation}
\end{linenomath}

Physically, $\alpha$ is the probability that a CAS operation crosses the bracket boundary $+$ and produces an interaction. The ratio of realized configurations (Shilov boundary) to total possible configurations ($D_5$) yields the low-energy electromagnetic coupling constant.

\emph{Note on uniqueness of the Shilov boundary}: The Shilov boundary is uniquely defined as the minimal closed subset of a bounded domain on which the maximum principle holds~\cite{hua1963}. Alternative boundary concepts such as the Bergman boundary or the distinguished boundary exist; however, the fact that the coupling constant in Wyler's~\cite{wyler1969} original construction is given by the Shilov boundary volume ratio is consistent with the physical meaning of this boundary as ``the minimal set on which the values of holomorphic functions are determined''---that is, the minimal boundary that specifies the realized configurations. The present paper independently determines the signature $(5,2)$ (Theorem~\ref{thm:cost-signature}); the Shilov boundary is the unique minimal boundary on which holomorphic functions attain their maxima~\cite{hua1963} and is therefore not an arbitrary choice but a mathematical necessity.

\textit{Why the Shilov boundary volume ratio equals $\alpha$---the self-descriptive structure of the Banya equation.} This question has remained open since Wyler (1969) and is the core of Robertson's criticism R1 (``Why should $D_5$ be related to electromagnetism?''). In the present axiomatic system, this identity is not an additional assumption but a consequence of the self-descriptive structure of Axiom~\ref{ax:banya} (the Banya equation).

In the Banya equation $\delta^2 = (\mathrm{time} + \mathrm{space})^2 + (\mathrm{observer} + \mathrm{superposition})^2$, the left-hand side $\delta$ is a single quantity of change. This $\delta$ can be described in the structural domain (as a geometric volume ratio) or in the cost domain (as a coupling constant). The two descriptions are expressions in different domains but refer to the same $\delta$, and therefore necessarily yield the same numerical value. This is the principle of domain transformation.

Concretely: (i)~In the geometric domain, the Shilov boundary volume ratio of $D_5$, determined by $(5,2)$, is ``the fraction of possible configurations in which cost is incurred'' (§\ref{sec:derivation}). (ii)~In the cost domain, $\alpha$ is ``the probability that a CAS operation crosses the bracket boundary'' (§\ref{sec:derivation}). (iii)~Both descriptions are projections of the same $\delta$ (quantity of change) onto different domains. Their numerical identity is therefore not a coincidence but a necessary consequence of the self-descriptive structure of the Banya equation.

This structure reflects the methodology of the axiomatic system itself: the Banya equation is a tool for viewing the same object from multiple domains, and the value of $\delta$ is preserved across domain transformations. This is the principle of ``axiom mining''---discovering hidden information through domain transformation and substitution. The identity of the Shilov boundary volume ratio with $\alpha$ is one instance of this conservation principle.

\subsection{Numerical Result}

\begin{linenomath}
\begin{equation}
\frac{1}{\alpha} = 137.036\,082.
\label{eq:result}
\end{equation}
\end{linenomath}

CODATA 2022 experimental value~\cite{codata2022}: $1/\alpha_{\mathrm{exp}} = 137.035\,999\,177(21)$.

\begin{linenomath}
\begin{equation}
\frac{|\alpha^{-1}_{\mathrm{derived}} - \alpha^{-1}_{\mathrm{exp}}|}{\alpha^{-1}_{\mathrm{exp}}} = 6 \times 10^{-7}.
\label{eq:error}
\end{equation}
\end{linenomath}

\subsection{Factor Tracing}

Every numerical factor in the Wyler formula is determined by the signature $(5,2)$:

\begin{table}[H]
\begin{adjustwidth}{-\extralength}{0cm}
\caption{Axiomatic status of each factor in the Wyler formula.\label{tab:factors}}
\small
\begin{tabularx}{\fulllength}{cp{2.2cm}X}
\toprule
\textbf{Class.} & \textbf{Factor} & \textbf{Remarks} \\
\midrule
\textbf{I} & $9$ & Residual cost $13{-}4{=}9$. Analytically corresponds to dim\,SO(5)$-$dim\,SO(2)${}=10{-}1$. The unique factor produced by the axiomatic system (Definition~\ref{def:residual}) \\
\textbf{II} & $\pi^5/(2^4 \cdot 5!)$ & Hua volume of $D_5$. Theorem~\ref{thm:cost-signature}${}\to{}$(5,2)${}\to{}D_5{}\to{}$Hua Ch.\,IV identity. Determined by mathematical identity (Section~\ref{sec:derivation}) \\
\textbf{II} & $(\cdot)^{1/4}$ & Axiom~\ref{ax:banya} orthogonality $\to$ independence $\to$ geometric mean $= (\cdot)^{1/4}$. Derivation complete (Sections~\ref{sec:derivation},~\ref{sec:prediction}) \\
\textbf{II} & $8\pi^4$ & Mathematically axiom-determined from the $(5,2) \to \mathrm{SO}(5)\times\mathrm{SO}(2)$ group volume (Section~\ref{sec:derivation}) \\
\bottomrule
\end{tabularx}
\parbox{\fulllength}{All four factors are accounted for by the axiomatic structure. For the classification criteria, see Table~\ref{tab:axiom-results}.}
\end{adjustwidth}
\end{table}

\subsubsection{Forward Chain of $\pi^5/5!$ (Strong Result)}

\emph{The exponent $5$ in $\pi^5$ equals the number of irreversible axes.} This is not a coincidence but a definitional consequence of the Cartan classification. We make the derivation chain explicit:
\begin{enumerate}
\item By Theorem~\ref{thm:cost-signature} (uniqueness of the cost partition), the 7 axes are partitioned into 5 irreversible $+$ 2 reversible.
\item This partition fixes the quadratic-form signature $(p, q) = (5, 2)$.
\item The maximal compact subgroup of $\mathrm{SO}(p, q) = \mathrm{SO}(5, 2)$ is $\mathrm{SO}(5) \times \mathrm{SO}(2)$.
\item The quotient space $D_p = D_5 = \mathrm{SO}_0(5,2)/[\mathrm{SO}(5)\times\mathrm{SO}(2)]$ (where $\mathrm{SO}_0(5,2)$ is the identity component of $\mathrm{SO}(5,2)$) is a Cartan type IV bounded symmetric domain, and \emph{its complex dimension is exactly $n = 5$} (Helgason~\cite{helgason1978}, Hua~\cite{hua1963} Ch.~IV).
\item Substituting $n = 5$ into the type IV$_n$ Hua volume formula from Hua~\cite{hua1963} Ch.~IV, $\mathrm{Vol}_{\text{Hua}}(D_n^{\text{IV}}) = \pi^n/(2^{n-1} n!)$, yields $\pi^5/(2^4 \cdot 5!) = \pi^5/1920$ (Eq.~\eqref{eq:hua-volume}).
\end{enumerate}

At each step, ``$5$'' is the same number---\emph{the number of irreversible axes}---propagated through mathematical abstraction (signature $\to$ group $\to$ quotient space $\to$ volume). Accordingly, the $\pi^5$ in the Wyler formula is not a purely geometric constant; within the present axiomatic system it admits a direct physical interpretation as \emph{the unit-sphere volume of the complex 5-dimensional space collectively spanned by the 5 irreversible axes}. It is not the case that each complex dimension individually contributes a factor of $\pi$; rather, all 5 axes collectively form a single 5-dimensional space whose volume is $\pi^5 / 5!$. The remaining factors $9$, $8\pi^4$, and $2^4$ in the table are likewise described within the axiomatic structure, through Theorem~\ref{thm:cost-signature} and its automatic group-theoretic consequences.

\emph{$\pi^5/(2^4 \cdot 5!) = D_5$ Hua volume: direct identity (Hua Ch.~IV)}. The factor $\pi^5/(2^4 \cdot 5!)$ in the Wyler formula~\eqref{eq:wyler-formula} is \emph{not two separate factors} ($\pi^5/5!$ and $2^4$) but a \emph{single object}: the Hua volume of $D_5$. (``Hua volume'' follows Hua~\cite{hua1963} Ch.~IV; the Wyler--Robertson literature~\cite{wyler1969,robertson1971} calls it ``Bergman volume.'')

The Hua volume formula for the type-IV$_n$ Lie ball $D_n^{\text{IV}} = \mathrm{SO}_0(n,2)/[\mathrm{SO}(n)\times\mathrm{SO}(2)]$ is
\begin{linenomath}
\begin{equation}
\mathrm{Vol}_{\text{Hua}}(D_n^{\text{IV}}) = \frac{\pi^n}{2^{n-1} \cdot n!},
\label{eq:hua-volume}
\end{equation}
\end{linenomath}
and substituting $n=5$ gives $\pi^5/(2^4 \cdot 5!) = \pi^5/1920$, identical to the Wyler factor. Once the axioms force $(5,2) \to D_5$ via Theorem~\ref{thm:cost-signature}, this volume is a mathematical identity of standard Lie group theory (Hua Ch.~IV, Helgason~\cite{helgason1978}) and is automatically determined (classification~II).

\textit{Axiomatic status of the four factors.} Each factor in the Wyler formula~\eqref{eq:wyler-formula} receives an account within the axiomatic structure, at differing epistemic levels (Table~\ref{tab:factors}): $9$ is axiom-derived from the cost accounting (I), $\pi^5/(2^4 \cdot 5!)$ is axiom-determined as the $D_5$ Hua volume (II), $8\pi^4$ is mathematically axiom-determined from the $(5,2) \to \mathrm{SO}(5)\times\mathrm{SO}(2)$ group volume (II), and $(\cdot)^{1/4}$ is axiom-determined as the geometric mean of 4 orthogonal independent axes (II).

\subsubsection{Status of the $(\cdot)^{1/4}$ Power --- Geometric Mean of Orthogonal Independent Axes}

\emph{Status of the $1/4$ power}. Axiom~\ref{ax:banya} declares the mutual orthogonality of the 4 domain axes. Orthogonality $=$ independence (as stated in the Hartley information paragraph of Section~\ref{sec:framework}). For $n$ independent axes, the per-axis contribution is the geometric mean $= (\text{total})^{1/n}$ --- this is a mathematical identity following from the definition of independence, not an assumption. With $n = 4$ (domain axis count), the $1/4$ power is mathematically determined by the orthogonality of Axiom~\ref{ax:banya} (classification~II). The $1/4$ adopted by Wyler on dimensional-analysis grounds is independently derived in the present axiomatic system.

\emph{Ratio interpretation: $\alpha$ = cost / structure (power division).} By Proposition~\ref{ax:dof} (complete descriptive degrees of freedom), the degrees of freedom of the axiomatic system split into two orthogonal categories: \emph{structure} (DATA category, $\{1, 2, 3, 4, 7\}$) and \emph{cost} (OPERATOR category, $\{1, 4, 5, 9, 13\}$). The 4 domain axes (Axiom~\ref{ax:banya}) correspond to the \emph{4 of the structure category}; since these 4 axes are orthogonal (independent), they combine via geometric mean; the total cost accumulated across 4 axes is the per-axis cost raised to the 4th power, and conversely the per-axis cost is the total cost raised to the $1/4$ power:
\begin{linenomath}
\begin{equation}
\alpha = (\text{cumulative cost over 4 axes})^{1/\text{structure 4}}.
\label{eq:cost-ratio}
\end{equation}
\end{linenomath}
From this perspective, the $(\cdot)^{1/4}$ power in the Wyler formula is \emph{the inverse operation of the multiplicative combination of 4 domain axes}, and other powers ($1/2$, $1/3$, $1/5$, $\ldots$) are inconsistent with the axiomatic system's domain axis count of $4$. The $1/4$ power adopted by Wyler on dimensional-analysis grounds is \emph{independently rediscovered} in the cost ratio interpretation of the present axiomatic system, and the convergence of two independent paths weakens the hypothesis that this power is accidental (the $\alpha^4$ factor interpretation immediately after Eq.~\eqref{eq:eta-B} in Section~\ref{sec:baryogenesis} also rests on the same 4-domain-axis geometric-mean identity).

\emph{Categorical distinction: why 4-fold and not 7-fold}. The total internal DOF count of the present axiomatic system is $7 = 4 + 3$ (Proposition~\ref{ax:dof}), but the geometric-mean combination above applies only to the \emph{target space dimensions} (4 domain axes) and not to the \emph{access operations} (3 CAS steps, which combine by sequential summation).

The 4 domain axes (Axiom~\ref{ax:banya}) are the CAS's \emph{access targets} (target space) and candidates for multiplicative combination --- the observer domain is the target of CAS branch filtering, the superposition domain is the target of CAS RLU indexing (RLU definition immediately after Axiom~\ref{ax:cas}), and the time and space domains are the targets of CAS writing. In contrast, the 3 CAS steps (Axiom~\ref{ax:cas}) combine by sequential summation (this is why the cost enumeration of Definition~\ref{lem:cost13} counts $13 = 8\text{ reads} + 5\text{ writes}$ additively).

Consequently, the ``7-fold alternative $(\cdot)^{1/7}$'' is a category error that conflates two different composition laws, and the $n$ in $(\cdot)^{1/n}$ is naturally restricted to the \emph{number of target space dimensions} ($= 4$ domain axes).

\textit{Mathematical basis of the multiplicative combination.} The Banya equation (Axiom~\ref{ax:banya}) has an additive norm ($\delta^2 = (\cdot)^2 + (\cdot)^2$), but volumes are products of independent dimensions---these are different tiers. Orthogonality $=$ independence, and the per-dimension contribution to a volume of $n$ independent dimensions is the geometric mean $= (\cdot)^{1/n}$. This is a mathematical property of orthogonal spaces, not an additional assumption. With $n = 4$ (domain axis count), $(\cdot)^{1/4}$ is mathematically determined by the orthogonality of Axiom~\ref{ax:banya} (classification~II). Wyler's $2^4$ is not a separate factor but part of the Hua volume $\pi^n/(2^{n-1} \cdot n!)$: substituting $n=5$ yields $2^{5-1} = 2^4$ automatically (Eq.~\eqref{eq:hua-volume} and the ``$\pi^5/(2^4 \cdot 5!) = D_5$ Hua volume direct identity'' paragraph above). A separate mapping of $2^4$ is therefore unnecessary.

\subsubsection{Summary of the Derivation Chain}

The status of the Wyler formula factors within the axiomatic system is classified as follows.

(i) $9 = \dim\mathrm{SO}(5)-\dim\mathrm{SO}(2)$ analytically corresponds to the residual cost from the axiomatic system's cost accounting (axiom-derived).

(ii) $\pi^5/(2^4 \cdot 5!)$ is the $D_5$ Hua volume (Eq.~\eqref{eq:hua-volume}), axiom-determined once $(5,2) \to D_5$ is forced (classification~II).

(iii) The denominator $8\pi^4$ is mathematically axiom-determined from the $(5,2) \to \mathrm{SO}(5)\times\mathrm{SO}(2)$ group volume (classification~II).

(iv) The $(\cdot)^{1/4}$ power is the geometric mean of 4 orthogonal independent axes (Axiom~\ref{ax:banya}), mathematically axiom-determined (classification~II). Derivation chain:
\begin{linenomath}
\begin{equation}
\text{Axioms~1--4 + Proposition~1} \;\xrightarrow{\text{Theorem~1}}\; (5,2) \;\xrightarrow{\text{group theory}}\; D_5 \;\xrightarrow{\text{Wyler}}\; \alpha.
\end{equation}
\end{linenomath}

%%%%%%%%%%%%%%%%%%%%%%%%%%%%%%%%%%%%%%%%%%
\section{Structural Consistency Check: Weinberg Angle}
\label{sec:prediction}

As a check of internal consistency, we derive an independent second physical constant from the same axiomatic system. Here we derive the Weinberg angle $\sin^2\theta_W$.

\subsection{Main Result: Minimal Form}

The forward derivation of $\sin^2\theta_W$ from the present axiomatic system yields Eq.~\eqref{eq:weinberg-main}:
\begin{linenomath}
\begin{equation}
\boxed{\sin^2\theta_W = \frac{7}{2 + 9\pi} = 0.23122.}
\label{eq:weinberg-main}
\end{equation}
\end{linenomath}

The origins of the four structural constants are summarized in Table~\ref{tab:weinberg-constants}. The $+$ in the denominator ($2 + 9\pi$) is an \emph{orthogonal composition} of the reversible~2 (cost reading) and the irreversible~$9\pi$ (norm reading)---matching Axiom~\ref{ax:banya}'s orthogonal sum. No double counting. This combination rule of cost reading and norm reading is formalized in the complete axiomatic system~\cite{han2026} as the \emph{Axiom~4 proposition ``Reading Cost in Norm''} (v1.6) (cost $\perp$ norm orthogonal sum is expressed as standard arithmetic addition, consistent with the framework's $+$).

\begin{table}[H]
\caption{Four structural constants in $\sin^2\theta_W = 7/(2+9\pi)$.\label{tab:weinberg-constants}}
\small
\begin{tabular}{clll}
\toprule
\textbf{Symbol} & \textbf{Value} & \textbf{Origin} & \textbf{Reading type} \\
\midrule
$7$ & Internal DOF & Theorem~\ref{thm:cost-signature}: $5$ irrev.\ $+$ $2$ rev. & Cost reading (integer) \\
$2$ & Reversible axes & Theorem~\ref{thm:cost-signature}: observer, superposition & Cost reading (integer) \\
$9$ & Residual cost & Definition~\ref{def:residual}: $13-4$ (irreversible) & Cost reading (integer) \\
$\pi$ & Hemispheric arc & Axiom~\ref{ax:cost}$+$\ref{ax:cas}: irrev.\ $\to$ forward only & Norm reading (geometric) \\
\bottomrule
\end{tabular}
\end{table}

Experimental value at the $M_Z$ scale ($\overline{\mathrm{MS}}$ scheme, PDG 2024 ``Electroweak model and constraints on new physics'' section, $\hat s_Z^2 = \sin^2\hat\theta_W(M_Z)_{\overline{\mathrm{MS}}}$): $\sin^2\theta_W = 0.23122 \pm 0.00004$~\cite{pdg2024}. Experimental uncertainty: $0.017\%$; theory--experiment relative deviation: $4 \times 10^{-6}$. We emphasize that the comparison is restricted to a single definition---$\overline{\mathrm{MS}}$ at $M_Z$---rather than the \emph{on-shell} definition $1 - M_W^2/M_Z^2 \approx 0.22337$ or the \emph{effective leptonic} definition $\sin^2\theta_{\text{eff}}^{\text{lept}} \approx 0.23155$. The other fit rows of the ``$\sin^2\theta$ values'' table in PDG 2024 ($\bar s_Z^2 \approx 0.23129$ all-data fit, $0.23122$ LHC fit, etc.) are \emph{effective leptonic} definitions and therefore not the comparison target ($\hat s_Z^2$ and $\bar s_Z^2$ are distinct definitions).

\textit{Note on energy scale and renormalization scheme.}
The experimental value of $\sin^2\theta_W$ depends on the energy scale and the renormalization scheme---the on-shell definition gives approximately $0.2234$, whereas $\overline{\mathrm{MS}}$ at $M_Z$ gives approximately $0.23122$, a difference of about $3.5\%$ between the two. The numerical value $0.23122$ of the present formula agrees with the latter ($\overline{\mathrm{MS}}$ at $M_Z$). We organize in four items how this agreement is situated within the present axiomatic system.

\begin{enumerate}
\item \emph{The axiomatic system contains no energy scale---a statement of fact.} Axioms~\ref{ax:banya}--\ref{ax:delta} of this paper contain no dimensionful quantities such as energy, mass, or length. Costs are dimensionless integers (Axiom~\ref{ax:cost}), and the structural constants are likewise dimensionless (Proposition~\ref{ax:dof}, Definition~\ref{lem:cost13}, Definition~\ref{def:residual}). All quantities produced by the axioms are therefore dimensionless ratios that precede any energy scale.

\item \emph{System time vs.\ domain time---an interpretive placement.} In the present axiomatic system a single CAS cycle (rest $\to$ R $\to$ C $\to$ S $\to$ DATA commit) is one tick of ``system time,'' and the 13 boundary crossings of Definition~\ref{lem:cost13} occur within this single tick. By contrast, the ``energy scale'' that we measure (e.g.\ $M_Z = 91.2$~GeV) is a quantity defined on \emph{domain time}, after rendering into DATA (classical bracket). The two temporal layers operate at different granularities, and the present axiomatic system operates exclusively at the system-time layer.

\item \emph{Agreement with $\overline{\mathrm{MS}}$ at $M_Z$.} Which measurement scheme on the domain-time side the structural ratio $7/(2+9\pi)$ produced at the system-time layer will agree with does not follow automatically from the axioms. The alignment is natural: the $\overline{\mathrm{MS}}$ scheme is a ``threshold-free'' definition that explicitly removes physical particle-mass thresholds, placing it \emph{aligned with} the system-time layer of the present axiomatic system (where the concept of threshold is absent). The on-shell definition, which incorporates physical particle masses as thresholds, lies closer to the domain-time side. That the present formula agrees with $\overline{\mathrm{MS}}$ at $M_Z$ rather than the on-shell definition is consistent with this placement. Because Axiom~3 of the full axiomatic system~\cite{han2026} explicitly declares DATA (discrete) / OPERATOR (continuous), the system/domain time separation is an immediate consequence of Axiom~\ref{ax:banya} (classification~I; see Section~\ref{sec:axioms}). Axiom~3 of the full system~\cite{han2026} further specifies this separation.

\item \emph{System $\to$ domain time map.} The separation of system time ($\delta$) and domain time (time axis) is an immediate consequence of Axiom~\ref{ax:banya} (Section~\ref{sec:axioms}). Specific correspondences in physical equations are obtained by substituting known physical quantities into the Banya equation.
\end{enumerate}

The result is the fact that ``a dimensionless ratio produced by the axiomatic system agrees to four significant figures with one definition ($\overline{\mathrm{MS}}$ at $M_Z$) of the measured value.''

\subsection{Structural Uniqueness of the Formula}

We answer, from within the axiomatic system, the question ``Why this particular combination $7/(2+9\pi)$?'' In five steps we show that the degrees of freedom in choosing operations are narrowed by the two reading principles of the axiomatic system together with standard arithmetic. \emph{Steps (1)--(3)} proceed automatically: orthogonal structure $\to$ unit sphere $\to$ irreversible hemispheric arc $\to$ norm reading $9\pi$ of the residual cost~9. \emph{Step~(4)} uses orthogonal composition~$+$ as the combination rule for the two reading principles (cost reading vs.\ norm reading)---this combination rule follows from applying Axiom~\ref{ax:banya} (norm reading of $+$) and Axiom~\ref{ax:cost} (cost reading of $+$) to the same structural symbol. \emph{Step~(5)} is the natural consequence of applying the ratio (= division) of standard arithmetic.

\emph{(1) $\pi$ is unique.} The three-axis orthogonality of the CAS (Axiom~\ref{ax:cas}: $\mathrm{R} \perp \mathrm{C} \perp \mathrm{S}$) defines a Euclidean inner-product space. The only transcendental constant that can arise from the unit sphere of a finite-dimensional real inner-product space is $\pi$---this is a direct consequence of the sphere-measure formula $\mathrm{Vol}(S^{n-1}) = 2\pi^{n/2}/\Gamma(n/2)$. Other transcendental constants ($e$, $\ln 2$, etc.) cannot emerge from an orthogonal norm structure.

\emph{(2) The hemisphere ($\pi$, not $2\pi$) is unique.} As shown in the ``geometric measure of the residual cost'' paragraph immediately following Definition~\ref{lem:cost13} and Definition~\ref{def:residual}, each $+$ crossing is irreversible by the forward-directedness of Axiom~\ref{ax:cas}, and the arc length of a single irreversible crossing is the hemisphere~$\pi$ (were it reversible the full sphere $2\pi$ would apply, but only the forward direction is permitted, hence half). $\pi/2$ is a single-axis partial crossing and therefore not the complete measure of one $+$ crossing. The geometric cost of a single $+$ crossing is thus exactly~$\pi$.

\emph{(3) $9 \times \pi$ is unique.} The residual cost~$9$ (Definition~\ref{def:residual}) consists of 9 independent irreversible $+$ crossings ($8$~reads + $1$~commit). The arc length of each crossing is $\pi$ by step~(2). The combination of $9$ independent arcs is the sum of arc lengths ($1$-dimensional measure), giving $9\pi$.

This is a \emph{geometrically distinct quantity} from the $5$-dimensional complex ball volume $\pi^5/5!$ (multiplicative measure) of $D_5$ that appears in the derivation of~$\alpha$---arc length is a $1$-dimensional measure for which addition is natural, whereas volume is a dimensional product measure for which multiplication is natural. In differential geometry the two combination rules (addition / multiplication) are determined by the dimensionality of the measure, and it is not a contradiction that the same axiomatic system, when applied to \emph{different geometric objects}, gives rise to different combination rules.

\textit{Forward reference}: the present $9\pi$ is a \emph{raw use case} of the axiomatic system (expressing the cost reading~9 as the norm reading~$9\pi$), whereas the RLU coupled-access form ($(4 + 1/\pi)$) introduced in Section~\ref{sec:baryogenesis} of this paper belongs to a \emph{coupled use case} of the axiomatic system (adding the RLU mechanism for object identification) (see ``Use-case classification'' immediately after Definition~\ref{lem:cost13}). The two forms are a natural classification of axiomatic-system mechanisms and do not contradict each other.

\emph{(4) The $+$ in $2 + 9\pi$ (arithmetic addition, consistent with the axiomatic-system $+$).} The two terms of the denominator belong to two distinct reading categories: $2$ is the number of reversible axes (\emph{cost reading} integer, Theorem~\ref{thm:cost-signature}), and $9\pi$ is the \emph{norm reading} of the irreversible residual cost (arc length, Definition~\ref{def:residual} + ``geometric measure''). The combination of two distinct readings is expressed by \emph{standard arithmetic addition}~($+$), which is \emph{consistent} with the orthogonal composition~$+$ of the axiomatic-system axioms---because the axiomatic-system $+$ (orthogonal composition, Axiom~\ref{ax:banya}) is the natural expression of standard arithmetic~$+$ at the axiomatic-system level (see Reader's note, common misconception~\#3). Multiplication implies sharing the same axis and is therefore inappropriate for an orthogonal reading combination.

\emph{(5) Division (ratio, standard arithmetic).} $\sin^2\theta_W$ is a dimensionless mixing ratio and must lie in $[0, 1]$. A \emph{ratio} is expressed as division in standard mathematics; this is not ``defined'' by the axioms but is a \emph{commonplace arithmetic operation}---the present axiomatic system supplies the numerator and denominator \emph{quantities} ($7$ and $2 + 9\pi$), and the division itself is standard arithmetic. Addition ($7 + (2+9\pi) \approx 30.3$) and multiplication ($7 \times (2+9\pi) \approx 212$) both fall outside $[0,1]$ and are therefore excluded (arithmetic fact). Because the denominator exceeds the numerator ($2 + 9\pi \approx 30.27 > 7$), the result naturally lies in the $[0,1]$ range. This is exactly the share/ratio form of a standard mixing angle ($\sin^2\theta_W$ = ``one component${}^2$ / total${}^2$'').

\emph{Exclusion of alternatives.}

$(7-2)/(9\pi) \approx 0.177$: $(7-2)$ is not arithmetic subtraction but \emph{set subtraction (= filter)}---a standard set operation that filters out the reversible part~2 from 7 (= 5 irreversible + 2 reversible, Theorem~\ref{thm:cost-signature}), leaving only the irreversible~5. This set-subtraction/filter operation is not defined by the axioms but is a \emph{commonplace set operation} applied to the axiomatic-system quantities (7, 2).

However, this form \emph{discards} the reversible part~(2) from the result---this undermines the ``irreversible/reversible ratio'' meaning of a mixing angle (mixing is a ratio of two domains, not a quantity of one domain alone). The canonical form $7/(2+9\pi)$ of this paper preserves the full~7 (cost reading) in the numerator and includes both the reversible~2 (cost reading) and the irreversible~$9\pi$ (norm reading) in the denominator, thereby \emph{preserving} the complete quantity structure of the axiomatic system. Accordingly, $(7-2)/(9\pi)$ is arithmetically valid but incompatible with the mixing interpretation and is excluded.

$7/(9\pi - 2) \approx 0.260$: subtracting the reversible cost-reading~$2$ from the irreversible norm-reading~$9\pi$ is a \emph{cross-reading filter}; unlike set subtraction within a single reading, it is meaningless (a filter removes a subset from within one set). Removing either~$2$ or~$9\pi$ from the denominator omits part of the descriptive degrees of freedom and is therefore incomplete.

Conclusion: in the formula $7/(2+9\pi)$, the origin of each number (axioms), the uniqueness of the geometric constant~($\pi$), the necessity of the hemisphere (irreversibility), the additivity of arc lengths, the category combination rule ($+$ as orthogonal composition of two readings), and the ratio operation (standard arithmetic ratio)---all of these are forced by the two reading principles of the axiomatic system together with standard arithmetic. All five steps are a forward chain, and none of them involves a ``choice.'' The ``mechanism by which the axioms automatically produce this ratio'' is the combination of the two reading principles with standard arithmetic, and this is of the same strength as the forward chain in which the Wyler volume ratio yields the coupling constant in the derivation of~$\alpha$.

\emph{Relation of the present derivation path to other candidates.} The full axiomatic system~\cite{han2026} records, in addition to the present $7/(2+9\pi)$, several other dimensionless candidate expressions for $\sin^2\theta_W$ (e.g.\ a purely geometric form $(4\pi^2-3)/(16\pi^2)$, an information-theoretic form $1/\log_2 20$, etc.). These candidates also approach the experimental value 0.23122 to within $0.1\%$.

However, the derivation of this paper is \emph{explicitly} based on the cost enumeration of Definition~\ref{lem:cost13} and the residual cost~$9$ of Definition~\ref{def:residual}. The other candidates do not use this enumeration and belong to derivation paths different from the axiomatic chain of this paper (Axiom~\ref{ax:cost} $\to$ Definition~\ref{lem:cost13} $\to$ Definition~\ref{def:residual} $\to$ residual cost $9 \times$ hemisphere $\pi$). This paper presents the form that is forced on this enumeration-based path; evaluation of the other candidates lies outside the scope of this paper.

The derivation path of this paper is tied to a single cost enumeration---that of Definition~\ref{lem:cost13} and Definition~\ref{def:residual}---and within this tie there is no free choice.

\subsection{Significance of Independence}

$\alpha$ and $\sin^2\theta_W$ emerge from \emph{distinct forward chains} within the same axiomatic system:
\begin{itemize}
\item $\alpha$: signature $(5,2) \to \mathrm{SO}_0(5,2) \to D_5$ volume ratio. Each factor is accounted for at a different epistemic level: $9$ is axiom-derived~(I), $\pi^5/(2^4 \cdot 5!)$ Hua volume is axiom-determined~(II), $8\pi^4$ group volume is axiom-determined~(II), $(\cdot)^{1/4}$ geometric mean is axiom-determined~(II) (see Table~\ref{tab:factors}).
\item $\sin^2\theta_W$: cost reading of 7 (= 5 irreversible + 2 reversible, Theorem~\ref{thm:cost-signature}) and the standard arithmetic ratio $7/(2 + 9\pi)$ with the cost + norm reading combination $2 + 9\pi$ in the denominator.
\end{itemize}
The derivation paths are independent. Because $\alpha$ does not enter the formula for $\sin^2\theta_W$, the independence is complete. That the results of two forward chains from a single set of axioms converge on experimental values with relative deviations of $6 \times 10^{-7}$ and $4 \times 10^{-6}$, respectively, is difficult to dismiss as numerological coincidence.

%%%%%%%%%%%%%%%%%%%%%%%%%%%%%%%%%%%%%%%%%%
\section{Response to Robertson/Gilmore Criticisms}
\label{sec:robertson}

We address three criticisms~\cite{robertson1971,gilmore1972} directed at the Wyler formula.

\subsection{(R1) Lack of physical motivation: ``Why should $D_5$ be related to electromagnetism?''~\cite{robertson1971}}

\emph{Response.} $D_5$ is the space of transformations connecting the irreversible domain (where cost is incurred) to the reversible domain (where no cost is incurred). An electromagnetic interaction, in the present axiomatic system, is the process in which the CAS crosses the bracket boundary from OPERATOR to DATA and modifies the state. $D_5$ represents all possible configurations of this process, and $\alpha$ is the fraction that is actually realised. The physical motivation is the interpretation as ``all possible ways of paying cost.''

\subsection{(R2) Arbitrariness of group selection: ``Why $\mathrm{SO}(5,2)$?''~\cite{robertson1971}}

\emph{Response.} By Theorem~\ref{thm:cost-signature}, the signature $(5,2)$ is forced by the axioms. Once the signature is determined, the preservation group can only be $\mathrm{SO}(5,2)$. The group is not arbitrary---the axioms force the signature, and the signature forces the group.

\subsection{(R3) Non-uniqueness of measure selection: ``A different symmetric space yields a different value''~\cite{gilmore1972}}

\emph{Response.} Gilmore's criticism is that ``there are degrees of freedom in choosing the symmetric space.'' In the present axiomatic system the symmetric space is not chosen---it is derived from the axioms. Because the irreversible cost classification (Proof~A of Theorem~\ref{thm:cost-signature}) is completely fixed by the axioms, the signature $(5,2)$ and the resulting $D_5$ are the unique outcome. There are no degrees of freedom to choose ``a different symmetric space.''

Concretely: Gilmore's counterexamples hold when one assumes ``a dimension other than~7'' or ``a signature other than $(5,2)$.'' In the present axiomatic system, however, the dimension~7 is fixed by Proposition~\ref{ax:dof} (descriptive degrees of freedom: 4~domain axes from Axiom~\ref{ax:banya} $+$ 3~CAS stages from Axiom~\ref{ax:cas}), and the signature $(5,2)$ is fixed by Theorem~\ref{thm:cost-signature}. All of Gilmore's degrees of freedom are thereby eliminated.

%%%%%%%%%%%%%%%%%%%%%%%%%%%%%%%%%%%%%%%%%%

%%%%%%%%%%%%%%%%%%%%%%%%%%%%%%%%%%%%%%%%%%
\section{Auxiliary Results}
\label{sec:auxiliary}

\subsection{Independent structural match: Hamming error-correcting code}
\label{sec:predictions}

This paper focuses on the derivations of $\alpha$ and $\sin^2\theta_W$. To show that these two derivations are not single-constant fits, the most direct approach is to confirm that the core numbers of the same axiomatic system also appear in an information-theoretic domain unrelated to either derivation. In this section we present one such independent result from the information-theory literature that uses neither $\alpha$ nor $\sin^2\theta_W$.

By Proposition~\ref{ax:dof} (descriptive degrees of freedom), the total number of internal degrees of freedom of the present axiomatic system is 4 (domain axes) $+$ 3 (CAS stages) $= 7$. This number 7 exactly matches a well-established result in classical information theory---the minimum code length to protect a 4-bit message against single-bit errors~\cite{hamming1950}:
\begin{linenomath}
\begin{equation}
n_{\mathrm{Hamming}} = 7.
\end{equation}
\end{linenomath}
The Hamming code $[7, 4, 3]$ consists of 4 message bits and 3 parity bits. This decomposition is structurally identical to the $7 = 4 + 3$ decomposition of the present axiomatic system (4 domain axes of Axiom~\ref{ax:banya} + 3 CAS stages of Axiom~\ref{ax:cas}). Hamming's (1950) derivation is based on the parity check matrix and minimum distance analysis, and has no conceptual connection to CAS or the present axiomatic system. This match is therefore not a causal borrowing but the fact that two independent derivations converge on the same decomposition $4+3=7$. The numerology hypothesis (coincidence of small integers) cannot explain why the decomposition into 4 and 3 also matches.

\emph{Natural extension to the quantum domain}. The Steane code $[[7,1,3]]$~\cite{steane1996} in quantum information theory is a quantum error-correcting code built on the present Hamming code $[7,4,3]$ via the CSS (Calderbank--Shor--Steane) construction, encoding 1 logical qubit in 7 physical qubits. That is, the $7 = 4+3$ decomposition of the present axiomatic system matches at the \emph{classical Hamming level}, and the Steane code is the \emph{quantization} of this classical structure. The match thus extends naturally from classical information theory (Hamming, 1950) to quantum information (Steane, 1996).

The status of this correspondence: the axioms force $4{+}3{=}7$ (classification~I, Proposition~\ref{ax:dof}), and the Hamming bound---a theorem of classical information theory---uniquely determines that the perfect single-error-correcting code for 4-bit messages is $[7,4,3]$ with 3 parity bits. This is the same pattern by which the Hua volume (classification~II) is determined once $(5,2) \to D_5$ is forced: the axioms fix the structure, and external mathematics completes the identification. Classification~II.

\emph{Classical vs.\ quantum level specification}. This match is confined to the \emph{4 message + 3 parity decomposition} of the classical Hamming code $[7,4,3]$ and the \emph{4 domain + 3 CAS stages decomposition} of the axiomatic system. As noted above, the quantum Steane code $[[7,1,3]]$ is built on this classical Hamming code via the CSS construction, and its \emph{quantum-level} decomposition is 1 logical qubit + 6 stabilizer generators (3 X-type + 3 Z-type), which structurally differs from the axiomatic system's $4+3$. The $7 = 4+3$ match of the present axiomatic system is therefore a structural matching with the \emph{classical information-theory level} of cost enumeration and does not directly verify the quantum-mechanical content of the axiomatic system --- the quantum aspect of the axiomatic system (in particular the Born rule) lies outside the 4-axiom scope of this paper but is derived in the complete system~\cite{han2026} (Axiom~15 $\to$ Axiom~3 $\to$ Axiom~2; see Section~\ref{sec:discussion}).

\emph{Cl(0,7) algebraic deepening --- complete axiomatic system v1.6 integration}. By the \emph{7-axis orthogonal Clifford form proposition} (a structural consequence of Axiom~1) added in the complete axiomatic system~\cite{han2026} v1.6, the 7 degrees of freedom of the present axiomatic system are explicitly identified not as a mere dimension count but as the \emph{7 generators of the Cl(0,7) Clifford algebra}. That is, $\{\gamma_1, \ldots, \gamma_7\}$ decompose into 4 domain generators (time, space, observer, superposition; Axiom~\ref{ax:banya}) and 3 CAS-stage generators (R, C, S; Axiom~\ref{ax:cas}), satisfying $\gamma_k^2 = -1$ and $\gamma_i \gamma_j = -\gamma_j \gamma_i$ ($i \ne j$). The multivector dimension of this algebra is
\begin{linenomath}
\begin{equation}
\dim \mathrm{Cl}(0,7) = \sum_{k=0}^{7} \binom{7}{k} = 1 + 7 + 21 + 35 + 35 + 21 + 7 + 1 = 2^7 = 128,
\label{eq:cl7-dimension}
\end{equation}
\end{linenomath}
which exactly matches the present axiomatic system's \emph{data type 128 = T(16) + 1 = 8 + 128 + 1, the 128 part} (Axiom 9 proposition; the state space dimension of data type 137).

By this mapping, the agreement with the Hamming code $[7,4,3]$ is strengthened from a mere \emph{structural matching} (numerology resistance level: agreement up to the $4+3$ decomposition) to an \emph{algebraic isomorphism}. Specifically: (i) the 7 code bits of the Hamming 7-bit code space $\leftrightarrow$ the 7 generators of Cl(0,7); (ii) the 4 message bits of the Hamming code $\leftrightarrow$ the 4 domain generators of the present axiomatic system; (iii) the 3 parity bits of Hamming $\leftrightarrow$ the 3 CAS-stage generators of the present axiomatic system; (iv) the $2^7 = 128$ codewords of the Hamming code space (information + parity) $\leftrightarrow$ the 128 objects of the Cl(0,7) multivector space dimension --- all in one-to-one correspondence. This isomorphism is a candidate for potential promotion from this paper's classification~III ``structural matching'' to classification~II ``algebraic identity automatically confirmed by the axioms''; the formal proof (e.g., the equivalence of the Hamming code's syndrome decoding with the grade projection of Cl(0,7)) is left to follow-up work. This deepening does not exceed the 4-axiom scope of this paper (the Cl(0,7) proposition is a structural consequence proposition derived from Axiom~1, not a new axiom addition).


\subsection{Independent production candidate: baryon-to-photon ratio ($\eta_B$)}

\label{sec:baryogenesis}

The Standard Model does not naturally produce the cosmic baryon-to-photon ratio $\eta_B \approx 6 \times 10^{-10}$---that is, ``why does roughly 1 baryon survive per $\sim 10^9$ photons?'' The three necessary conditions for matter--antimatter asymmetry (baryon number violation, C/CP violation, departure from thermal equilibrium) are known, but a mechanism that derives the precise numerical value of $\eta_B$ from these conditions has not been found within the Standard Model. In this section we present a production of $\eta_B$ using only the two quantities ($\alpha$ and $\sin^2\theta_W$, both forward chain) produced by the present axiomatic system. The correction form $(4+1/\pi)$ of this result is at axiom-described level (classification~III) ($\alpha$ and $\sin^2\theta_W$ are forward; RLU itself is axiomatically required for CAS, the $+$ in $(4+1/\pi)$ is uniquely determined by Axiom~\ref{ax:banya}'s orthogonal composition, and the formal systematization of the full RLU mechanism is classification~III; the prefactor 2 is directly derived from the meta-level orthogonal category count of the axiomatic system (A1c)).

\emph{Formula}.
\begin{linenomath}
\begin{equation}
\eta_B = \alpha^4 \cdot \sin^2\theta_W \cdot \left[ 1 - 2\!\left(4 + \frac{1}{\pi}\right)\alpha \right] = 6.14 \times 10^{-10}.
\label{eq:eta-B}
\end{equation}
\end{linenomath}

\textit{Origin of each factor in the paper's vocabulary}:
\begin{itemize}
\item $\alpha$: from Section~\ref{sec:derivation} via signature $(5,2) \to D_5 \to$ Wyler formula. All 4 factors produced from axiomatic structure (Table~\ref{tab:factors}).
\item The exponent $4$ in $\alpha^4$: directly derived from Axiom~\ref{ax:cas} (CAS) via \emph{two complementary interpretations}, both pointing to the same 4 of the axiomatic system (detailed justification in the ``Structural meaning of $\alpha^4$'' paragraph of this section).
\item $\sin^2\theta_W$: produced in Section~\ref{sec:prediction} (forward chain, 5 steps).
\item The correction $(4 + 1/\pi)$: the \emph{coupled access} expression between RLU (defined immediately after Axiom~\ref{ax:cas}) and an object (``ball'') registered in a DATA slot. $4$ = number of domain components covered by one access (the 4 orthogonal axes of Axiom~\ref{ax:banya}; if even 1 axis is missing, a different object is identified; a \emph{cost reading} integer); $1/\pi$ = \emph{the identification expression of RLU itself} (see immediately after Axiom~\ref{ax:cas}): the basic unit of the RLU mechanism that distinguishes 1 object over the angular range $\pi$---not the reciprocal of the arc length but an identification ratio inherent in the RLU definition, belonging to a \emph{separate category} that falls under neither cost reading nor norm reading. The combination ($+$) of the two terms joins two \emph{orthogonal categories} (the domain axis count from cost reading and the identification expression defined by the RLU mechanism) via the orthogonal sum $+$ of Axiom~\ref{ax:banya} --- not arithmetic addition (see Reader's note, common misconception \#3).
\item The prefactor $\mathbf{2}$ in $-2 \times (\ldots) \alpha$: the \emph{meta-level orthogonal category count} of the axiomatic system. Two equivalent interpretations, both derived directly from the axioms of the axiomatic system: (i) two-bracket orthogonality of Axiom~\ref{ax:banya}: DATA bracket $\perp$ OPERATOR bracket $\Rightarrow$ bracket count $= 2$. (ii) structure $\perp$ cost orthogonality of Proposition~\ref{ax:dof}: the degrees of freedom partition orthogonally into a structure category (DATA) and a cost category (OPERATOR) $\Rightarrow$ meta category count $= 2$. Both interpretations are different expressions of the same axiomatic-system orthogonality and yield the same result ($2$). Since RLU coupled access must cover both orthogonal aspects within a single access (both the structure component and the cost component, or equivalently both DATA-side information and OPERATOR-side information), the correction contribution accumulates \emph{twofold} --- this is the axiomatic-system-internal forcing of the prefactor $2$.
\end{itemize}

\emph{Structural meaning of $\alpha^4$: two axiomatic-system interpretations via ball value 4 and CAS state space}.

The exponent $4$ in $\alpha^4$ is directly derived from Axiom~\ref{ax:cas} (CAS) via \emph{two complementary interpretations}, both pointing to the same 4 of the axiomatic system.

\emph{Interpretation (i): ball value view (cost side)}. By Definition~\ref{def:residual}, creating 1 ``ball'' (the minimum mass unit of the axiomatic system) requires 4 writes: 1 time write (timestamp) + 3 space writes ($x, y, z$ coordinates). This is the \emph{ball value 4} (see Definition~\ref{def:residual}). That is, 4 is the \emph{cost count for completing 1 ball}.

\emph{Interpretation (ii): CAS state-space view (state side)}. The monotone bit progression of the 3 CAS stages ($\mathrm{R}, \mathrm{C}, \mathrm{S}$) passes through \emph{4 states}: $000 \to 001 \to 011 \to 111$ ($000$ is the CAS rest state, the starting point of the ``rest$\to$R$\to$C$\to$S$\to$DATA commit'' cycle of Definition~\ref{lem:cost13}). The count of possible monotone progression states is exactly \emph{4}. That is, 4 is the \emph{count of the CAS ball-completion progression state space}.

\emph{Relation to Axiom 2}: Axiom~\ref{ax:cas} specifies only the \emph{3 stages} R, C, S, and the $000$ ``rest state'' is an operational definition introduced in Definition~\ref{lem:cost13} as the starting point of this cycle, not an additional stage of Axiom~\ref{ax:cas}. Interpretation (ii) is a \emph{restatement} of interpretation (i)'s ball value view as ``starting point + 3 stages = 4 positions'' (classification~III).

The two interpretations \emph{converge on the number 4} but, by the structure $\perp$ cost orthogonality of the axiomatic system (see Section~\ref{sec:discussion}, ``Constraining the search space'' paragraph), are quantities of \emph{different categories} --- (i) is a cost-category 4 (ball value, Definition~\ref{def:residual}), (ii) is a state-space-side 4 (CAS monotone state count). Both are directly derived from Axiom~\ref{ax:cas}, and the fact that two independent routes \emph{converge on the same numerical value 4} strengthens the axiomatic-system-internal motivation for the exponent $4$ in $\alpha^4$. The convergence of two independent routes to the same numerical value 4 constitutes strong internal evidence; formal unification of the two categories is addressed in the complete system~\cite{han2026} (see the earlier delimitation regarding the multiplicative-coupling assumption of 4 domain axes).

\emph{Baryon and the $111$ state}. CAS Swap is ``the final stage that registers matter in DATA'' (Axiom~\ref{ax:cas}), and upon completion of this stage the CAS state reaches $111$. This $111$ state is a \emph{completed ball} in the axiomatic system, corresponding to a baryon in physics (a particle with mass that persists) --- that is, baryon $\equiv$ $111$ CAS state $\equiv$ object registered by consuming all ball value 4.

\emph{Physical interpretation of $\alpha^4$}. $\alpha$ is the cost ratio of one cycle (see Section~\ref{sec:derivation}, ``ratio interpretation''), and $\alpha^4$ is this ratio to the 4th power $=$ \emph{the combined cost ratio of ball value 4} $=$ \emph{the cost ratio of 1 ball that has completed the 4-state CAS progression}. This is compatible with the \emph{standard scaling of rare processes} of standard physics ($\sim (\text{small parameter})^N$, $N$ = number of process-completion steps) --- baryon survival requires 4-step completion, so its cost ratio appears as $\alpha^4$. Other powers ($\alpha, \alpha^2, \alpha^3, \ldots$) do not match the ball value 4 of the axiomatic system and are therefore not permitted at the axiom level.

\emph{Baryon versus photon}. $\eta_B$ is the \emph{baryon-photon ratio}. A baryon is a DATA-registered object (having reached the $111$ state by consuming ball value 4); a photon is an OPERATOR-traversing object (traversing the axiomatic system without DATA registration). The two objects belong to different categories, so their ratio $\eta_B$ expresses the \emph{DATA-side registration cost ratio} and naturally contains $\alpha^4$ (the 4th power of ball value 4).

\emph{Relation to standard baryogenesis --- CAS mapping of Sakharov's three conditions}. In the complete axiomatic system~\cite{han2026}, Sakharov's~\cite{sakharov1967} three necessary conditions are automatically satisfied by the CAS structure:
\begin{enumerate}
\item \emph{Baryon number non-conservation (B violation)}: when CAS Swap crosses a domain boundary (+), the FSM state is rearranged and baryon-number-violating paths open (Axiom~\ref{ax:cost}).
\item \emph{C/CP violation}: CAS is irreversible --- the order $\mathrm{R} \to \mathrm{C} \to \mathrm{S}$ is forced and the reverse $\mathrm{S} \to \mathrm{C} \to \mathrm{R}$ is undefined (Axiom~\ref{ax:cas}). Because CAS$^{-1}$ does not exist, CP symmetry cannot be restored.
\item \emph{Departure from thermal equilibrium}: the recovery of the residual cost~$9$ by RLU is not instantaneous but proceeds via geometric decay (Section~\ref{sec:axioms}). This delay corresponds to departure from thermal equilibrium.
\end{enumerate}
The three conditions emerge from three different axioms. When CAS exists, Sakharov's conditions are structurally embedded.

\emph{Unique predictions}: (a)~The magnitude of CP violation $J_\text{CKM} \approx 3 \times 10^{-5}$ must be derivable from the phase volume of the irreversible CAS cycle --- in the complete axiomatic system~\cite{han2026}, $J_\text{CKM} = 3.10 \times 10^{-5}$ (experimental value $(3.08 \pm 0.15) \times 10^{-5}$, deviation $0.62\%$) is already produced. (b)~Whether the numerical coincidence $2^3 = 8$ (CAS states) and the baryon octet is structural is recorded as a structural prediction.

\emph{RLU coupled access interpretation}.

We present an axiomatic-system-internal explanation for \emph{why} the correction $(4 + 1/\pi)$ takes this form. $(4 + 1/\pi)$ itself is the formula expression of coupled access between RLU and a domain object --- in the same sense that ``$qE$ is the definition itself of the electric-force coupling.''

\emph{Forcing by the principle of least action}. As stated in Section~\ref{sec:framework}, ``Least action and minimum cost'' subsection, the present axiomatic system takes as its core principle that all change follows the \emph{minimum cost path}. This is the axiomatic-system-vocabulary expression of the principle of least action in physics. This principle \emph{forces} RLU coupling:

\begin{itemize}
\item \emph{Raw access (without coupling)}: reading the 4 domain components of one object requires one $+$ traversal per component = total $4$ $+$ traversals = cost $4$.
\item \emph{Coupled access (via RLU)}: the RLU index bundles the 4 components into a single object. Access to all 4 components via a \emph{single} $+$ traversal on this single object. Cost $\approx 1$ core access $+ \, 1/\pi$ geometric weight $\approx 1.318$.
\item \emph{Cost difference}: $4 \to 1.318$ --- approximately $70\%$ cost savings.
\end{itemize}

The principle of least action ($=$ minimum cost) \emph{forces the selection} of coupled access over raw access. RLU coupling is therefore not a \emph{choice} but a \emph{necessity} in the axiomatic system.

\emph{Computer science analogy and key difference}. This mechanism has the \emph{same cost-saving structure} as \emph{pointer dereference} optimization in computer science, but the \emph{identification mechanisms differ}. In C:
\begin{itemize}
\item Raw method: read each of the 4 components via a separate memory access (4 reads).
\item Dereference method: access the entire struct via a single pointer dereference (\texttt{*ptr}) (1 read + small overhead).
\end{itemize}
The RLU index plays the role of this ``pointer,'' accessing all 4 domain components in a single $+$ traversal. \emph{Key difference}: a CS pointer uses a \emph{logical address}, but RLU does \emph{not use logical addresses}. Instead, RLU is an \emph{addressless angular index} that distinguishes targets by \emph{angular coordinate (theta angle)} in the OPERATOR bracket (see the RLU definition immediately after Axiom~\ref{ax:cas}). In both cases the core cost saving is covering 4 components in 1 access, but RLU uses angular distinction due to the ``addressless'' property of the axiomatic system --- this is why $1/\pi$ appears (1 unit distinction over the angular range $\pi$). Minimum cost in the present axiomatic system, least action in physics, and dereference optimization in computer science --- three domains expressing the \emph{same cost-saving principle}, though the identification mechanism (addressless angular vs.\ addressed integer) differs across systems.

\emph{Decomposition of $(4 + 1/\pi)$}. The cost encoding of this forced coupled access:
\begin{itemize}
\item $4$ = ``number of domain components covered'' (\emph{cost reading}: the integer 4 axes of Axiom~\ref{ax:banya}) --- all 4 axes must be included in one access for object identification.
\item $1/\pi$ = \emph{the identification expression of RLU itself}: the \emph{essential definition} of the mechanism by which RLU distinguishes targets by angular coordinate (theta angle) without logical addresses (see the RLU definition immediately after Axiom~\ref{ax:cas}). That is, a unit ratio of identifying 1 object over the angular range $\pi$, not derived arithmetically as ``the reciprocal of the arc length'' but a quantity that the RLU mechanism itself introduces into the axiomatic system. This quantity belongs to a \emph{separate category} (the RLU mechanism definition category) falling under neither cost reading nor norm reading, and together with the two reading principles constitutes the quantity-production toolkit of the axiomatic system.
\item $+$ = combines two \emph{different categories} ($4$ is the domain axis count from cost reading; $1/\pi$ is the identification expression defined by the RLU mechanism) within a single coupled access. Combined via the \emph{orthogonal sum} $+$ of Axiom~\ref{ax:banya}, not arithmetic addition --- the two terms are not two quantities of the same type but quantities of two \emph{orthogonal categories}, so the ``type error'' concern from the quantum information community does not arise (arithmetic addition operates within the same type; the $+$ of the present axiomatic system operates between orthogonal categories --- see Reader's note, common misconception \#3).
\end{itemize}

\emph{Difference from $9\pi$ in $\sin^2\theta_W$ --- different use cases of the same axiomatic system}. The $(4 + 1/\pi)$ of this section and the $9\pi$ of $\sin^2\theta_W = 7/(2 + 9\pi)$ in Section~\ref{sec:prediction} may \emph{appear} to have different forms, but both express \emph{different use cases} of the axiomatic system:

\begin{itemize}
\item $9\pi$ in $\sin^2\theta_W$ = \emph{sum of raw $+$ traversals}. Simple accumulation of the arc length $\pi$ for each of the $9$ irreversible $+$ traversals ($9 \times \pi$). Raw cost \emph{before} coupling. The mixing angle is an \emph{internal structural ratio} of the axiomatic system, so the raw form is natural.
\item $(4 + 1/\pi)$ in $\eta_B$ = \emph{encoding of coupled 1-access}. Access to 4 domain components via a single $+$ traversal through RLU. Minimum cost \emph{after} coupling. Since baryogenesis centers on \emph{object identification} (surviving baryons), the RLU coupled form is natural.
\end{itemize}

The difference in form is therefore the \emph{natural classification of mechanisms} in the axiomatic system: \emph{internal structural ratios} take the raw $n\pi$ form, while \emph{object-identification-related quantities} take the coupled $(m + 1/\pi)$ form.

\emph{Numerical results and comparison}:
\begin{linenomath}
\begin{align}
\alpha^4 &= (1/137.036)^4 = 2.834 \times 10^{-9}, \nonumber \\
\alpha^4 \cdot \sin^2\theta_W &= 6.55 \times 10^{-10}, \nonumber \\
1 - 2(4+1/\pi)\alpha &= 1 - 0.0630 = 0.937, \nonumber \\
\eta_B^{\text{derived}} &= 6.55 \times 10^{-10} \times 0.937 = 6.14 \times 10^{-10}.
\end{align}
\end{linenomath}

Planck 2018 measurement~\cite{pdg2024}: $\eta_B^{\text{exp}} = (6.12 \pm 0.04) \times 10^{-10}$ (converted from $\Omega_b h^2$). The difference between the present production and the measured value is $0.5\sigma$.

\textit{Status.} Each component of $(4+1/\pi)$ is forced by the axiomatic structure: $4$ is the domain axis count (Axiom~\ref{ax:banya}), $1/\pi$ is the RLU identification unit (defined after Axiom~\ref{ax:cas}), and $+$ is the orthogonal composition (Axiom~\ref{ax:banya}). No other combination operation exists in the axiomatic system, so this form is unique within the axiomatic structure. The least-cost principle forces coupled access (Section~\ref{sec:framework}), and the numerical result agrees with the measured value at $0.5\sigma$. Prefactor 2 is directly derived from the two-bracket orthogonality of Axiom~\ref{ax:banya} (classification~II). Specifically:

\begin{enumerate}
\item \emph{Forward part}: the product $\alpha^4 \cdot \sin^2\theta_W$ itself is the direct product of two quantities derived in the paper, and the exponent $4 = $ number of domain axes is natural (forward chain).

\item \emph{RLU mechanism extension part}: the $1/\pi$ in the correction $(4 + 1/\pi)$ is the angular discrimination density of RLU, unlike $9\pi$ in $\sin^2\theta_W$ (norm reading), and requires an \emph{extension} of the two reading principles (the reciprocal of norm reading does not fit directly into the standard two readings). This extension belongs to the axiom-described tier (classification~III); RLU itself is axiomatically required, and the $+$ combination is uniquely determined by Axiom~\ref{ax:banya}'s orthogonal composition.

\item \emph{Axiomatic-system-internal interpretation of prefactor 2}: the prefactor $2$ is motivated by the \emph{meta-level orthogonal category count} of the axiomatic system. Two equivalent interpretations: (i) two-bracket orthogonality of Axiom~\ref{ax:banya} (DATA $\perp$ OPERATOR $\Rightarrow$ bracket count $= 2$), or (ii) structure $\perp$ cost orthogonality of Proposition~\ref{ax:dof} (meta category count $= 2$). The interpretation is that since RLU coupled access must cover both orthogonal aspects, the correction contribution accumulates \emph{twofold}. The multiplicative coupling ``two aspects covered $\Rightarrow$ twofold accumulation'' follows the same ``orthogonality$\to$independence$\to$multiplication'' pattern as $(\cdot)^{1/4}$: independent dimensions contribute multiplicatively, which is a mathematical property of orthogonal spaces (classification~II, Table~\ref{tab:axiom-results}).

\item \emph{Overall status}: the $\eta_B$ production is ``$\alpha^4 \times \sin^2\theta_W$ part is forward chain ($\alpha$ and $\sin^2\theta_W$ are both forward), the combination form $(4 + 1/\pi)$ has each component forced by the axiomatic structure and the combination is unique within the axiomatic system's orthogonal composition.''

\item \emph{Remaining scope}: the complete classification of external candidate axiom sets (rule selection problem) is an open problem common to digital physics in general (classification~IV). Within the axiomatic structure of this paper, each component and combination form of $(4+1/\pi)$ is forced, and $\alpha$, $\sin^2\theta_W$, and $\eta_B$ are all products of the axiomatic structure.
\end{enumerate}

\emph{Forward-testability}. $\eta_B$ measurements will be refined by future CMB precision experiments (CMB-S4, LiteBIRD, and other next-generation experiments) and precision big-bang nucleosynthesis measurements. The present $0.5\sigma$ agreement is verifiable/falsifiable by future precision improvements. $\eta_B$ is one of the forward-testable items of this paper (4 additional predictions are listed in Table~\ref{tab:predictions}), and despite its axiom-described status (classification~III), considering that $\eta_B$ is a cosmological initial condition that the Standard Model does not naturally produce, the fact that the present axiomatic system reaches within $0.5\sigma$ of the measured value using only the two quantities derived in the present paper, ($\alpha$, $\sin^2\theta_W$) demands subsequent forward verifications.

\emph{Meaning and limitations of forward-testability (falsifiability)}.
The results of this paper fall into three categories.

\begin{center}
\small
\begin{tabular}{p{2.2cm}p{3.5cm}p{2.5cm}p{3.2cm}}
\toprule
\textbf{Category} & \textbf{Result} & \textbf{Type} & \textbf{Verification} \\
\midrule
Forward chain derivation & $\alpha = 1/137.036$, $\sin^2\theta_W = 0.23122$ & post-dictive matching & Agreement with measured values ($6 \times 10^{-7}$, $4 \times 10^{-6}$) \\
Structural match & Hamming $[7,4,3]$: $7 = 4+3$ decomposition match & post-dictive matching & Structural, not quantitative \\
Axiom-described (III) & $\eta_B = 6.14 \times 10^{-10}$ & forward-testable & Verifiable/falsifiable by future CMB. Currently $0.5\sigma$ \\
\bottomrule
\end{tabular}
\end{center}

Under Popper's falsifiability criterion, (i) and (ii) are ``the axiomatic system reproduces known values'' (forward derivation or mathematically determined correspondence), and (iii) is ``the axiomatic system is verifiable by future measurements.'' The grounds for the claim that this paper ``differs from curve fitting'' are (a)~the search-space restriction argument of Section~\ref{sec:discussion}, (b)~the axiom-internal forcing of the two forward-chain derivations ($\alpha$, $\sin^2\theta_W$), and (c)~the future verifiability of $\eta_B$. The complete axiomatic system~\cite{han2026} records 130 unique predictions, and additional forward items will be presented in follow-up papers.

%%%%%%%%%%%%%%%%%%%%%%%%%%%%%%%%%%%%%%%%%%
\section{Discussion}
\label{sec:discussion}

\subsection{The energy-scale problem}

The present derivation corresponds to $\alpha$ in the low-energy limit ($q^2 \to 0$). The reason is as follows: the axiomatic system computes the ``base cost'' of a CAS operation, and this base cost corresponds to the measured value in the low-energy limit. The experimentally measured $\alpha = 1/137.036$ in QED is a renormalized coupling constant that runs with energy, yielding $\alpha(M_Z) \approx 1/128$. That the present derivation reproduces this low-energy limiting value follows from the Wyler volume ratio reflecting the full geometric structure of $D_5$. Since Axiom~3 of the complete axiomatic system~\cite{han2026} declares the DATA (discrete) / OPERATOR (continuous) distinction, the system/domain time separation is an immediate consequence of Axiom~\ref{ax:banya} (classification~I; see Section~\ref{sec:axioms}), and Axiom~3 of the complete system~\cite{han2026} further specifies the discrete-to-continuous transition. 
\textit{Note on scheme dependence of $\sin^2\theta_W$.} The value $\sin^2\theta_W = 7/(2+9\pi) = 0.23122$ derived from the present axiomatic system agrees with the experimental value in the $\overline{\mathrm{MS}}$ @ $M_Z$ scheme ($\approx 0.23122$) to a relative deviation of approximately $4 \times 10^{-6}$. By contrast, the on-shell scheme value is approximately $0.2234$, differing by about 3.5\%.

The $\overline{\mathrm{MS}}$ scheme provides a ``threshold-free'' value from which physical particle-mass thresholds have been removed, and the structural ratios of the present axiomatic system operate at a level prior to the definition of energy scales. Since Axiom~3 of the complete axiomatic system~\cite{han2026} declares the DATA (discrete) / OPERATOR (continuous) distinction, the correspondence with the threshold-free scheme is consistent with the axiomatic structure. The system/domain time separation is an immediate consequence of Axiom~\ref{ax:banya} (Section~\ref{sec:axioms}); specific correspondences in physical equations are obtained by substituting known physical quantities into the Banya equation.

\subsubsection{CAS cost mechanism for $\alpha$ running}
\label{sec:running}

In the complete axiomatic system~\cite{han2026}, the energy dependence (running) of $\alpha$ is described as the accumulation of CAS Compare-false events. Each time CAS returns false at the Compare stage, a virtual pair is generated (one loop), and the accumulation count $N$ corresponds to the energy scale:
\begin{linenomath}
\begin{equation}
\alpha(N) = \frac{\alpha}{1 - \frac{\alpha N}{3\pi}}.
\label{eq:running}
\end{equation}
\end{linenomath}
In the denominator, $3 =$ CAS 3 stages (R, C, S) and $\pi =$ one-cycle phase of the d-ring. The standard QED 1-loop $\beta$-function coefficient is $2/(3\pi)$, where the denominator $3$ is the CAS step count. In Eq.~\eqref{eq:running}, $N$ is the effective accumulation count including charge degrees of freedom, corresponding to $2\sum_f Q_f^2 \ln(\mu/m_f)$. The 2-loop coefficient $\beta_1 = -1/4$ exactly matches $-1/4 =$ the reciprocal of the number of domain axes on which Swap operates~\cite{han2026}.

For QCD, the 1-loop $\beta_0 = 7/(4\pi)$ has numerator $7 =$ the CAS complete descriptive degrees of freedom (Proposition~\ref{ax:dof}), and $11N_c - 2n_f = 21 = 3 \times 7$ is CAS 3 stages $\times$ 7 degrees of freedom~\cite{han2026}.

\emph{Unique prediction}: if $\beta$-function coefficients at $n$-loop are functions of CAS structural numbers, then $\beta_2$, $\beta_3$, and higher-order coefficients must systematically contain CAS structural numbers (3, 4, 7, 9, 13). Confirmation of this pattern would establish the CAS origin of running couplings; failure would reveal the limits of this correspondence.

\subsubsection{Structural interpretation of $137 = 128 + 9$}

The internal degrees of freedom of the present axiomatic system can be represented with 7 bits (Proposition~\ref{ax:dof}: 4 domain axes + 3 CAS stages $= 7$; Axiom~\ref{ax:delta}: $\delta$ is a global flag outside this 7-bit \emph{finite state machine (FSM)}). The axiomatic system possesses the following three integer invariants internal to the system: (i) $128 = 2^7$ (number of possible states of the 7-bit FSM); (ii) $9 = 13 - 4$ (residual cost arithmetic, Definition~\ref{def:residual}); (iii) the integer part $137$ of $1/\alpha$ from the Wyler derivation (Section~\ref{sec:derivation}). These three integers satisfy the arithmetic relation $137 = 128 + 9$. By the \emph{Cl(0,7) Clifford form proposition} of the complete axiomatic system v1.6~\cite{han2026}, the $128 = 2^7$ in (i) is algebraically identified, beyond a 7-bit state count, with the \emph{multivector space dimension of Cl(0,7)} (Eq.~\eqref{eq:cl7-dimension}) --- that is, the algebraic essence of data type 128 is explicitly the multivector algebra over the 7 generators (see ``Cl(0,7) algebraic deepening'' in Section~\ref{sec:predictions}).

Standard QED RG running~\cite{pdg2024} gives $1/\alpha(0) \approx 137.036$ (Thomson limit), $1/\alpha(M_Z) \approx 127.918$ ($M_Z$ scale), and $\Delta\alpha^{-1} \approx 9.118$ (vacuum polarization), obeying $1/\alpha(0) = 1/\alpha(M_Z) + \Delta\alpha^{-1}$. At the level of integer parts, $128 + 9 = 137$, coinciding with the arithmetic of the axiomatic system.

\emph{Interpretation via CAS running}. In the CAS running mechanism of Section~\ref{sec:running}, the accumulation count $N$ of Compare-false events corresponds to $\Delta\alpha^{-1}$. When the residual cost~$9$ (Definition~\ref{def:residual}) is exhausted by $N = 9$ Compare-false events, $\Delta\alpha^{-1} = 9$, and $1/\alpha(0) - 9 = 128 = 2^7$ is the FSM state count. That is, $128$ is the \emph{ground-state count} of the FSM after all residual cost has been consumed, and $9$ is the \emph{total cost consumed}. The relation $137 = 128 + 9$ reads as ``ground states + residual cost = low-energy limit.''

\emph{Unique prediction}: the nearest integer to $1/\alpha(M_Z)$ being exactly $128 = 2^7$ is a consequence of the 7-bit FSM structure. If future precision measurements of $\alpha(M_Z)$ shift the integer part away from $128$, this interpretation is falsified.

\subsection{Constraining the search space: blocking numerology}

The key distinction from numerology is ``how many candidates were searched before selecting the one that fits.'' In the present system, the search space is not open---it is closed by the axioms.

The descriptive degrees of freedom of Proposition~\ref{ax:dof} fall into two categories, and the numbers admissible in each category are constrained by the following filtering rules:
\begin{enumerate}
\item \emph{Structure registry} (DATA category): numbers derived directly from the axiomatic structure (number of axes, stages, brackets). Admissible structure numbers: $\{1, 2, 3, 4, 7\}$ ($1$ bit basis, $2$ brackets, $3$ CAS stages, $4$ domain axes, $7$ internal degrees of freedom).
\item \emph{Cost registry} (OPERATOR category): numbers derived from Definition~\ref{lem:cost13} and Definition~\ref{def:residual}. Admissible cost numbers: $\{1, 4, 5, 9, 13\}$ ($1$ unit traversal cost, $4$ ball value, $5$ irreversible axis count or write traversal count, $9$ residual cost, $13$ total cost).
\item \emph{Category separation}: even when two numbers are numerically equal, they represent different quantities if they belong to different categories (e.g., the structural $9 = 7+2$ and the cost $9 = 13-4$ are distinct quantities; the $\sin^2\theta_W$ formula of this paper uses the latter, i.e.\ the cost $9$).
\item \textbf{Independence}: numbers that factorize into preceding numbers are excluded (e.g., $6 = 2 \times 3$ and hence is not an independent DOF).
\item \emph{Continuous-quantity cutoff}: infinitesimals of normed spaces are not discrete and therefore are not registered as degrees of freedom.
\end{enumerate}

Under these rules the total number of admissible numbers in this paper is only $10$ (structure $5$ + cost $5$). Formulas are constructed not from arbitrary combinations of $\pi$, $e$, and integers but only from a finite set derived from the axioms. Because the search space is closed, the degrees of freedom available for post-hoc fitting are severely limited.

\emph{Posterior probability estimate.} Even though the number set is closed, freedom remains in the choice of operations (four arithmetic operations, exponentiation, inclusion or exclusion of $\pi$), so the size of the formula space must be estimated. From the $10$ structure/cost constants, with 4 binary operations $\times$ nesting depth 2--3 $\times$ $\pi$ included/excluded $= O(10^3)$ candidate formulas are possible. The probability that one of these matches a particular physical constant to within a relative deviation of $10^{-6}$ is $\sim 10^{-3}$.

The essential point: the $\alpha$ derivation is a forward chain starting from Theorem~\ref{thm:cost-signature} (axioms $\to$ signature $\to$ $D_5$ $\to$ Wyler volume ratio), not ``a chance hit from the $O(10^3)$ formula space'' but rather ``a group-theoretically forced consequence.'' The chance hypothesis of a single $\alpha$ hit probability $\sim 10^{-3}$ is therefore rejected by the very existence of the forward chain.

The $\sin^2\theta_W$ result likewise follows from a forward chain (all 5 steps forced by the two reading principles of the axiomatic system; see Section~\ref{sec:prediction}), with the \emph{same strength} of forward-chain status as the $\alpha$ derivation. Neither result is a chance hit from the $O(10^3)$ candidate formula space; both follow from the forward derivation of the axiomatic system. The chance hypothesis that two independent forward derivations simultaneously agree with experimental values is rejected even more strongly. This estimate is an order-of-magnitude upper bound.

\emph{Status of axiom choice and response to the circularity critique.} That Theorem~\ref{thm:cost-signature} follows ``automatically'' from the axioms is a structural strength of the axiomatic system. In every axiomatic system, theorems are logical consequences of the axioms---calling this circular would equally apply to the derivation of ``the sum of interior angles of a triangle $= 180^{\circ}$'' from Euclid's parallel postulate. The critique that the answer is ``built into'' the axioms is a meta-critique applicable to axiomatic systems in general, does not apply to this paper.

The contribution of this paper is not the bare fact that the $(5,2)$ partition emerges from the axioms, but rather \emph{why these axioms are physically justified}. The justification is threefold: (a)~the 3-stage CAS is the minimum-cost atomic operation that preserves causality: the lower bound (2 steps or fewer violates causality) and the upper bound (4 steps or more violates least cost) force exactly 3 steps; LL/SC shares the same 3-stage structure and yields the same $(5,2)$; TAS (2 stages) fails due to the absence of Compare (Section~\ref{sec:framework}). Within the axiomatic structure, the 3-step form is unique. (b)~The DATA/OPERATOR separation is the distinction between a ``discrete record space'' and a ``continuous operational domain''; the moment the axiomatic system distinguishes discrete (DATA) from continuous (OPERATOR), this separation is inevitable. (c)~The resulting $\alpha$ and $\sin^2\theta_W$ agree with experimental values to relative deviations of $6 \times 10^{-7}$ and $4 \times 10^{-6}$, respectively.

\emph{Logical minimal convergence.} Each component of the present axiomatic system is a candidate for the logical minimal convergence of the minimum-cost principle. To change a state while preserving causality, the minimum is three stages---``read$\to$compare$\to$write'' ($2$ stages would be writing without comparing, violating causality). For a 3-stage operator to function, the operand and the operational domain must be separated, making the DATA/OPERATOR distinction inevitable. The operand (DATA) requires at minimum 2 axes (time, space), and the operational domain (OPERATOR) likewise requires at minimum 2 axes (observer, superposition)---without an observer there is nowhere to receive the branch of Compare, and without superposition there is no path for CAS to reference DATA. By atomicity and orthogonality, the cost can only take values $\{0, +1\}$. Combining all these minimal conditions yields $4 + 3 = 7$ degrees of freedom and the irreversible/reversible partition $(5,2)$.

\emph{Note on the status of the convergence argument.} Each step of this convergence---minimum of 2 axes, binary atomicity, DATA/OPERATOR separation---is forced by the axiomatic structure, and their combination converges to $(5,2)$ as proved by Theorem~\ref{thm:cost-signature}. see the rule selection discussion in Section~\ref{sec:limitations}; within the axiomatic structure the convergence is unique. The axioms were not reverse-engineered to match $(5,2)$; rather, following the minimum-cost principle led to $(5,2)$ as the unique convergence point. The a posteriori verification of this convergence is the experimental agreement of $\alpha$ and $\sin^2\theta_W$.

Robertson's critique ``why $\mathrm{SO}(5,2)$?'' is elevated in this paper to ``why these axioms?'' This is a meta-question shared by every axiomatic system---Euclid's five postulates, the Zermelo--Fraenkel axioms, the von Neumann axioms of quantum mechanics. The necessity of axioms cannot be proved---but when the consequences of axioms agree with experiment, the axiomatic system is justified a posteriori.

\subsection{Transition from discrete axioms to continuous physical quantities}

The present axiomatic system is discrete (cost $+1$, in bit units), yet the final results---$\alpha$, $\sin^2\theta_W$, etc.---are continuous real values. Where does this transition occur?

The transition point is step~2$\to$step~3 of Theorem~\ref{thm:cost-signature}. Once the cost classification (discrete: 5 vs.\ 2) determines the metric signature $(5,2)$, the group $\mathrm{SO}(5,2)$ (or, when defining the quotient space $D_5$, the identity component $\mathrm{SO}_0(5,2)$) that preserves that signature is a continuous Lie group. That is, \emph{the discrete classification selects a continuous group}. All subsequent calculations---the quotient space $D_5$ and the volume ratio---are continuous mathematics.

This is analogous to how, in lattice gauge theory, a discrete lattice approaches the continuum limit.

The viewpoint that physical laws emerge from discrete structure is shared by the computational-universe programs of Zuse~\cite{zuse1969}, Fredkin~\cite{fredkin2003}, and Wolfram~\cite{wolfram2002}; Wheeler's ``it from bit''~\cite{wheeler1990}; Lloyd's computational-capacity estimate~\cite{lloyd2002}; and Tegmark's mathematical-universe hypothesis~\cite{tegmark2008}. In particular, 't~Hooft's cellular automaton interpretation of QM~\cite{thooft2016} is the most directly comparable program, seeking to derive quantum mechanics from a discrete deterministic model. The difference is that 't~Hooft aims to reproduce quantum mechanics, whereas the present system aims to derive a metric signature from an irreversible cost structure---both the starting point and the destination differ. For the thermodynamic cost of irreversible operations, see Landauer~\cite{landauer1961}. Our treatment of the fundamental status of dimensionless constants follows Duff et al.~\cite{duff2002}.

The difference from lattice gauge theory is: there, the lattice spacing $a \to 0$ limit is required, whereas here it is not a limit but a \emph{classification} (5+2) that determines the group. The discrete structure fixes the integer parameters (dimension, signature) of the continuous group, and the continuous fine structure (volume ratio) is produced automatically by group theory.

Since Axiom~3 of the complete axiomatic system~\cite{han2026} declares the DATA (discrete) / OPERATOR (continuous) distinction, this transition is part of the axiomatic structure (classification~I). The discrete-to-continuous transition is already within the axioms: CAS manages costs on the OPERATOR sphere (continuous) via RLU angular coordinates, and CAS Swap records results in DATA (discrete). Both directions---discrete$\to$continuous (RLU sphere management) and continuous$\to$discrete (Swap recording)---are contained in Axioms~\ref{ax:banya}, \ref{ax:cas}, and \ref{ax:cost}.

\subsection{Axiomatic origin of structural constants}

Table~\ref{tab:constants} makes explicit the axiomatic origin of the structural constants appearing in the derivations of $\alpha$ and $\sin^2\theta_W$. All numbers originate from the structure of Axioms~\ref{ax:banya}--\ref{ax:delta}:

\begin{table}[H]
\begin{adjustwidth}{-\extralength}{0cm}
\caption{Structural constants appearing in the derivations of $\alpha$ and $\sin^2\theta_W$ and their axiomatic origins.\label{tab:constants}}
\small
\begin{tabularx}{\fulllength}{p{2.5cm}p{3.5cm}X}
\toprule
\textbf{Number (category)} & \textbf{Role} & \textbf{Remarks} \\
\midrule
$2$ (structure) & bracket count, $\sin^2\theta_W$ denominator & DATA/OPERATOR two brackets (Axiom~\ref{ax:banya}) \\
$3$ (structure) & CAS stage count & R, C, S (Axiom~\ref{ax:cas}) \\
$4$ (structure/cost) & $2^4$ domain axes; ball value & Structure: domain axis count (Axiom~\ref{ax:banya}). Cost: 1 time + 3 space writes (Definition~\ref{def:residual}). The two 4s are quantities of different categories \\
$5$ (cost) & irreversible axis count & $= 7-2$ (Theorem~\ref{thm:cost-signature}) \\
$7$ (structure/cost) & $\sin^2\theta_W$ numerator & Structure: $4+3$ (Proposition~\ref{ax:dof}). Cost: $5+2$ (Theorem~\ref{thm:cost-signature}) \\
$9$ (cost) & $\sin^2\theta_W$ denominator & Residual cost $13-4$ (Definition~\ref{def:residual}) \\
$13$ (cost) & cost enumeration & $8$ reads + $5$ writes (Definition~\ref{lem:cost13}) \\
$\pi$ & Wyler formula, $\sin^2\theta_W$ & hemispheric arc length of irreversible traversal (Axiom~\ref{ax:cost} + Axiom~\ref{ax:cas}) \\
\bottomrule
\end{tabularx}
\end{adjustwidth}
\end{table}

\subsection{Distinction from Eddington}

Several attempts exist to derive $\alpha$ from first principles. Table~\ref{tab:prior-art} compares the major prior work with the present paper. The most precise measurement of $\alpha$ was performed by Morel et al.~\cite{morel2020} with an uncertainty of $81$\,ppt.

\begin{table}[H]
\caption{Prior attempts at deriving $\alpha$ from first principles.\label{tab:prior-art}}
\begin{adjustwidth}{-\extralength}{0cm}
\small
\begin{tabularx}{\fulllength}{p{2.5cm}p{3.5cm}p{3.0cm}X}
\toprule
\textbf{Approach} & \textbf{Method} & \textbf{Result} & \textbf{Failure / distinction} \\
\midrule
Eddington~\cite{eddington1946} & 16-fold algebraic combinatorics & $\alpha^{-1} = 136 \to 137$ & Reverse-engineering suspicion; no independent check; no axiom necessity \\
Koide~\cite{koide1983} & Lepton mass ratio pattern & Numerical coincidence & Origin not provided \\
Connes~\cite{connes1996} & Noncommutative geometry & Constant derivation attempt & $\alpha$ not numerically reproduced \\
Lisi~\cite{lisi2007} & $E_8$ unification & All-particle embedding & Refuted by Distler--Garibaldi~\cite{distler2010} \\
Bars~\cite{bars2006} & 2T-physics $(d,2)$ & Same $\mathrm{SO}(5,2)$ & $(5,2)$ is spacetime signature; this paper: cost category. $\alpha$ not derived \\
\midrule
\textbf{This paper} & CAS cost structure & $1/\alpha = 137.036\,082$ & (i) Signature derived first, (ii) $\sin^2\theta_W$ independent check, (iii) search-space blocked \\
\bottomrule
\end{tabularx}
\end{adjustwidth}
\end{table}

\subsection{Scope and limitations}
\label{sec:limitations}

This paper states its limitations explicitly throughout the text. Table~\ref{tab:limitations} collects and classifies all of them.

\begin{table}[H]
\begin{adjustwidth}{-\extralength}{0cm}
\caption{Scope and limitations --- classified summary.\label{tab:limitations}}
\small
\begin{tabularx}{\fulllength}{cp{4.2cm}X}
\toprule
\textbf{Class} & \textbf{Item} & \textbf{Remarks} \\
\midrule
\multicolumn{3}{l}{\emph{III. Axiom-described --- derivation path complete in text; formal systematization reserved for dedicated papers}} \\
\midrule
III & RLU mechanism formal systematization & Full mechanism described. $+$ uniqueness in $(4+1/\pi)$ established by Axiom~\ref{ax:banya}. RLU formal systematization in~\cite{han2026} (Section~\ref{sec:baryogenesis}) \\
III & $\sin^2\theta_W$ alternative candidates & $1/\log_2 20$ etc.\ recorded. $7/(2+9\pi)$ uniqueness established by 5-step argument (Section~\ref{sec:prediction}). Only exclusion of other candidates reserved for dedicated paper \\
III & $\alpha$ precision $\Delta = 0.000\,083$ & Possible correction at order $\alpha^3 \approx 4 \times 10^{-7}$ (Section~\ref{sec:limitations}) \\
III & CAS 3-step uniqueness & Lower bound (causality) + upper bound (least cost) force 3 steps (Section~\ref{sec:framework}). Only complete external classification reserved for dedicated paper (Section~\ref{sec:framework}) \\
III & $\sin^2\theta_W$/$\eta_B$ form difference & Raw use case ($9\pi$) vs.\ coupled use case ($(4+1/\pi)$). No contradiction (Section~\ref{sec:prediction}, Section~\ref{sec:baryogenesis}) \\
III & SU(3)$\times$SU(2)$\times$U(1) correspondence & R,C,S $\to$ U(1),SU(2),SU(3) mapping complete. 12 gauge bosons derived. Coupling-constant ratio quantification in~\cite{han2026} (Section~\ref{sec:discussion}) \\
\midrule
\multicolumn{3}{l}{\emph{IV. Out of scope --- not included in the scope of this paper}} \\
\midrule
IV & Rule selection problem & Open meta-problem of digital physics in general (Section~\ref{sec:limitations}) \\
\bottomrule
\end{tabularx}
\parbox{\fulllength}{Classification criteria: see Table~\ref{tab:axiom-results}.}
\end{adjustwidth}
\end{table}

This paper adopts the format of ``explicit delimitation + verifiable partial results.'' The detailed analysis of the last item in the table, ``$\alpha$ measurement precision $\Delta$,'' is given in the immediately following paragraph.

\emph{Detailed analysis of $\alpha$ precision $\Delta$.} The difference $\Delta = 0.000\,083$ between $1/\alpha = 137.036\,082$ produced by the Wyler formula and the experimental value $137.035\,999\,177$ arises at the fourth decimal place. Two possible origins of this discrepancy are:
\begin{enumerate}
\item \emph{Discrete--continuous gap}: the present axiomatic system is discrete (cost $+1$), yet the Wyler volume ratio is continuous geometry. The transition from discrete structure to continuous volume ratio naturally generates higher-order corrections.
\item \emph{Higher-order cost correction}: a correction of order $\alpha^3 \approx 4 \times 10^{-7}$ is expected within the system, consistent with the magnitude of $\Delta$ ($6 \times 10^{-7}$).
\end{enumerate}
The magnitude of $\Delta$ ($6 \times 10^{-7}$) is of the same order as $\alpha^3 \approx 4 \times 10^{-7}$. The present axiomatic system is discrete while experimental measurements are continuous real numbers; a residual from this discrete--continuous gap is a structural characteristic of a discrete axiomatic system. That the residual is of order $\alpha^3$ is consistent with the existence of a higher-order cost correction.

\emph{The rule selection problem.} The present axiomatic system confronts a fundamental problem common to all axiomatic approaches in digital physics---Zuse's \textit{Rechnender Raum}~\cite{zuse1969}, 't Hooft's cellular automaton interpretation~\cite{thooft2016}, Wolfram's hypergraph rewriting~\cite{wolfram2002}: \emph{why these axioms? Why not others?}

The answer offered in this paper---minimum cost + causality preservation---is the design principle of the axiomatic system, verified by its results. This paper provides two partial answers: (a) the search space is restricted to 10 numbers (5 structural + 5 cost), closing the degrees of freedom for axiom selection (Section~\ref{sec:discussion}), and (b) the convergence of $\alpha$ and $\sin^2\theta_W$ to experimental values at relative deviations of $6 \times 10^{-7}$ and $4 \times 10^{-6}$ constitutes empirical verification of this axiom set. Complete classification of external candidate axiom sets is an open problem common to digital physics in general and lies outside the scope of this paper.

\subsection{Status of the present axiomatic system: relation to existing physics}

We state anticipated critiques of the present axiomatic system and our responses. We first summarize the positioning relative to existing axiomatic reconstruction programs in Table~\ref{tab:comparison}. The distinctive feature of the present axiomatic system is that it \emph{produces numerical values from axioms}. Other programs reconstruct the structure of quantum mechanics but do not produce physical constants.

\begin{table}[H]
\begin{adjustwidth}{-\extralength}{0cm}
\caption{Comparison with axiomatic reconstruction programs --- positioning of the present axiomatic system.\label{tab:comparison}}
\small
\begin{tabularx}{\fulllength}{p{3.2cm}p{2.3cm}p{2.3cm}p{2.3cm}p{2.5cm}X}
\toprule
 & \emph{Hardy 2001} & \emph{CDP 2011} & \emph{Spekkens 2007} & \emph{Masanes--M\"uller 2011} & \emph{This work} \\
\midrule
Axiom count & 5 & 6 & 2 & 5 & \emph{4+1} \\
Core principle & continuous reversibility & purification & knowledge balance & physical requirements & \emph{minimum cost = least action} \\
Reconstruction target & full QM & full QM & partial QM phenomenology & QM information bounds & \emph{coupling constants $\alpha$, $\sin^2\theta_W$} \\
Numerical output & none & none & none & none & \emph{$\alpha$, $\sin^2\theta_W$, $\eta_B$} \\
Experimental comparison & none & none & none & none & \emph{$6 \times 10^{-7}$, $4 \times 10^{-6}$, $0.5\sigma$} \\
Free parameters & 0 & 0 & 0 & 0 & \emph{0} \\
Program position & full reconstruction & full reconstruction & partial & partial & \emph{cost-theoretic branch} \\
\bottomrule
\end{tabularx}
\end{adjustwidth}
\end{table}

\subsubsection{``Where is the Hilbert space?''}

The present axiomatic system does not target the Hilbert-space formalism of quantum mechanics; it targets numerical output. The reason is as follows.

The Banya frame is a \emph{virtual experimental universe} that models nature's least-action principle as minimum-cost operations. The design goal of this virtual universe is not ``to reproduce existing physical formalisms internally'' but rather ``to construct a minimum-cost circuit under axiomatic constraints and to check whether the cost pattern of that circuit agrees with actual physical constants.''

By analogy: a flight simulator does not internally reproduce the metal alloys or the fluid-dynamics equations of an actual aircraft. The goal of the simulator is that the output (flight trajectory) for a given input (control stick) matches that of the actual aircraft. Even if the internal implementation differs from the actual aircraft, the simulation succeeds if the input--output relation agrees.

Likewise: the internal implementation of the present axiomatic system consists of CAS (conditional state transitions) and a cost structure. This is a \emph{different language} from Hilbert spaces, Lagrangians, and gauge groups. However, if the output---the numerical values of physical constants---agrees with experiment, the model succeeds. The Hilbert space is \emph{one language} for describing nature, not \emph{the only language}.

Zurek's quantum Darwinism~\cite{zurek2009} also derives classicality from the information structure of observers and environments, showing that descriptions that do not pass through Hilbert space can be physically valid. More generally, Hardy's ``Quantum Theory From Five Reasonable Axioms''~\cite{hardy2001} and Chiribella, D'Ariano, and Perinotti's ``Informational derivation of quantum theory''~\cite{cdp2011} are axiomatic programs that \emph{derive} the entirety of quantum mechanics from information-theoretic primitives --- the present axiomatic system's ``derivation of $\alpha$ from CAS cost'' belongs to the same scholarly tradition (information-theoretic axiomatization), with the distinction that it focuses on a specific coupling constant, $\alpha$.

\subsubsection{``Where does the gauge group SU(3)$\times$SU(2)$\times$U(1) come from?''}

The qualitative mapping is presented below. The quantitative derivation of coupling-constant ratios is completed in the full system~\cite{han2026} and will be presented in dedicated papers.

The Banya equation $\delta^2 = (\mathrm{time}+\mathrm{space})^2 + (\mathrm{observer}+\mathrm{superposition})^2$ is a domain-transformation structure. Projecting the same axioms onto different domains yields different physical sectors. $\alpha$ is the result of projecting the CAS cost path onto the electromagnetic domain, and $\alpha_s$ is the result of projecting the same cost path onto the strong domain. The gauge group structure is an effective description of this domain transformation---not a target that the axioms must reproduce, but an observation-side description produced by the domain projection of the axioms.

Concretely: in the complete axiomatic system~\cite{han2026}, the 3 CAS stages map to the Standard Model gauge group: Read (internal DOF 1) $\to$ U(1) (generators $1^2 - 1 = 0$; photon has no self-interaction), Compare (internal DOF 2) $\to$ SU(2) (generators $2^2 - 1 = 3$: $W^+, W^-, Z$), Swap (internal DOF 3) $\to$ SU(3) (generators $3^2 - 1 = 8$: 8 gluons). Total gauge bosons $= 0 + 3 + 8 + 1(\text{photon}) = 12 = \binom{4}{2} \times 2$ (choose 2 from 4 domain axes $\times$ 2 directions)~\cite{han2026}. \emph{Unique prediction}: a 4th gauge interaction ($Z'$, etc.) does not exist because CAS has no 4th stage --- discovery of $Z'$ or $W'$ below 10~TeV at LHC Run 3/4 would falsify this correspondence. The quantitative derivation of coupling-constant ratios is a task for subsequent papers.

\subsubsection{``Is this a physical theory or a computational model?''}

Both. More precisely: it is a \emph{computational physical model} constructed under the assumption that the operations of nature obey minimum cost.

Conventional physics describes nature in the mathematical language of continuous manifolds, differential equations, and Hilbert spaces. The present axiomatic system describes the same nature in a \emph{different mathematical language}: discrete operations, cost functions, and state transitions. The two languages need not be mutually translatable---``reproducing'' differential equations as discrete operations is not the goal. The goal is to check whether the two descriptions produce the same numerical output (physical constants).

The reason this difference is unavoidable is as follows. If one designs a circuit faithful to the minimum-cost principle, its internal structure is necessarily discrete and order-dependent---counting cost requires stages, stages require order, and order implies discrete rather than continuous. The language of minimum-cost circuits therefore necessarily differs from the language of continuous physics. This difference is a structural consequence of faithfully modeling the minimum-cost principle.

\subsubsection{``Why read/compare/write---this differs from quantum measurement''}

Correct. It differs. And this difference is the point.

In quantum mechanics, ``measurement'' is an irreversible process accompanied by wavefunction collapse. In the present system, ``read (Read)$\to$compare (Compare)$\to$write (Swap)'' is the indivisible 3-stage CAS operation. These two descriptions, in different languages, possess \emph{structural analogy}; the formal isomorphism is derived in the full axiomatic system~\cite{han2026}: the Born rule $|\psi|^2$ is produced via the chain Axiom~15 ($\delta$ 128 states) $\to$ Axiom~3 (FSM) $\to$ Axiom~2 (CAS) as ``the tally of $\delta$'s free selection within the FSM,'' constituting a CAS version of Gleason's theorem~\cite{han2026}. Because Axiom~15 ($\delta$'s 128-state structure) is not among the 4 axioms of the present paper, the formal derivation of the Born rule lies outside the present scope.

Table~\ref{tab:qm-cas} states the differences explicitly.

\begin{table}[H]
\caption{Structural comparison: quantum mechanics vs.\ the present axiomatic system.\label{tab:qm-cas}}
\small
\begin{tabular}{lll}
\toprule
\textbf{Dimension} & \textbf{Quantum mechanics} & \textbf{Present system} \\
\midrule
Measurement & Projection operator $P_k$ & Read $\to$ Compare $\to$ Swap \\
Probability/cost rule & $|\langle k|\psi\rangle|^2$ & Cost $\{0, +1\}$ \\
State space & Continuous Hilbert space & Discrete FSM (7-bit) \\
Collapse mechanism & Wavefunction collapse & $\delta$-firing screen rendering completion \\
Irreversibility origin & Measurement postulate & CAS order forcing (Axiom~\ref{ax:cas}) \\
\bottomrule
\end{tabular}
\end{table}

The two descriptions are not in a translation relation. They are \emph{different-axis projections} of the same natural phenomenon---projecting the Banya equation onto the classical bracket (DATA) yields the Hilbert-space description, and projecting onto the quantum bracket (OPERATOR) yields the CAS description. The question ``which is real?'' is itself ill-posed---both projections are different aspects of the same $\delta$.

More concretely, in the branching structure of CAS from Axioms~\ref{ax:cas} and \ref{ax:cost}: if Compare returns true, Swap executes, paying cost $+1$, and a definite state is written to DATA---this corresponds to wavefunction collapse. If Compare returns false, Swap does not execute and superposition is maintained---cost $0$, reversible. That is, superposition maintenance $=$ Compare false $=$ cost $0$, and collapse $=$ Compare true $=$ cost $+1$. The core of this structure is that quantum is the default; classical results from paying cost.

\emph{Correspondence with quantum non-demolition (QND) measurement}: in the complete axiomatic system, the execution condition for Swap is twofold---Compare true \emph{AND} isWriteAble~\cite{han2026}. When Compare is true but isWriteAble is false, the state has been observed but not altered. This is the structural correspondence to QND measurement.

\subsubsection{How does the continuous emerge from the discrete?}

(\emph{Note}: this subsubsection is related to the ``Transition from discrete axioms to continuous physical quantities'' subsection of Section~\ref{sec:discussion}, but the two address \emph{different aspects} --- the earlier subsection concerns the \emph{mathematical} transition (Theorem 1 $\to$ Lie group $\mathrm{SO}(5,2)$ $\to$ continuous volume ratio), while this subsubsection concerns the \emph{physical} emergence (cost recovery by RLU proceeds continuously in the continuous domain of the OPERATOR bracket, generating discrete records in DATA).)

Contention creates order, and order creates discreteness (Section~\ref{sec:framework}). When multiple local states exist simultaneously, causality preservation forces an ordering device (lock) to operate, and cost $+1$ arises wherever the CAS sequence $\mathrm{R} \to \mathrm{C} \to \mathrm{S}$ is forced; when cost arises, it is recorded in DATA in discrete slot units.

Where, then, does the continuous come from? It comes from the OPERATOR bracket. In Axiom~\ref{ax:banya}, the OPERATOR (quantum bracket) is a continuous domain in which all states coexist simultaneously. The cost accumulated by CAS (Axiom~\ref{ax:cost}) is recovered at the OPERATOR level according to the eviction rule of RLU (analogous to LRU), and this recovery process itself is continuous. When the continuous cost distribution passes through CAS Swap and is written to DATA, it is discretized (Axiom~\ref{ax:cost}: crossing $+$ triggers discretization). In the present axiomatic system, no quantity called ``force'' is defined --- the only physical quantity is \emph{change}, and the expression of change is \emph{cost}.

The differential calculus used in physics is an approximation of this process: when sufficiently many discrete records (local states) accumulate, the discrete distribution appears continuous and the differential becomes a valid approximation. This is the same logic by which the continuum limit is recovered in lattice gauge theory when the lattice spacing is sufficiently small. The difference is: in the present system, discreteness is not ``the starting point of an approximation'' but ``a structural consequence generated by ordering.'' Causal set theory likewise pursues a program in which continuous spacetime emerges from discrete causal structure.

\subsubsection{Structural correspondence of superposition and collapse}

We make explicit the structural correspondence between quantum-mechanical state vectors (Hilbert space) and the present system.

In Axiom~\ref{ax:banya}, the superposition of the quantum bracket (OPERATOR) is ``an addressless index spread simultaneously over all states.'' This is structurally analogous to a state vector $|\psi\rangle = \sum c_k |k\rangle$ in Hilbert space being spread simultaneously over all basis states $|k\rangle$. This ``addressless index'' property is made concrete by the RLU indexing mechanism (defined immediately after Axiom~\ref{ax:cas}) as \emph{angular coordinate (theta angle) identification} --- RLU distinguishes targets by angular position in the OPERATOR bracket without logical addresses, which is why the indexing of the present axiomatic system distributes cost over the arc length $\pi$.

The cost rule of Axiom~\ref{ax:cost} provides the ``collapse'' mechanism. If CAS Compare returns true, cost $+1$ is paid and Swap executes, writing a definite state to DATA---this corresponds to wavefunction collapse. If Compare returns false, Swap does not execute and superposition is maintained---cost $0$, reversible. That is:
\begin{itemize}
\item Superposition maintained $=$ Compare false $=$ cost $0$ $=$ remains in OPERATOR.
\item Collapse (write) $=$ Compare true $=$ cost $+1$ $=$ written to DATA.
\end{itemize}

\emph{Structural isomorphism with the Born rule --- derived in the complete system}. The 4 axioms of this paper do not \emph{directly} derive the Born rule; this paragraph is recorded as an out-of-scope observation pointing toward the complete axiomatic system.

\emph{Match in formula structure}. In standard quantum mechanics, the measurement probability is $P(k) = |c_k|^2 = a_k^2 + b_k^2$ (complex amplitude $c_k = a_k + i b_k$), which is the \emph{norm squared of two orthogonal real components}. The Banya equation (Axiom~\ref{ax:banya}) $\delta^2 = (\text{time}+\text{space})^2 + (\text{observer}+\text{superposition})^2$ is a \emph{sum of norm squares of two orthogonal categories}, and its formula structure \emph{matches} the form $|c_k|^2 = a_k^2 + b_k^2$ at the $n = 2$ level (two categories/brackets).

\emph{Scope delimitation for this paper}. This paper focuses on the \emph{self-contained derivation of $\alpha$} within the 4 axioms + 1 proposition. The formal derivation of the Born rule is completed in the full system~\cite{han2026} (Axiom~15 $\to$ Axiom~3 $\to$ Axiom~2, a CAS version of Gleason's theorem). Because Axiom~15 is not among the 4 axioms of this paper, the Born rule lies outside the 4-axiom scope.

\emph{Related prior art and formal differences}. The ``measurement problem'' of quantum mechanics has remained unresolved since the founding of quantum mechanics, and attempts to render the Born rule a \emph{derivable consequence} rather than a \emph{basic postulate} constitute an active research area. The present axiomatic system is positioned \emph{adjacent to} both the \emph{Hardy/CDP full reconstruction program} and the \emph{Spekkens/Masanes/Pawlowski partial reconstruction program} without coinciding directly with either, occupying a \emph{cost-theoretic branch}.

(i) \emph{Full reconstruction program}: Hardy's ``Quantum Theory From Five Reasonable Axioms''~\cite{hardy2001} \emph{fully reconstructs} finite-dimensional QM (state space, measurement, time evolution) from 5 operational axioms. Chiribella, D'Ariano, and Perinotti~\cite{cdp2011} add ``purification'' to 6 informational axioms to obtain the same result. Hardy's fifth axiom\footnote{Hardy 2001 original: ``Continuity: there exists a continuous reversible transformation on a system between any two pure states of that system.''} (``continuous reversible transformation between pure states'') conflicts with the CAS 3-stage irreversible progression of the present axiomatic system in the same way as Masanes--M\"uller's continuous reversibility, and this constitutes a categorical boundary common to \emph{both Hardy and Masanes--M\"uller}.

(ii) \emph{Partial reconstruction program}: Spekkens' epistemic toy model~\cite{spekkens2007} \emph{partially reproduces} a large part of QM phenomenology (no-cloning, interference, teleportation, entanglement monogamy) from 2 principles (knowledge balance + 4-state classical). Masanes--M\"uller~\cite{masanes2011} derive QM from 5 physical requirements. Pawlowski et al.'s information causality~\cite{pawlowski2009} derives QM information bounds from a single principle.

\emph{Spekkens 4-state vs.\ the present axiomatic system's $4+3$: structural comparison}. Spekkens' ``4 states'' is the \emph{ontic state space of a single-element system} (e.g., 4 epistemic states that a coin can take), with the knowledge balance principle forcing 1 question of knowledge per system. The present axiomatic system's ``$4+3 = 7$'' is the internal DOF of a \emph{single CAS cycle} (4 domain axes + 3 CAS stages), and the two 4s belong to different categories --- Spekkens' 4 is an \emph{external ontic state count}, while the axiomatic system's 4 is an \emph{internal degree-of-freedom count}. The \emph{numerical coincidence} of the two 4s is not accidental but arises from the \emph{same binary partition structure} of ``2-bit ontic state space'' and ``2-dimensional DATA bracket + 2-dimensional OPERATOR bracket''; this is a comparison target for subsequent work; a detailed comparison is reserved for a dedicated study. Similarly, the \emph{continuous reversibility} principle among Masanes--M\"uller's 5 requirements~\cite{masanes2011} directly conflicts with the CAS atomicity (3-stage irreversible progression) of the present axiomatic system, revealing the essential difference that the present axiomatic system starts from a discrete CAS cycle rather than a connected continuous transformation group.

(iii) \emph{Position of the present axiomatic system}: the present axiomatic system coincides directly with neither program --- the \emph{reconstruction target} differs. Hardy/CDP reconstruct the \emph{formal structure} of QM; Spekkens/Masanes reproduce \emph{parts of QM phenomenology}; the present axiomatic system provides the \emph{single coupling constant $\alpha$ and its cost-theoretic motivation}. That is, ``the reconstruction target is theory $\to$ phenomenology $\to$ single constant'' --- three different categories, and the present axiomatic system operates from a \emph{third starting point}: the \emph{cost classification $\leftrightarrow$ metric signature correspondence}.

(iv) \emph{Formal limitations}: the present axiomatic system operates outside (a) the operational framework, (b) GPT (generalized probabilistic theory), and (c) convex/probabilistic structure --- compared with Hardy/CDP's \emph{full reconstruction}, it is categorically different, occupying a \emph{cost-theoretic branch} that uniquely produces physical constants.

(v) Hardy/CDP's axioms are all directly reducible to \emph{operational measurement outcomes}, whereas the CAS axioms of the present system are \emph{computer-science abstractions} whose operational motivation relies on the structural isomorphism with the Landauer bound~\cite{landauer1961} (Section~\ref{sec:axioms}). The formula-structure isomorphism in this paragraph is an out-of-scope observation toward a formal bridge with the complete axiomatic system, and a precise comparison with the above programs will appear in dedicated papers.

\subsubsection{Uniqueness of the cost function $\{0, +1\}$}

Why is the cost exactly $\{0, +1\}$ and not $+0.5$ or $+2$?

Why $+0.5$ is impossible: each CAS stage is atomic (indivisible) (Axiom~\ref{ax:cas}). Atomicity means that intermediate states are unobservable---splitting a stage that cannot be divided externally into 0.5 internally contradicts the definition of atomicity. This argument rests not on ``DATA is discrete'' but on ``CAS stages are atomic.'' The discreteness of DATA itself is a declaration of Axiom~\ref{ax:banya}, but atomicity is a structural property independently defined in Axiom~\ref{ax:cas}.

Why $+2$ or more is impossible: a single $+$ traversal is exactly one orthogonal-axis transition (Axiom~\ref{ax:cost}). The 3 CAS axes are mutually orthogonal (Axiom~\ref{ax:cas}: $\mathrm{R} \perp \mathrm{C} \perp \mathrm{S}$), and the ordering device (lock) permits only one transition at a time. If two axes were activated simultaneously in a single transition, orthogonality would be broken and the order dependence ``compare cannot execute without read'' (Axiom~\ref{ax:cas}) would be violated. Therefore the cost per $+$ traversal is exactly $+1$.

Result: the range of the cost function is $\{0, +1\}$, and this form is forced by Axiom~\ref{ax:banya} (DATA discrete), Axiom~\ref{ax:cas} (atomicity + orthogonality), and Axiom~\ref{ax:cost} (order $=$ cost). Any other form of cost function violates at least one of these three axioms.

\begin{table}[H]
\begin{adjustwidth}{-\extralength}{0cm}
\caption{Complete results of the axiomatic system --- results produced by 4 axioms and 1 proposition.\label{tab:axiom-results}}
\scriptsize
\begin{tabularx}{\fulllength}{cp{3.2cm}X}
\toprule
\textbf{Classification} & \textbf{Result} & \textbf{Remarks} \\
\midrule
\textbf{I} & $(5,2)$ signature uniqueness & Axioms 1--4 $\to$ Theorem 1 (§\ref{sec:theorem}) \\
\textbf{I} & $7 = 4+3$ internal degrees of freedom & Axioms 1+2 $\to$ Proposition 1. Hamming $[7,4,3]$ coincidence (§\ref{sec:axioms}) \\
\textbf{I} & Cost function $\{0,+1\}$ uniqueness & Axioms 1,2,3: atomicity+orthogonality+order (§\ref{sec:discussion}) \\
\textbf{I} & Exclusion of SO(7), SO(4,3), and 6 others & Theorem 1 Part (C) (§\ref{sec:theorem}) \\
\textbf{I} & $1/\alpha = 137.036\,082$ & Theorem 1 $\to$ $D_5$ $\to$ Wyler volume ratio. $6 \times 10^{-7}$ vs.\ CODATA (§\ref{sec:derivation}) \\
\textbf{I} & $\sin^2\theta_W = 0.23122$ & 5-step forward chain. $4 \times 10^{-6}$ vs.\ $\overline{\text{MS}}$ at $M_Z$ (§\ref{sec:prediction}) \\
\textbf{I} & $9$, $13 = 8+5$, ball value $4 = 1+3$ & Residual cost $13-4$, cost accounting (Definition~\ref{lem:cost13}, Definition~\ref{def:residual}) \\
\textbf{I} & $\pi$, $9\pi$, $7/(2+9\pi)$ uniqueness & CAS orthogonality $\to$ cost+norm reading + arithmetic (§\ref{sec:prediction}) \\
\textbf{I} & Independence of $\alpha$ and $\sin^2\theta_W$ & Distinct forward chains (§\ref{sec:prediction}) \\
\textbf{I} & 3-dimensional space ($x+y+z$) & Observational decomposition of the space axis declared by Axiom 1. Explicit in Definition~\ref{lem:cost13} (§\ref{sec:axioms}) \\
\textbf{I} & $\alpha$ = interaction probability & Theorem 1 $\to$ $D_5$ $\to$ Shilov boundary/$D_5$ volume ratio $=$ bracket-boundary $+$ crossing probability. Direct reading in axiomatic terminology (§\ref{sec:derivation}) \\
\textbf{I} & 13 boundary distinctions & $8$ reads $+$ $5$ writes. Unique enumeration forced by axioms (Definition~\ref{lem:cost13}) \\
\textbf{I} & Energy scale $\leftrightarrow$ domain time & Immediate consequence of Axiom~\ref{ax:banya}: $\delta$ (system time) $\neq$ time axis (domain time). Axiom 3 of complete system~\cite{han2026} further specifies (§\ref{sec:axioms}) \\
\midrule
\textbf{II} & $\pi^5/(2^4 \cdot 5!) = D_5$ Hua volume & $(5,2) \to D_5 \to$ Hua Ch.\,IV identity. Mathematical identity (§\ref{sec:derivation}) \\
\textbf{II} & $D_5$, $\mathrm{SO}(5)\times\mathrm{SO}(2)$ determined & $(5,2) \to$ unique bounded symmetric domain (§\ref{sec:derivation}) \\
\textbf{II} & $8\pi^4$ & Axiom-determined from $(5,2) \to \mathrm{SO}(5)\times\mathrm{SO}(2)$ group volume (Table~\ref{tab:factors}) \\
\textbf{II} & $(\cdot)^{1/4}$ geometric mean of 4 orthogonal axes & Axiom~\ref{ax:banya} orthogonality $\to$ independence $\to$ geometric mean $= (\cdot)^{1/4}$. Mathematically axiom-determined (§\ref{sec:derivation}) \\
\textbf{II} & Hartley information $\Delta H = 5$ bits & Axiom~\ref{ax:banya} (orthogonal independence) + Theorem~\ref{thm:cost-signature} (5 irreversible) $\to$ Hartley formula: $\log_2 128 - \log_2 4 = 5$. Mathematically axiom-determined (§\ref{sec:framework}) \\
\textbf{II} & Hamming $[7,4,3]$ structural correspondence & Proposition 1: $4{+}3{=}7$ (I) $\to$ Hamming bound: unique perfect code for 4-bit messages $= [7,4,3]$ (mathematical theorem). Same pattern as Hua volume (§\ref{sec:predictions}) \\
\textbf{II} & Steane $[[7,1,3]]$ quantum extension & Hamming $[7,4,3]$ (II) $\to$ CSS construction (Calderbank--Shor--Steane 1996; automatic mathematical procedure) $\to$ $[[7,1,3]]$. No choice (§\ref{sec:predictions}) \\
\textbf{II} & CAS $\leftrightarrow$ Landauer erasure & CAS Swap overwrites previous state (Axiom~\ref{ax:cas}) $=$ Landauer information erasure by definition. Structurally isomorphic (§\ref{sec:framework}) \\
\textbf{II} & $\eta_B$ prefactor 2 & Bracket count $= 2$ (Axiom~\ref{ax:banya}); orthogonality $=$ independence $\to$ multiplicative combination. Same pattern as $(\cdot)^{1/4}$ geometric mean (§\ref{sec:baryogenesis}) \\
\midrule
\textbf{III} & $\eta_B = 6.14 \times 10^{-10}$ & $\alpha^4 \sin^2\theta_W [1{-}2(4{+}1/\pi)\alpha]$. Formula and all inputs complete in text. $0.5\sigma$ vs.\ Planck. Only RLU formal systematization in~\cite{han2026} (§\ref{sec:baryogenesis}) \\
\textbf{III} & Quantum measurement $\leftrightarrow$ R$\to$C$\to$S & Compare true $\to$ Swap $\to$ DATA record (§\ref{sec:discussion}) \\
\textbf{III} & Born rule $|c_k|^2$ formal coincidence & $\delta^2 \leftrightarrow |c_k|^2 = a_k^2+b_k^2$ (§\ref{sec:discussion}) \\
\textbf{III} & R,C,S $\to$ U(1),SU(2),SU(3) & CAS 3-step $\to$ 3 gauge conservation types (§\ref{sec:discussion}) \\
\textbf{III} & $128 = 2^7 \to 137 = 128+9$ & 7-bit FSM + residual cost 9. QED RG integer part coincidence (§\ref{sec:discussion}) \\
\textbf{III} & SU(3)$\times$SU(2)$\times$U(1) correspondence & R,C,S $\to$ U(1),SU(2),SU(3) mapping complete. 12 gauge bosons derived. Coupling ratio quantification in~\cite{han2026} (§\ref{sec:discussion}) \\
\textbf{III} & Alternative $\sin^2\theta_W$ candidates exist & $1/\log_2 20$ etc.\ recorded. Cost enumeration path adopted (§\ref{sec:prediction}) \\
\textbf{III} & CAS 3-step uniqueness & Lower bound (causality) + upper bound (least cost) force 3 steps (Section~\ref{sec:framework}). Only complete external classification reserved for dedicated paper (§\ref{sec:framework}) \\
\textbf{III} & $\sin^2\theta_W$/$\eta_B$ form difference & Raw ($9\pi$) vs.\ coupled ($(4{+}1/\pi)$). Not contradictory (§\ref{sec:prediction}, §\ref{sec:baryogenesis}) \\
\textbf{III} & $\alpha$ precision $\Delta = 0.000\,083$ & $\alpha^3 \approx 4 \times 10^{-7}$ order correction possible (§\ref{sec:limitations}) \\
\textbf{III} & Landauer heat lower bound $5 \times kT\ln 2$ & Minimum heat dissipation per CAS cycle when tier-2 structural correspondence holds (§\ref{sec:framework}) \\
\textbf{III} & QND measurement correspondence & Compare true $+$ isWriteAble false $=$ observed but not altered (§\ref{sec:discussion}) \\
\textbf{III} & CAS running mechanism & $\alpha(N) = \alpha/(1-\alpha N/(3\pi))$. $\beta$ denominator $3$ = CAS stages; QCD numerator $7$ = complete DOF. Prediction: $\beta_2$, $\beta_3$ contain CAS numbers (§\ref{sec:running}) \\
\textbf{III} & Sakharov 3 conditions = CAS embedded & B violation $\leftarrow$ Swap domain crossing (Axiom~\ref{ax:cost}); CP violation $\leftarrow$ CAS irreversibility (Axiom~\ref{ax:cas}); non-equilibrium $\leftarrow$ RLU delay (§\ref{sec:baryogenesis}) \\
\textbf{III} & $J_\text{CKM} = 3.10 \times 10^{-5}$ & CAS irreversible phase volume. Exp.\ $(3.08 \pm 0.15) \times 10^{-5}$, deviation $0.62\%$~\cite{han2026} (§\ref{sec:baryogenesis}) \\
\midrule
\textbf{IV} & Rule selection problem & Open meta-problem common to digital physics in general (§\ref{sec:limitations}) \\
\bottomrule
\end{tabularx}
\parbox{\fulllength}{\emph{Classification criteria} --- I: Formally derived from axioms. II: The axioms force the structure, which is then mathematically determined. III: Derivation path and interpretation given in the text; formal systematization reserved for dedicated papers. IV: Not included in the scope of the present paper. The classification criteria for all subsequent tables are the same as in this table.}
\end{adjustwidth}
\end{table}

%%%%%%%%%%%%%%%%%%%%%%%%%%%%%%%%%%%%%%%%%%
\section{Conclusion}
\label{sec:conclusion}

Four axioms and one proposition independently produce the fine-structure constant $\alpha$ and the Weinberg angle $\sin^2\theta_W$; the baryon-to-photon ratio $\eta_B$ is obtained as an independent forward-testable production at axiom-described level (classification~III). \textbf{We show that axioms describe nature.}

Within an axiomatic system whose primitive object is the CAS operation, we have traced a structural chain that reaches the low-energy limit $\alpha = 1/137.036\,082$ of the fine-structure constant. All four factors of the Wyler formula are accounted for within the axiomatic structure: $9$ is axiom-derived~(I), $\pi^5/(2^4 \cdot 5!)$ is axiom-determined~(II), $(\cdot)^{1/4}$ is axiom-determined as the geometric mean of 4 independent axes~(II), and $8\pi^4$ is the group volume fixed by the signature $(5,2)$ and is therefore mathematically axiom-determined~(II).

The causal chain of the derivation is: contention $\to$ ordering $\to$ discreteness $\to$ cost classification $\to$ signature $(5,2)$ $\to$ $\mathrm{SO}_0(5,2) \to D_5 \to$ Wyler formula (each factor accounted for at its epistemic level) $\to \alpha$. The key step is the cost--signature correspondence (Definition~\ref{def:cost-sig} and Theorem~\ref{thm:cost-signature}): the irreversible/reversible classification of costs maps naturally onto the positive and negative signs of a quadratic form, and because this classification is determined by the axioms, the signature $(5,2)$ follows. The result is $\alpha = 1/137.036\,082$ (relative deviation $6 \times 10^{-7}$ from the experimental value).

The three criticisms of Robertson~(1971) and Gilmore~(1972)---absence of physical motivation, arbitrariness of the group choice, and non-uniqueness of the measure---have been answered item by item. In particular, part~(C) of Theorem~\ref{thm:cost-signature} formally excludes all partitions other than $(5,2)$ among the eight candidates, and among $21$-dimensional simple Lie groups we have explicitly ruled out $(7,0)$ ($\mathrm{SO}(7)$, the compact form) and $(4,3)$ ($\mathrm{SO}(4,3)$, the split form) on the grounds that they are incompatible with the principle of deriving the metric sign from the \emph{cost category}---this constitutes the direct reply to Robertson's~(R2) ``arbitrariness of the group choice'' criticism.

A different combination of structural constants within the same axiomatic system yields the Weinberg angle $\sin^2\theta_W = 7/(2+9\pi) = 0.23122$ (relative deviation $4 \times 10^{-6}$ from the experimental value). The distinction between the energy scale and the domain time is an immediate consequence of Axiom~\ref{ax:banya} (classification~I; see Section~\ref{sec:axioms}).

Furthermore, in Section~\ref{sec:predictions} the internal degree-of-freedom count $7 = 4 + 3$ of the present axiomatic system was shown to coincide structurally with the decomposition into four information bits and three parity bits of the classical Hamming code $[7,4,3]$ (the quantum-information community's Steane code $[[7,1,3]]$ extends naturally as a CSS construction built on this Hamming code).

Finally, in Section~\ref{sec:baryogenesis} the baryon-to-photon ratio $\eta_B = 6.14 \times 10^{-10}$ was produced from the two quantities derived in the present paper, $\alpha$ and $\sin^2\theta_W$ (both forward chains) alone, and was shown to agree with the Planck measurement to within $0.5\sigma$---a forward item that can be verified or falsified by future CMB precision measurements (the $+$ in $(4+1/\pi)$ is uniquely determined by Axiom~\ref{ax:banya}'s orthogonal composition; the formal systematization of the full RLU mechanism is classification~III).

The $\alpha$ derivation process of this paper yields 5 unique predictions (Table~\ref{tab:predictions}). These predictions are not fitted to match $\alpha$; they are deductive consequences that automatically follow from the same CAS cost structure. Each prediction has a specific experiment and an explicit falsification condition. Falsification of any one negates the corresponding CAS structure.

\begin{table}[H]
\caption{Unique predictions of this paper --- deductive consequences of the CAS cost structure.\label{tab:predictions}}
\begin{adjustwidth}{-\extralength}{0cm}
\small
\begin{tabularx}{\fulllength}{p{0.4cm}Xp{3.0cm}p{4.0cm}}
\toprule
\textbf{\#} & \textbf{Prediction (axiomatic basis)} & \textbf{Source} & \textbf{Falsification condition} \\
\midrule
1 & \textbf{CAS structural numbers in higher-loop $\beta$-coefficients ($\overline{\mathrm{MS}}$, $\mu = M_Z$).} 1-loop: QED denominator $3$ = CAS stages; QCD numerator $7$ = complete descriptive DOF. 2-loop: $\beta_1 = -1/4 = -1/$(domain~4). The prediction: under $\overline{\mathrm{MS}}$ scheme, $\mu = M_Z$, with standard PDG normalization, the numerators and denominators of the rational form of $\beta_2, \beta_3$ must be expressible as products/sums/differences of the CAS set $\{3, 4, 7, 9, 13\}$ (no spontaneous appearance of primes outside this set). & §\ref{sec:running} & A prime $\geq 11$ (excluding 13) appears spontaneously in the standard rational form of $\overline{\mathrm{MS}}$ $\beta_2$ or $\beta_3$ \\
\midrule
2 & \textbf{Integer part of $1/\alpha(M_Z) = 128 = 2^7$.} CAS running exhausts residual cost $9$ via Compare-false accumulation: $137 - 9 = 128$. $128 = 2^7$ is the FSM ground-state count. ``Ground states + residual cost = low-energy limit.'' & §\ref{sec:discussion} & Precision measurement of $\alpha(M_Z)$ yields nearest integer $\neq 128$ \\
\midrule
3 & \textbf{Jarlskog invariant $J_\text{CKM} = 3.10 \times 10^{-5}$.} CP violation magnitude is determined by the phase volume of the irreversible CAS cycle (R$\to$C$\to$S, reverse undefined). Produced in the complete system~\cite{han2026}. Experimental value $(3.08 \pm 0.15) \times 10^{-5}$, deviation $0.62\%$. & §\ref{sec:baryogenesis}, \cite{han2026} & Experimental value deviates from $3.10 \times 10^{-5}$ by $> 3\sigma$ \\
\midrule
4 & \textbf{No 4th gauge interaction ($Z'$, etc.).} CAS = 3 stages (R, C, S). Read$\to$U(1), Compare$\to$SU(2), Swap$\to$SU(3). No 4th stage $\Rightarrow$ no 4th gauge group. Total gauge bosons $= 12 = \binom{4}{2} \times 2$. & §\ref{sec:discussion} & Discovery of $Z'$ or $W'$ below 10~TeV at LHC Run 3/4 \\
\midrule
5 & \textbf{Sakharov's 3 conditions embedded in CAS.} B violation $\leftarrow$ FSM Swap domain crossing (Axiom~\ref{ax:cost}). CP violation $\leftarrow$ CAS irreversibility (Axiom~\ref{ax:cas}). Non-equilibrium $\leftarrow$ RLU cost-recovery delay. Three conditions from three different axioms. & §\ref{sec:baryogenesis} & Any of the 3 conditions requires an independent mechanism outside CAS \\
\bottomrule
\end{tabularx}
\end{adjustwidth}
\end{table}

Table~\ref{tab:scope} compares the scope of the present paper (4~axioms + 1~proposition) with that of the complete axiomatic system~\cite{han2026} (15~axioms). The present paper addresses only the minimal axiom subset required for the $\alpha$ derivation. The complete axiomatic system~\cite{han2026} has already produced coupling-constant hierarchies ($\alpha_s = 0.1183$, deviation $0.3\%$), all 6 quark masses, the cosmological constant ($\Lambda l_p^2 \sim \alpha^{57}$), and others; their formal academic presentation will appear in dedicated papers.

\begin{table}[H]
\caption{Scope of derivations --- this paper versus the complete system.\label{tab:scope}}
\begin{tabular}{lrr}
\toprule
\textbf{Item} & \textbf{This paper} & \textbf{Complete system~\cite{han2026}} \\
\midrule
Number of axioms & 4 + 1 prop. & 15 \\
\midrule
\multicolumn{3}{l}{\emph{Derivation items (Table~\ref{tab:axiom-results} classification)}} \\
\quad Experimentally matched ($<1\%$) & 3 & 155 \\
\quad Structural/mathematical results & 35 & --- \\
\quad Awaiting verification & --- & 941 \\
\quad Subtotal & 38 & 1,096 \\
\midrule
\multicolumn{3}{l}{\emph{Unique predictions (experimentally testable)}} \\
\quad Experimental match & --- & 18 \\
\quad Awaiting verification & 5 & 111 \\
\quad Hypothesis & --- & 1 \\
\quad Subtotal & 5 & 130 \\
\bottomrule
\end{tabular}
\end{table}

%%%%%%%%%%%%%%%%%%%%%%%%%%%%%%%%%%%%%%%%%%
\vspace{6pt}

%%%%%%%%%%%%%%%%%%%%%%%%%%%%%%%%%%%%%%%%%%
\authorcontributions{Conceptualization, methodology, formal analysis, investigation, writing of the original draft, and editing: Hyukjin Han. The author has read and agreed to the published version of the manuscript.}

\funding{This research received no external funding.}

\institutionalreview{Not applicable.}

\informedconsent{Not applicable.}

\dataavailability{No new data were created or analyzed in this study. The complete axiomatic system and all derivations are available at \url{https://doi.org/10.5281/zenodo.19850973}.}

\acknowledgments{During the preparation of this manuscript, the author used Claude (Anthropic) as an aid for \LaTeX{} typesetting and manuscript formatting. All axioms, theorems, proofs, derivations, and conclusions were developed by the author, who takes full responsibility for the content of this publication.}

\conflictsofinterest{The author declares no conflict of interest.}

%%%%%%%%%%%%%%%%%%%%%%%%%%%%%%%%%%%%%%%%%%
\abbreviations{Abbreviations}{
The following abbreviations are used in this manuscript:
\\

\noindent
\begin{tabular}{@{}ll}
CAS & Compare-And-Swap (conditional state transition) \\
DOF & Degrees of freedom \\
QED & Quantum electrodynamics \\
R, C, S & Read, Compare, Swap \\
RLU & Addressless virtual indexing unit \\
FSM & Finite state machine
\end{tabular}
}

%%%%%%%%%%%%%%%%%%%%%%%%%%%%%%%%%%%%%%%%%%
\appendixstart
\appendix
\section[\appendixname~\thesection]{}
\subsection{Explicit computation of the Wyler formula}
\label{app:computation}

Step-by-step numerical evaluation of Eq.~\eqref{eq:wyler-formula}:

$9/(8\pi^4) = 9/779.273 = 0.011\,549$. \quad $\pi^5/(2^4 \cdot 5!) = 306.020/1920 = 0.159\,385$. \quad $(0.159\,385)^{1/4} = 0.631\,85$. \quad $\alpha = 0.011\,549 \times 0.631\,85 = 0.007\,297\,4$. \quad $1/\alpha = 137.036\,082$.

%%%%%%%%%%%%%%%%%%%%%%%%%%%%%%%%%%%%%%%%%%
\begin{adjustwidth}{-\extralength}{0cm}

\reftitle{References}

\begin{thebibliography}{999}

\bibitem{feynman1985}
Feynman,~R.P. \textit{QED: The Strange Theory of Light and Matter}; Princeton University Press: Princeton, NJ, USA, 1985; p.~129.

\bibitem{wyler1969}
Wyler,~A. L'espace sym\'{e}trique du groupe des \'{e}quations de Maxwell. \textit{C. R. Acad. Sci. Paris} \emph{1969}, \textit{269A}, 743--745.

\bibitem{wyler1971}
Wyler,~A. Les groupes des potentiels de Coulomb et de Yukawa. \textit{C. R. Acad. Sci. Paris} \emph{1971}, \textit{272A}, 186--188.

\bibitem{robertson1971}
Robertson,~B. Wyler's expression for the fine-structure constant. \textit{Phys. Rev. Lett.} \emph{1971}, \textit{27}, 1545--1547.

\bibitem{gilmore1972}
Gilmore,~R. Scaling of Wyler's expression for $\alpha$. \textit{Phys. Rev. Lett.} \emph{1972}, \textit{28}, 462--464.

\bibitem{helgason1978}
Helgason,~S. \textit{Differential Geometry, Lie Groups, and Symmetric Spaces}; Academic Press: New York, NY, USA, 1978.

\bibitem{hua1963}
Hua,~L.K. \textit{Harmonic Analysis of Functions of Several Complex Variables in the Classical Domains}; American Mathematical Society: Providence, RI, USA, 1963.

\bibitem{codata2022}
Tiesinga,~E.; Mohr,~P.J.; Newell,~D.B.; Taylor,~B.N. CODATA Recommended Values of the Fundamental Physical Constants: 2022. \textit{J. Phys. Chem. Ref. Data} \emph{2024}, \textit{53}, 031301.

\bibitem{eddington1946}
Eddington,~A.S. \textit{Fundamental Theory}; Cambridge University Press: Cambridge, UK, 1946.

\bibitem{han2026}
Han,~H. Banya Framework: Axiom-Based Science Mining Engine---Comprehensive Report, v1.6. 2026. Available online: \url{https://doi.org/10.5281/zenodo.19850973} (accessed on 28 April 2026).

\bibitem{pdg2024}
Workman,~R.L.; et~al. Review of Particle Physics. \textit{Phys. Rev. D} \emph{2024}, \textit{110}, 030001.

\bibitem{koide1983}
Koide,~Y. New view of quark and lepton mass hierarchy. \textit{Phys. Rev. D} \emph{1983}, \textit{28}, 252--254.

\bibitem{zuse1969}
Zuse,~K. \textit{Rechnender Raum (Calculating Space)}; MIT Technical Translation, 1970.

\bibitem{fredkin2003}
Fredkin,~E. An introduction to digital philosophy. \textit{Int. J. Theor. Phys.} \emph{2003}, \textit{42}, 189--247.

\bibitem{lloyd2002}
Lloyd,~S. Computational capacity of the universe. \textit{Phys. Rev. Lett.} \emph{2002}, \textit{88}, 237901.

\bibitem{wheeler1990}
Wheeler,~J.A. Information, physics, quantum: The search for links. In \textit{Complexity, Entropy, and the Physics of Information}; Zurek, W.H., Ed.; Addison-Wesley: Redwood City, CA, USA, 1990; pp.~3--28.

\bibitem{landauer1961}
Landauer,~R. Irreversibility and heat generation in the computing process. \textit{IBM J. Res. Dev.} \emph{1961}, \textit{5}, 183--191.

\bibitem{duff2002}
Duff,~M.J.; Okun,~L.B.; Veneziano,~G. Trialogue on the number of fundamental constants. \textit{J. High Energy Phys.} \emph{2002}, \textit{2002}, 023.

\bibitem{bars2006}
Bars,~I. Two-time physics in field theory. \textit{Phys. Rev. D} \emph{2006}, \textit{74}, 085019.

\bibitem{connes1996}
Connes,~A. Gravity coupled with matter and the foundation of non-commutative geometry. \textit{Commun. Math. Phys.} \emph{1996}, \textit{182}, 155--176.

\bibitem{lisi2007}
Lisi,~A.G. An exceptionally simple theory of everything. arXiv:0711.0770, 2007.

\bibitem{distler2010}
Distler,~J.; Garibaldi,~S. There is no ``Theory of Everything'' inside $E_8$. \textit{Commun. Math. Phys.} \emph{2010}, \textit{298}, 419--436.

\bibitem{morel2020}
Morel,~L.; Yao,~Z.; Clad\'{e},~P.; Guellati-Kh\'{e}lifa,~S. Determination of the fine-structure constant with an accuracy of 81 parts per trillion. \textit{Nature} \emph{2020}, \textit{588}, 61--65.

\bibitem{thooft2016}
't~Hooft,~G. \textit{The Cellular Automaton Interpretation of Quantum Mechanics}; Springer: Cham, Switzerland, 2016.

\bibitem{wolfram2002}
Wolfram,~S. \textit{A New Kind of Science}; Wolfram Media: Champaign, IL, USA, 2002.

\bibitem{tegmark2008}
Tegmark,~M. The mathematical universe. \textit{Found. Phys.} \emph{2008}, \textit{38}, 101--150.

\bibitem{bennett1973}
Bennett,~C.H. Logical reversibility of computation. \textit{IBM J. Res. Dev.} \emph{1973}, \textit{17}, 525--532.

\bibitem{herlihy1991}
Herlihy,~M. Wait-free synchronization. \textit{ACM Trans. Program. Lang. Syst.} \emph{1991}, \textit{13}, 124--149.

\bibitem{zurek2009}
Zurek,~W.H. Quantum Darwinism. \textit{Nat. Phys.} \emph{2009}, \textit{5}, 181--188.

\bibitem{hamming1950}
Hamming,~R.W. Error detecting and error correcting codes. \textit{Bell Syst. Tech. J.} \emph{1950}, \textit{29}, 147--160.

\bibitem{steane1996}
Steane,~A.M. Error Correcting Codes in Quantum Theory. \textit{Phys. Rev. Lett.} \emph{1996}, \textit{77}, 793--797.

\bibitem{hardy2001}
Hardy,~L. Quantum Theory From Five Reasonable Axioms. \textit{arXiv preprint} \emph{2001}, quant-ph/0101012.

\bibitem{cdp2011}
Chiribella,~G.; D'Ariano,~G.M.; Perinotti,~P. Informational derivation of quantum theory. \textit{Phys. Rev. A} \emph{2011}, \textit{84}, 012311.

\bibitem{spekkens2007}
Spekkens,~R.W. Evidence for the epistemic view of quantum states: A toy theory. \textit{Phys. Rev. A} \emph{2007}, \textit{75}, 032110.

\bibitem{masanes2011}
Masanes,~L.; M\"uller,~M.P. A derivation of quantum theory from physical requirements. \textit{New J. Phys.} \emph{2011}, \textit{13}, 063001.

\bibitem{pawlowski2009}
Paw\l owski,~M.; Paterek,~T.; Kaszlikowski,~D.; Scarani,~V.; Winter,~A.; \.Zukowski,~M. Information causality as a physical principle. \textit{Nature} \emph{2009}, \textit{461}, 1101--1104.

\bibitem{hartley1928}
Hartley,~R.V.L. Transmission of Information. \textit{Bell Syst. Tech. J.} \emph{1928}, \textit{7}, 535--563.

\bibitem{sakharov1967}
Sakharov,~A.D. Violation of CP invariance, C asymmetry, and baryon asymmetry of the universe. \textit{JETP Lett.} \emph{1967}, \textit{5}, 24--27.

\end{thebibliography}

% Preprint mode: skip MDPI Publisher's Note
% \PublishersNote{}
\end{adjustwidth}
\end{document}
